%%%%%%%%%%%%%%%%%%%%%%%%%%%%%%%%%%%%%%%%%%%%%%%%%%%%%%%%%%%%%%%%%%%%%%%%%%%%%%
%  PRÉSENTATION BEAMER 16:9 -- gabarit INSA CENTRE VAL DE LOIRE : main.tex
%
%  Ce fichier est réutilisable pour n'importe quelle présentation : il contient le
%  thème (couleurs INSA, page de titre, pied de page, encadrés, styles TikZ,
%  pictogrammes). Le contenu est dans des fichiers séparés, importés avec \include.
%
%  - Nouvelle présentation : modifier le bloc « Informations du document », remplacer
%    la liste des \include (un fichier chap_<nom>.tex par partie du plan ; chaque
%    fichier commence par \section{...} et ne contient ni préambule ni
%    \begin{document}), puis mettre à jour sources.bib.
%  - Une partie du plan = un \section = un fichier. L'ordre des \include donne l'ordre
%    du sommaire. Pour ne compiler que certaines parties (la numérotation est
%    conservée) : décommenter \includeonly{...} plus bas.
%  - Dans une diapositive :
%       \begin{frame}{Titre} ... \end{frame}
%       \intervenant{NOM}   (facultatif) « Intervenant : NOM » en haut à droite
%       \etiquette{Texte}   (facultatif) remplace l'étiquette rouge en haut à gauche,
%                           qui reprend sinon le titre de la section en cours
%       \srcnote{...}       note de source en bas à gauche, ex. \srcnote{Source : \cite{cle}}
%  - Encadrés : idee (rose, bord rouge), pointcle (pêche) et sa variante compacte pcompact
%    (diapositives très remplies), remarque[titre] (gris), encadre{titre} (rouge), carte (beige ;
%    carte[height=3cm,valign=top] pour aligner plusieurs cartes) avec \gros{chiffre ou mot-clé}
%    en grand. Listes compactes dans les encadrés : puces et numeros. Coche / croix en début
%    de ligne : \oui, \non. Tableaux : tabular ou tabularx avec \rowcolor{insarose}.
%  - Ponctuation française : babel-french est désactivé pour « : ; ! ? » (conflit avec TikZ) ;
%    écrire donc texte~: suite, texte~?, texte\,!, texte\,; (espace insécable à la main).
%  - Dans les schémas TikZ d'une diapositive : toujours des accolades autour du corps d'un
%    \foreach (le corps d'une diapositive est lu avant d'être exécuté).
%  - Couleurs : insarouge, insagris (charte INSA), insaorange, insapeche, insapeche2,
%    insabeige, insasafran, insarose (relevées sur la présentation Trove).
%  - Styles TikZ des schémas : pcli, pfoc, psvc, pdat, pgris (pastilles), etiqf,
%    fluxo, fluxr, fluxd, callout, num, acteur, vie, msg, rep (+ \seqmsg, \seqrep).
%  - Pictogrammes TikZ : \pic at (x,y) {icone tel}; liste en fin de préambule.
%  - Page de titre : \insatitleframe (logo images/logo_insa_blanc.png ; à défaut, un
%    texte « INSA » est composé).
%  - Mise en page prévue pour aspectratio=169. Texte de 10 pt : les plus petits textes
%    des schémas restent à \footnotesize (8 pt) ou plus.
%  - Sources : seul le nom de domaine est imprimé dans la liste (package xstring, voir « Bibliographie »),
%    le lien complet reste cliquable dans le PDF. La liste est répartie sur trois diapositives d'après
%    le numéro de citation (bornes dans les filtres srcA, srcB, srcC de la section « Bibliographie »).
%  - Compilation : pdflatex main, biber main, pdflatex main, pdflatex main.
%    (deux passes de pdflatex au minimum : sommaire, décor TikZ de chaque page)
%  - Fichiers du projet : main.tex (gabarit et ordre des parties), un fichier par partie
%    (chap_idee, chap_fonctionnement, chap_marche, chap_concurrence, chap_marketing, chap_finances,
%    chap_france, etude, chap_swot, chap_quanti, conclusion), sources.bib, guillotime.sty et
%    images/. questionnaire.tex est un document A4 autonome (pdflatex questionnaire), à imprimer
%    ou à recopier dans un formulaire en ligne : ses numéros de question sont ceux de chap_quanti.
%%%%%%%%%%%%%%%%%%%%%%%%%%%%%%%%%%%%%%%%%%%%%%%%%%%%%%%%%%%%%%%%%%%%%%%%%%%%%%
\documentclass[aspectratio=169,10pt,t]{beamer}

%%%%% Informations du document (à modifier pour chaque présentation) %%%%%
\newcommand{\nomproduit}{FreeWave}   % nom de travail du projet (à remplacer)
\newcommand{\doctitre}{\nomproduit}
\newcommand{\docsoustitre}{La messagerie radio libre,\\ sans opérateur ni abonnement}
\newcommand{\doctagline}{Ton réseau, tes règles}
\newcommand{\docauteur}{Groupe 4}
\newcommand{\docprof}{STI 5A ASL-IA}
\newcommand{\docecole}{INSA Centre Val de Loire}
%\newcommand{\docdate}{\today}
\newcommand{\docdate}{\aujourdhui}

%%%%% Langue, encodage, polices %%%%%
\usepackage[T1]{fontenc}
\usepackage[utf8]{inputenc}
\usepackage{lmodern}          % police machine à écrire (Latin Modern)
\usepackage{tgpagella}        % texte courant : serif proche de celui de la présentation Trove
\usepackage{tgheros}          % titres : sans serif (équivalent d'Arial / Liberation Sans)
% Pas de ligature en police machine à écrire : \texttt{--option} garde ses deux tirets.
\usepackage[protrusion=false,expansion=false,tracking=false,kerning=false,spacing=false]{microtype}
\DisableLigatures[-]{encoding = *, family = tt*}
\usepackage[french]{babel}
% Beamer lit le corps d'une diapositive avant de l'exécuter : les caractères actifs « : ; ! ? » de babel-french
% casseraient alors les coordonnées polaires (\a:1.7) et la syntaxe calc (\p!0.5!\q) de TikZ. On les désactive ;
% les espaces insécables se mettent à la main (texte~: suite, texte~?, texte\,!).
\AtBeginDocument{\shorthandoff{:;!?}}
\usefonttheme{serif}
\usepackage{amsmath}
\usepackage{amssymb}
\usepackage{pifont}           % \ding{51} (coche), \ding{55} (croix)
% symboles courants saisis en UTF-8 dans le texte
\DeclareUnicodeCharacter{2248}{\ensuremath{\approx}}
\DeclareUnicodeCharacter{2192}{\ensuremath{\rightarrow}}
\DeclareUnicodeCharacter{2264}{\ensuremath{\leq}}
\DeclareUnicodeCharacter{2265}{\ensuremath{\geq}}
\DeclareUnicodeCharacter{2212}{\ensuremath{-}}
\DeclareUnicodeCharacter{2260}{\ensuremath{\neq}}

%%%%% Charte graphique INSA CVL %%%%%
\usepackage{xcolor}
\definecolor{insarouge}{HTML}{E42618} % rouge du logo INSA (couleur principale)
\definecolor{insagris}{HTML}{5F5F5F}  % gris du logo INSA (couleur secondaire)
\definecolor{insaorange}{HTML}{C57F17}  % contour et flèches des services « voisins »
\definecolor{insapeche}{HTML}{FCDDB0}   % fond des services « voisins »
\definecolor{insapeche2}{HTML}{FCE2BB}  % fond des encadrés d'idée
% Couleurs relevées sur la page de titre et les encadrés de la présentation Trove
\definecolor{insabeige}{HTML}{F8F0EC}   % panneaux clairs de la page de titre
\definecolor{insasafran}{HTML}{F69F1D}  % blocs orange de la page de titre
\definecolor{insarose}{HTML}{F9CECC}    % fond des encadrés « idée » et des en-têtes de tableau

%%%%% Thème Beamer : structure, marges, listes %%%%%
\setbeamertemplate{navigation symbols}{}
\setbeamersize{text margin left=1.15cm,text margin right=1.15cm}
\setbeamercolor{structure}{fg=insarouge}
\setbeamercolor{normal text}{fg=black,bg=white}
\setbeamercolor{frametitle}{fg=black}
\setbeamercolor{block title}{fg=white,bg=insarouge}
\setbeamercolor{block body}{fg=black,bg=insarose}
\setbeamerfont{frametitle}{family=\sffamily,series=\bfseries,size=\fontsize{15}{18}\selectfont}
\setbeamerfont{framesubtitle}{family=\sffamily,series=\bfseries,size=\small}
\setbeamertemplate{itemize item}{\raisebox{0.05ex}{\textcolor{insarouge}{\footnotesize$\bullet$}}}
\setbeamertemplate{itemize subitem}{\textcolor{insagris}{\textendash}}
\setbeamertemplate{itemize subsubitem}{\textcolor{insagris}{\tiny$\bullet$}}
\setbeamertemplate{enumerate item}{\textcolor{insarouge}{\bfseries\insertenumlabel.}}
\setbeamerfont{itemize/enumerate body}{size=\normalsize}
\setbeamerfont{itemize/enumerate subbody}{size=\small}
\setlength{\leftmargini}{1.1em}
\setlength{\leftmarginii}{1.1em}
\addtolength{\itemsep}{0.1em}
\hypersetup{pdftitle={\doctitre}, pdfauthor={\docauteur}}

%%%%% Décor de chaque diapositive (étiquette de section, intervenant, pied de page rouge) %%%%%
% Unités du décor : 1/935 de la largeur et 1/527 de la hauteur de la page (mesures relevées sur la
% présentation Trove). \etiquette{...} et \intervenant{...} s'utilisent à l'intérieur d'une diapositive.
\usepackage{etoolbox}
\def\insatag{\insertsection}
\def\insainterv{}
\newcommand{\etiquette}[1]{\gdef\insatag{#1}}
\newcommand{\intervenant}[1]{\gdef\insainterv{#1}}
% les deux réglages ne valent que pour la diapositive en cours : remis à zéro au début de chaque frame
\AtBeginEnvironment{frame}{\gdef\insatag{\insertsection}\gdef\insainterv{}}
\newcommand{\insaunitx}{0.0010695\paperwidth}
\newcommand{\insaunity}{0.0018975\paperheight}
\setbeamertemplate{background}{%
  \begin{tikzpicture}[remember picture,overlay,shift={(current page.north west)}]
    \node[anchor=north west,inner sep=0pt,text=insarouge,font=\sffamily\bfseries\scriptsize]
      at (1.15cm,-0.28cm) {\insatag};
    \ifx\insainterv\empty\else
      \node[anchor=north east,inner sep=0pt,text=insagris!70,font=\scriptsize]
        at (\paperwidth-1.15cm,-0.28cm) {Intervenant : \insainterv};
    \fi
    \begin{scope}[shift={(current page.south east)},x=\insaunitx,y=-\insaunity]
      \foreach \i in {0,...,17} {%
        \draw[insarouge,line width={0.3pt+0.3pt*mod(\i,2)}]
          ({-55+3.0*\i+1},-109) -- ({-55+3.0*\i+1},-55);}
      \foreach \i in {0,...,4} {%
        \draw[insagris!55,line width=0.35pt] ({-122+11*\i},-55) -- ({-122+11*\i},0);}
      \fill[insarouge] (-55,-55) rectangle (0,0);
      \node[white,font=\bfseries\large] at (-27.5,-27.5) {\insertframenumber};
    \end{scope}
  \end{tikzpicture}}
\setbeamertemplate{headline}{}
% marge basse de 0,62 cm réservée aux notes de source (\srcnote) : le contenu ne descend pas plus bas
\setbeamertemplate{footline}{\hbox{\vrule width 0pt height 0pt depth 0.62cm}}
\setbeamertemplate{frametitle}{%
  \vskip0.52cm
  {\usebeamerfont{frametitle}\usebeamercolor[fg]{frametitle}\insertframetitle\par}%
  \vskip0.08cm}
\setbeamertemplate{framesubtitle}{%
  {\usebeamerfont{framesubtitle}\usebeamercolor[fg]{frametitle}\insertframesubtitle\par}%
  \vskip0.05cm}

% Note de source en bas à gauche de la diapositive : \srcnote{Sources : \cite{cle}}
\newcommand{\srcnote}[1]{%
  \begin{tikzpicture}[remember picture,overlay]
    \node[anchor=south west,inner sep=0pt,text=insagris,font=\tiny,text width=0.72\paperwidth,align=left]
      at ([xshift=1.15cm,yshift=0.22cm]current page.south west) {#1};
  \end{tikzpicture}}

%%%%% Sommaire : numéros rouges, pas de sous-parties %%%%%
\setbeamertemplate{section in toc}{%
  \leavevmode\strut\textcolor{insarouge}{\bfseries\inserttocsectionnumber.}\hspace{0.6em}\inserttocsection\strut\par\vspace{-0.85em}}
\setbeamerfont{section in toc}{size=\large}
\setbeamercolor{section in toc}{fg=black}

%%%%% Tableaux, graphiques, schémas %%%%%
\usepackage{array}
\usepackage{makecell}
\usepackage{booktabs}
\usepackage{colortbl}
\usepackage{tabularx}
\usepackage{multirow}
\usepackage{graphicx}
\graphicspath{{images/}}
\usepackage{tikz}
\usetikzlibrary{positioning, arrows.meta, fit, backgrounds, calc, shapes.geometric, babel,
                decorations.pathmorphing, decorations.pathreplacing, patterns}
\usepackage{pgfplots}
\pgfplotsset{compat=1.16}
\usepackage[skins]{tcolorbox}
\tcbset{boxrule=0.6pt,arc=2pt,left=6pt,right=6pt,top=4pt,bottom=4pt}

%%%%% Pour la date %%%%%
\usepackage[delan,parigot]{guillotime}

%%%%% Lignes de code (frames [fragile]) %%%%%
\usepackage{listings}
\lstset{
  basicstyle=\ttfamily\footnotesize,
  backgroundcolor=\color{insagris!8},
  frame=single,
  rulecolor=\color{insagris!60},
  breaklines=true,
  upquote=true,
  showstringspaces=false,
  literate={é}{{\'e}}1 {è}{{\`e}}1 {ê}{{\^e}}1 {ë}{{\"e}}1 {à}{{\`a}}1 {â}{{\^a}}1
           {î}{{\^i}}1 {ï}{{\"i}}1 {ô}{{\^o}}1 {ù}{{\`u}}1 {û}{{\^u}}1 {ç}{{\c c}}1
           {É}{{\'E}}1 {È}{{\`E}}1 {À}{{\`A}}1 {Ç}{{\c C}}1 {œ}{{\oe}}1
}

%%%%% Encadrés %%%%%
% Idée à retenir (rose, bord rouge) ; point clé en une phrase (pêche, sans bord) ; remarque grise ; encadré titré.
% (« definition » existe déjà dans Beamer : l'encadré titré s'appelle donc « encadre ».)
\newtcolorbox{idee}{colback=insarose,colframe=insarouge,boxrule=0.7pt,arc=2pt}
\newtcolorbox{pointcle}{colback=insapeche2,colframe=insapeche2,boxrule=0pt,arc=3pt}
% Variante compacte du point clé (marges haute et basse réduites), pour les diapositives très remplies.
\newtcolorbox{pcompact}{colback=insapeche2,colframe=insapeche2,boxrule=0pt,arc=3pt,top=2pt,bottom=2pt}
\newtcolorbox{remarque}[1][Remarque]{colback=insagris!7,colframe=insagris,
  fonttitle=\bfseries\sffamily\small,title={#1}}
\newtcolorbox{encadre}[1]{colback=insarouge!4,colframe=insarouge,
  fonttitle=\bfseries\sffamily\small,title={#1}}
% Carte beige (chiffre clé, vignette) ; \gros{...} : chiffre ou mot-clé en grand, rouge
\newtcolorbox{carte}[1][]{colback=insabeige,colframe=insabeige,boxrule=0pt,arc=3pt,
  left=5pt,right=5pt,top=3pt,bottom=3pt,#1}
\newcommand{\gros}[1]{{\sffamily\bfseries\large\color{insarouge}#1}}
% Coche rouge / croix grise en début de ligne : \oui Texte ...   \non Texte ...
\newcommand{\oui}{\textcolor{insarouge}{\ding{51}}\hspace{0.35em}}
\newcommand{\non}{\textcolor{insagris}{\ding{55}}\hspace{0.35em}}
% Listes compactes (corps en \footnotesize) pour les encadrés : \begin{puces}\item ...\end{puces}
% et \begin{numeros}\item ...\end{numeros}  (les listes beamer normales restent en taille normale)
\newenvironment{puces}
  {\begingroup\setbeamerfont{itemize/enumerate body}{size=\footnotesize}\begin{itemize}\setlength{\itemsep}{1pt}}
  {\end{itemize}\endgroup}
\newenvironment{numeros}
  {\begingroup\setbeamerfont{itemize/enumerate body}{size=\footnotesize}\begin{enumerate}\setlength{\itemsep}{1pt}}
  {\end{enumerate}\endgroup}

%%%%% Styles TikZ pour les schémas (repris du gabarit de cours) %%%%%
% svc    : boîte d'un composant rouge pâle          svcgris : composant secondaire
% bus    : barre horizontale                         flux    : flèche entre deux composants
\tikzset{
  svc/.style     = {draw=insarouge, fill=insarouge!8, thick, rounded corners=2pt,
                    minimum height=7mm, inner sep=3pt, align=center},
  svcgris/.style = {draw=insagris, fill=insagris!8, thick, rounded corners=2pt,
                    minimum height=7mm, inner sep=3pt, align=center},
  bus/.style     = {draw=insagris, fill=insagris!15, thick, minimum height=6mm, align=center},
  flux/.style    = {-{Stealth[length=2mm]}, thick, insagris},
}
% Schémas « à la Trove » : pastilles colorées par rôle, flèches étiquetées, encadré d'idée.
%   pcli  : client / utilisateur                  pfoc : élément au centre du schéma (rouge plein)
%   psvc  : élément voisin (pêche)                pdat : donnée, ressource (rose)
%   pgris : infrastructure, existant              etiqf: étiquette posée sur une flèche
%   fluxo / fluxr / fluxd : flèche orange / rouge / rouge pointillée
%   callout : encadré pêche sous le schéma         num : pastille numérotée d'une étape
% Diagrammes de séquence : acteur (en-tête), vie (ligne de vie), msg / rep (message / réponse)
\tikzset{
  pastille/.style = {rounded corners=0.3cm, minimum height=0.65cm, inner xsep=7pt, inner ysep=3pt,
                     thick, align=center, font=\footnotesize\bfseries},
  pcli/.style  = {pastille, draw=insarouge, fill=white, text=insarouge},
  pfoc/.style  = {pastille, draw=black, fill=insarouge, text=white, line width=1.2pt},
  psvc/.style  = {pastille, draw=insaorange, fill=insapeche, text=black},
  pdat/.style  = {pastille, draw=insarouge!90!black, fill=insarouge!30, text=black},
  pgris/.style = {pastille, draw=insagris, fill=insagris!12, text=black},
  etiqf/.style = {fill=white, inner sep=1.5pt, font=\scriptsize, text=insagris, align=center},
  fluxo/.style = {-{Stealth[length=2mm]}, thick, insaorange},
  fluxr/.style = {-{Stealth[length=2mm]}, thick, insarouge},
  fluxd/.style = {-{Stealth[length=2mm]}, thick, insarouge, dashed},
  callout/.style = {fill=insapeche2, rounded corners=3pt, inner sep=5pt, align=center,
                    font=\footnotesize},
  num/.style   = {circle, fill=insarouge, text=white, font=\scriptsize\bfseries, inner sep=1pt,
                  minimum size=4.5mm},
  acteur/.style = {draw=insarouge, fill=insarouge!8, thick, rounded corners=2pt, minimum height=0.7cm,
                   minimum width=2.2cm, align=center, font=\footnotesize\bfseries, inner sep=3pt},
  vie/.style   = {insagris!70, dashed, thin},
  msg/.style   = {-{Stealth[length=2mm]}, thick, insagris},
  rep/.style   = {-{Stealth[length=2mm]}, thick, insagris, dashed},
}
% Messages d'un diagramme de séquence : \seqmsg{x de départ}{x d'arrivée}{y}{numéro (ou vide)}{étiquette}
% (flèche pleine) et \seqrep (flèche pointillée = réponse).
\newcommand{\seqmsg}[5]{%
  \draw[msg] (#1,#3) -- (#2,#3) node[midway, above=1pt, font=\footnotesize, align=center, text=black, fill=white, inner sep=1pt] {#5};%
  \ifx\relax#4\relax\else\node[num] at ({#1+(#2>#1?-0.4:0.4)},#3) {#4};\fi}
\newcommand{\seqrep}[5]{%
  \draw[rep] (#1,#3) -- (#2,#3) node[midway, above=1pt, font=\footnotesize, align=center, text=black, fill=white, inner sep=1pt] {#5};%
  \ifx\relax#4\relax\else\node[num] at ({#1+(#2>#1?-0.4:0.4)},#3) {#4};\fi}

%%%%% Pictogrammes TikZ : \pic at (x,y) {icone ...}; (centrés, environ 1 cm de côté) %%%%%
% Ils prennent la couleur courante : \pic[insarouge] at (0,0) {icone tel}; \pic[scale=1.5] ...
% Liste : tel, tour, antenne, mont, immeuble, vague, cadenas, euro, code, usine, cairn,
%         croix, coche, tente, boite, personne, drapeau, batterie, pin, alerte
\tikzset{
  pics/icone tel/.style = {code={
    \draw[thick,fill=white,rounded corners=1.5pt] (-0.2,-0.36) rectangle (0.2,0.36);
    \draw[thin,fill=insagris!12] (-0.145,-0.26) rectangle (0.145,0.27);
    \fill (0,-0.31) circle (0.025);}},
  pics/icone tour/.style = {code={
    \draw[thick] (0,0.3) -- (-0.18,-0.38) (0,0.3) -- (0.18,-0.38)
                 (-0.1,-0.1) -- (0.1,-0.1) (-0.14,-0.24) -- (0.14,-0.24);
    \fill (0,0.32) circle (0.04);
    \foreach \r in {0.17,0.3} {
      \draw[thin] (0,0.32) ++(-35:\r) arc (-35:35:\r);
      \draw[thin] (0,0.32) ++(145:\r) arc (145:215:\r);}}},
  pics/icone antenne/.style = {code={
    \draw[thick,fill=white,rounded corners=1.5pt] (-0.34,-0.36) rectangle (0.3,0.0);
    \fill[insagris!35] (-0.27,-0.3) rectangle (-0.02,-0.08);
    \fill (0.12,-0.1) circle (0.035);
    \draw[thick] (0.2,0.0) -- (0.2,0.34);
    \fill (0.2,0.37) circle (0.035);
    \foreach \r in {0.14,0.24} {
      \draw[thin] (0.2,0.37) ++(-40:\r) arc (-40:40:\r);
      \draw[thin] (0.2,0.37) ++(140:\r) arc (140:220:\r);}}},
  pics/icone mont/.style = {code={
    \fill[insagris!40] (-0.5,-0.3) -- (-0.15,0.28) -- (0.2,-0.3) -- cycle;
    \fill[insagris] (-0.12,-0.3) -- (0.24,0.4) -- (0.6,-0.3) -- cycle;
    \fill[white] (0.24,0.4) -- (0.13,0.18) -- (0.2,0.2) -- (0.25,0.14) -- (0.31,0.2) -- (0.37,0.18) -- cycle;}},
  pics/icone immeuble/.style = {code={
    \draw[thick,fill=white] (-0.34,-0.36) -- (-0.34,0.22) -- (-0.18,0.38) -- (-0.08,0.16) -- (0.06,0.3)
      -- (0.14,0.06) -- (0.3,0.2) -- (0.3,-0.36) -- cycle;
    \foreach \x in {-0.2,0.02,0.18} \foreach \y in {-0.08,-0.25}
      \fill (\x-0.04,\y-0.04) rectangle (\x+0.04,\y+0.04);}},
  pics/icone vague/.style = {code={
    \foreach \y in {-0.22,0,0.22}
      \draw[thick,domain=-0.4:0.4,samples=24,smooth] plot (\x,{\y+0.06*sin(deg(\x*15))});}},
  pics/icone cadenas/.style = {code={
    \draw[thick] (-0.12,0.02) -- (-0.12,0.14) arc (180:0:0.12) -- (0.12,0.02);
    \draw[thick,fill=white,rounded corners=1pt] (-0.2,-0.32) rectangle (0.2,0.02);
    \fill (0,-0.15) circle (0.035);}},
  pics/icone euro/.style = {code={
    \draw[thick,fill=white] (0,0) circle (0.33);
    \node[font=\bfseries\small,text=black] at (0,0) {\texteuro};}},
  pics/icone code/.style = {code={
    \draw[thick,line cap=round,line join=round] (-0.1,0.2) -- (-0.32,0) -- (-0.1,-0.2)
         (0.1,0.2) -- (0.32,0) -- (0.1,-0.2) (0.07,0.27) -- (-0.07,-0.27);}},
  pics/icone usine/.style = {code={
    \draw[thick,fill=white] (-0.4,-0.3) -- (-0.4,0.06) -- (-0.2,0.2) -- (-0.2,0.06) -- (0,0.2) -- (0,0.06)
      -- (0.2,0.2) -- (0.2,-0.3) -- cycle;
    \draw[thick,fill=white] (0.2,-0.3) rectangle (0.34,0.36);
    \fill (-0.3,-0.12) rectangle (-0.22,-0.04);\fill (-0.1,-0.12) rectangle (-0.02,-0.04);
    \fill (0.07,-0.12) rectangle (0.15,-0.04);}},
  pics/icone cairn/.style = {code={
    \fill[insagris!70] (0,-0.34) ellipse (0.46 and 0.14);
    \fill[insagris!45] (0.03,-0.1) ellipse (0.36 and 0.12);
    \fill[insarouge]   (-0.02,0.1) ellipse (0.27 and 0.11);
    \fill[insagris!30] (0.02,0.27) ellipse (0.18 and 0.085);
    \fill[insasafran]  (0,0.4) ellipse (0.1 and 0.065);}},
  pics/icone croix/.style = {code={
    \fill[insarouge] (0,0) circle (0.17);
    \draw[white,line width=1.3pt,line cap=round] (-0.07,-0.07) -- (0.07,0.07) (-0.07,0.07) -- (0.07,-0.07);}},
  pics/icone coche/.style = {code={
    \fill (0,0) circle (0.17);
    \draw[white,line width=1.3pt,line cap=round,line join=round] (-0.08,0.0) -- (-0.025,-0.06) -- (0.08,0.07);}},
  pics/icone tente/.style = {code={
    \draw[thick,fill=white] (-0.42,-0.28) -- (0,0.32) -- (0.42,-0.28) -- cycle;
    \fill (-0.11,-0.28) -- (0,0.0) -- (0.11,-0.28) -- cycle;
    \draw[thick] (-0.5,-0.28) -- (0.5,-0.28);}},
  pics/icone boite/.style = {code={
    \draw[thick,fill=white] (-0.3,-0.3) rectangle (0.2,0.1);
    \draw[thick,fill=insagris!15] (-0.3,0.1) -- (-0.15,0.3) -- (0.35,0.3) -- (0.2,0.1) -- cycle;
    \draw[thick,fill=insagris!30] (0.2,0.1) -- (0.35,0.3) -- (0.35,-0.1) -- (0.2,-0.3) -- cycle;}},
  pics/icone personne/.style = {code={
    \fill (0,0.22) circle (0.1);
    \fill (-0.2,-0.32) -- (-0.2,-0.14) arc (180:0:0.2) -- (0.2,-0.32) -- cycle;}},
  pics/icone drapeau/.style = {code={
    \fill[blue!60!black] (-0.3,-0.2) rectangle (-0.1,0.2);
    \fill[white] (-0.1,-0.2) rectangle (0.1,0.2);
    \fill[insarouge] (0.1,-0.2) rectangle (0.3,0.2);
    \draw[insagris!60,thin] (-0.3,-0.2) rectangle (0.3,0.2);}},
  pics/icone batterie/.style = {code={
    \draw[thick,fill=white,rounded corners=1pt] (-0.34,-0.18) rectangle (0.26,0.18);
    \fill (0.26,-0.07) rectangle (0.33,0.07);
    \fill (-0.27,-0.11) rectangle (-0.07,0.11);\fill (-0.03,-0.11) rectangle (0.17,0.11);}},
  pics/icone pin/.style = {code={
    \fill (0,0.12) circle (0.2);
    \fill (-0.19,0.06) -- (0,-0.36) -- (0.19,0.06) -- cycle;
    \fill[white] (0,0.12) circle (0.08);}},
  pics/icone alerte/.style = {code={
    \draw[thick,fill=insapeche,line join=round] (0,0.34) -- (-0.4,-0.3) -- (0.4,-0.3) -- cycle;
    \node[font=\bfseries\small,text=black,inner sep=0pt] at (0,-0.08) {!};}},
}

%%%%% Page de titre (décor de la page de titre de la présentation Trove) %%%%%
% \insatitleframe : diapositive 1, sans pied de page ; logo, titre, sous-titre, auteur, école, date.
\newcommand{\insalogo}[1]{%
  \IfFileExists{images/logo_insa_blanc.png}%
    {\includegraphics[width=#1]{images/logo_insa_blanc.png}}%
    {\textcolor{white}{\sffamily\bfseries\Large INSA}\ \textcolor{white}{\sffamily\tiny CENTRE VAL DE LOIRE}}}
\newcommand{\insatitleframe}{%
  \begingroup
  \setbeamertemplate{background}{}%
  \begin{frame}[plain]
    \begin{tikzpicture}[remember picture,overlay,shift={(current page.north west)},x=\insaunitx,y=-\insaunity]
      % panneaux beiges, bloc rouge, blocs orange
      \fill[insabeige] (669,0) rectangle (935,295);
      \fill[insabeige] (609,235) rectangle (669,295);
      \fill[insabeige] (728,295) rectangle (935,527);
      \fill[insarouge] (0,0) -- (669,0) -- (669,235) -- (609,235) -- (609,295) -- (550,295) -- (550,353)
        -- (0,353) -- cycle;
      \fill[insasafran] (609,295) rectangle (728,353);
      \fill[insasafran] (609,353) rectangle (669,413);
      % lignes décoratives
      \foreach \i in {0,...,30} {\draw[insarouge,line width=0.3pt] ({842+3.0*\i},0) -- ({842+3.0*\i},295);}
      \foreach \x in {749,772,796,819} {\draw[insasafran,line width=0.4pt] (\x,295) -- (\x,527);}
      \draw[insagris!60,line width=0.4pt] (841,295) -- (841,527);
      \foreach \x in {65,76,87,98,109} {\draw[insarouge,line width=0.5pt] (\x,476) -- (\x,527);}
      % grille de « + »
      \foreach \i in {0,...,7} \foreach \j in {0,...,7} {%
        \draw[white,line width=0.4pt] ({503+21.1*\i-3},{16+21.1*\j}) -- ({503+21.1*\i+3},{16+21.1*\j})
                                      ({503+21.1*\i},{16+21.1*\j-3}) -- ({503+21.1*\i},{16+21.1*\j+3});}
      % logo, titre, sous-titre
      \node[anchor=north west,inner sep=0pt] at (65,58) {\insalogo{0.1647\paperwidth}};
      \node[anchor=north west,inner sep=0pt,text=white,text width=0.55\paperwidth,align=left] at (63,170)
        {{\sffamily\bfseries\fontsize{32}{34}\selectfont\doctitre\par}\vspace{0.25cm}%
         {\sffamily\fontsize{14}{17}\selectfont\docsoustitre\par}};
      % auteur, cursus, école, date
      \node[anchor=north west,inner sep=0pt,align=left,text width=0.62\paperwidth] at (63,392)
        {{\sffamily\fontsize{12}{15}\selectfont\textbf{\docauteur}\ ---\ \docprof\par}%
         {\sffamily\fontsize{10}{14}\selectfont\color{insagris}\docecole\ ---\ \docdate\par}};
      \node[font=\bfseries\large] at (908,497) {\insertframenumber};
    \end{tikzpicture}
  \end{frame}
  \endgroup}

%%%%% Bibliographie (sources.bib) %%%%%
\usepackage{csquotes}
\usepackage[backend=biber,style=numeric-comp,sorting=none,maxbibnames=3,minbibnames=1,
            doi=false,isbn=false,eprint=false]{biblatex}
\addbibresource{sources.bib}
\setcounter{biburlnumpenalty}{100}
\setcounter{biburlucpenalty}{100}
\setcounter{biburllcpenalty}{100}
\renewcommand*{\bibfont}{\fontsize{6}{7.2}\selectfont\raggedright}
\setlength{\bibitemsep}{0.1\baselineskip}
\AtEveryBibitem{\clearfield{urlyear}\clearfield{urlmonth}\clearfield{urlday}}
\urlstyle{same}
% Adresse imprimée = nom de domaine seulement (xstring) ; le lien complet reste cliquable dans le PDF.
% Cela permet de faire tenir de nombreuses sources sur une seule diapositive.
\usepackage{xstring}
\newcommand{\bibhost}[1]{\StrBehind{#1}{://}[\bibhosttmp]\StrBefore{\bibhosttmp}{/}[\bibhosttmpb]%
  \IfBeginWith{\bibhosttmpb}{www.}{\StrGobbleLeft{\bibhosttmpb}{4}}{\bibhosttmpb}}
\DeclareFieldFormat{url}{\href{#1}{\bibhost{#1}}}
\usepackage{multicol}
% La liste des sources est coupée en trois diapositives d'après le numéro de citation (ordre de première citation).
% \srcrange{n1}{n2} garde les sources numérotées de n1 à n2 ; ajuster les bornes si le nombre de sources change.
\newcommand{\srcrange}[2]{%
  \iffieldint{labelnumber}%
    {\ifboolexpe{test{\ifnumless{\thefield{labelnumber}}{#1}} or test{\ifnumgreater{\thefield{labelnumber}}{#2}}}{\skipentry}{}}%
    {\skipentry}}
\defbibcheck{srcA}{\srcrange{1}{25}}
\defbibcheck{srcB}{\srcrange{26}{50}}
\defbibcheck{srcC}{\srcrange{51}{999}}

\title{\doctitre}
\author{\docauteur}
\date{\docdate}

% Pour ne compiler que certaines parties (la numérotation est conservée) :
%\includeonly{chap_finances}

\begin{document}

%%%%% Page de titre : diapositive 1 %%%%%
\insatitleframe

%%%%% Sommaire (deux colonnes : parties 1 à 6, puis 7 à 11 ; ajuster si le nombre de parties change) %%%%%
\begin{frame}{Sommaire}
  \etiquette{Plan}
  \begin{columns}[T]
    \begin{column}{0.37\textwidth}
      \vspace{0.2cm}
      \tableofcontents[sections={1-6},hideallsubsections]
    \end{column}
    \begin{column}{0.37\textwidth}
      \vspace{0.2cm}
      \tableofcontents[sections={7-11},hideallsubsections]
    \end{column}
    \begin{column}{0.22\textwidth}
      \centering
      \vspace{0.3cm}
      \begin{tikzpicture}[scale=2.3]
        \pic at (0,0) {icone cairn};
      \end{tikzpicture}\\[0.2cm]
      {\sffamily\bfseries\color{insarouge}\doctagline}
    \end{column}
  \end{columns}
\end{frame}

%%%%% Parties du plan (un fichier par partie, chacun commence par \section) %%%%%
% L'ordre des \include donne l'ordre du sommaire : pour déplacer une partie, déplacer sa ligne.
% La partie SWOT cite les freins F1 à F7 de l'étude documentaire, qui doit donc la précéder.
%%%%%%%%%%%%%%%%%%%%%%%%%%%%%%%%%%%%%%%%%%%%%%%%%%%%%%%%%%%%%%%%%%%%%%%%%%%%%%
%  Partie 1 -- L'idée : le problème, notre réponse, réseau d'opérateur contre réseau maillé
%  (3 diapositives)
%  Note : dans les schémas TikZ, mettre des accolades à toutes les boucles \foreach
%  (le corps d'une diapositive est lu avant d'être exécuté).
%%%%%%%%%%%%%%%%%%%%%%%%%%%%%%%%%%%%%%%%%%%%%%%%%%%%%%%%%%%%%%%%%%%%%%%%%%%%%%
\section{L'idée}

%---------------------------------------------------------------------------
% Diapositive 3 : le problème (chaîne de dépendance)
%---------------------------------------------------------------------------
\begin{frame}{Quand le réseau s'arrête, tout s'arrête}
  % \intervenant{NOM}  % décommenter pour afficher l'intervenant
  \centering
  \begin{tikzpicture}[font=\footnotesize]
    \node[pcli,minimum height=0.9cm,text width=1.7cm] (toi) {Toi};
    \node[pgris,minimum height=0.9cm,text width=1.7cm,right=0.42cm of toi] (ant) {Antenne-relais};
    \node[pgris,minimum height=0.9cm,text width=1.7cm,right=0.42cm of ant] (op) {Réseau de l'opérateur};
    \node[pgris,minimum height=0.9cm,text width=1.7cm,right=0.42cm of op] (pf) {Plateforme (serveurs)};
    \node[pcli,minimum height=0.9cm,text width=1.7cm,right=0.42cm of pf] (ami) {Ton contact};
    \foreach \a/\b in {toi/ant,ant/op,op/pf,pf/ami} {\draw[fluxr] (\a) -- (\b);}
    % points de défaillance : une croix au-dessus de chaque maillon centralisé
    \foreach \n in {ant,op,pf} {\pic[scale=1.1] at ([yshift=0.6cm]\n.north) {icone croix};}
    \node[font=\footnotesize\bfseries,text=insarouge,anchor=east,align=right]
      at ([yshift=0.6cm,xshift=-0.4cm]ant.north) {trois points\\de défaillance};
    \node[font=\scriptsize,text=insagris,text width=2.2cm,align=center,below=0.08cm of ant]
      {tempête, séisme, coupure de courant};
    \node[font=\scriptsize,text=insagris,text width=2.2cm,align=center,below=0.08cm of op]
      {panne, saturation, coupure décidée};
    \node[font=\scriptsize,text=insagris,text width=2.2cm,align=center,below=0.08cm of pf]
      {serveur injoignable, compte suspendu};
  \end{tikzpicture}

  \vspace{0.2cm}
  \begin{columns}[T]
    \begin{column}{0.325\textwidth}
      \begin{carte}
        \gros{99,2 à 99,6\,\%}\\[1pt]
        {\footnotesize de couverture 4G selon l'opérateur (fin 2025), mais des zones blanches persistent.}
      \end{carte}
    \end{column}
    \begin{column}{0.325\textwidth}
      \begin{carte}
        \gros{17,88\,€ par mois}\\[1pt]
        {\footnotesize facture moyenne d'un forfait mobile (ARCEP, T4 2025).}
      \end{carte}
    \end{column}
    \begin{column}{0.325\textwidth}
      \begin{carte}
        \gros{Avant les secours}\\[1pt]
        {\footnotesize les communications tombent souvent avant même l'arrivée des secours.}
      \end{carte}
    \end{column}
  \end{columns}
  \srcnote{Sources~: \cite{arcep_mobile,arcep_prix,laresilience}}
\end{frame}

%---------------------------------------------------------------------------
% Diapositive 4 : notre idée (trois briques)
%---------------------------------------------------------------------------
\begin{frame}{Notre idée~: un kit radio que chacun contrôle}
  \centering
  \begin{tikzpicture}[font=\footnotesize]
    \tikzset{brique/.style={draw=insarouge,fill=white,thick,rounded corners=4pt,minimum width=2.8cm,
                            minimum height=2.2cm,text width=2.6cm,align=center,inner sep=3pt}}
    \node[brique] (a) at (0,0) {};
    \node[brique] (b) at (3.6,0) {};
    \node[brique] (c) at (7.2,0) {};
    \node[brique,fill=insarose] (d) at (10.8,0) {};
    \pic[insarouge,scale=1.3] at ([yshift=0.62cm]a.center) {icone antenne};
    \pic[insarouge,scale=1.3] at ([yshift=0.62cm]b.center) {icone code};
    \pic[insarouge,scale=1.3] at ([yshift=0.62cm]c.center) {icone personne};
    \pic[scale=1.2] at ([yshift=0.58cm]d.center) {icone cairn};
    \node[text width=2.6cm,align=center,anchor=north] at ([yshift=0.1cm]a.center)
      {{\bfseries Une antenne}\\Heltec V3, LoRa, préconfigurée};
    \node[text width=2.6cm,align=center,anchor=north] at ([yshift=0.1cm]b.center)
      {{\bfseries Un logiciel libre}\\code ouvert, sans compte ni serveur};
    \node[text width=2.6cm,align=center,anchor=north] at ([yshift=0.1cm]c.center)
      {{\bfseries Des voisins}\\chaque antenne relaie les autres};
    \node[text width=2.6cm,align=center,anchor=north] at ([yshift=0.1cm]d.center)
      {{\bfseries Ton réseau}\\une messagerie qui t'appartient};
    \node[font=\Huge\bfseries\sffamily,text=insarouge] at ($(a)!0.5!(b)$) {+};
    \node[font=\Huge\bfseries\sffamily,text=insarouge] at ($(b)!0.5!(c)$) {+};
    \node[font=\Huge\bfseries\sffamily,text=insarouge] at ($(c)!0.5!(d)$) {=};
  \end{tikzpicture}

  \vspace{0.12cm}
  \begin{columns}[T]
    \begin{column}{0.325\textwidth}
      \begin{carte}\centering\gros{≈ 69\,€}\\[1pt]{\footnotesize achat unique, sans abonnement}\end{carte}
    \end{column}
    \begin{column}{0.325\textwidth}
      \begin{carte}\centering\gros{Zéro compte}\\[1pt]{\footnotesize ni opérateur, ni réseau social}\end{carte}
    \end{column}
    \begin{column}{0.325\textwidth}
      \begin{carte}\centering\gros{GPL-3.0}\\[1pt]{\footnotesize logiciel libre, auditable, gratuit}\end{carte}
    \end{column}
  \end{columns}
  \vspace{0.08cm}
  \begin{idee}\centering\footnotesize
    Écrire à ses proches à plusieurs kilomètres, \textbf{sans réseau mobile, sans internet, sans abonnement}.
  \end{idee}
  \srcnote{Sources~: \cite{heltec_v3,meshtastic_fw}}
\end{frame}

%---------------------------------------------------------------------------
% Diapositive 5 : réseau d'opérateur (étoile) contre réseau maillé
%---------------------------------------------------------------------------
\begin{frame}{Réseau d'opérateur ou réseau maillé~?}
  \centering
  \begin{tikzpicture}[font=\footnotesize]
    % ---- à gauche : réseau en étoile (tout passe par le centre)
    \begin{scope}[shift={(-3.55,0)}]
      \foreach \a in {30,90,...,330} {
        \draw[insagris!70,thick] (\a:0.55) -- (\a:1.5);
        \pic[insagris,scale=0.62] at (\a:1.72) {icone personne};
      }
      \pic[insagris,scale=1.35] at (0,0.05) {icone tour};
      \pic[scale=1.15] at (0.62,0.5) {icone croix};
      \node[pgris,anchor=north] at (0,-2.3) {Réseau d'opérateur};
      \node[font=\scriptsize,text=insagris,align=center,anchor=north] at (0,-3.0)
        {le centre tombe~: plus personne ne communique};
    \end{scope}
    \node[font=\large\bfseries\sffamily,text=insagris] at (-0.55,0.1) {contre};
    % ---- à droite : réseau maillé (chaque antenne relaie)
    \begin{scope}[shift={(3.3,0)}]
      \coordinate (A) at (-1.85,-0.2); \coordinate (B) at (-1.0,0.95); \coordinate (C) at (-0.1,-0.25);
      \coordinate (D) at (0.85,0.95); \coordinate (E) at (1.8,0.05); \coordinate (F) at (-1.1,-1.3);
      \coordinate (G) at (1.0,-1.3);
      \foreach \p/\q in {A/B,A/F,B/C,C/F,C/G,F/G,G/E,C/D} {\draw[insagris!70,thick] (\p) -- (\q);}
      \foreach \p/\q in {B/D,D/E,C/E} {\draw[insagris!45,thick,dashed] (\p) -- (\q);}
      \foreach \p/\q in {B/D,C/E} {\pic[scale=0.8] at ($(\p)!0.5!(\q)$) {icone croix};}
      \draw[insarouge,line width=1.7pt,-{Stealth[length=2.6mm]}] (A) -- (B) -- (C) -- (G) -- (E);
      \foreach \p in {A,B,C,D,E,F,G} {
        \node[circle,draw=insarouge,fill=white,thick,minimum size=0.42cm,inner sep=0pt] at (\p) {};}
      \node[font=\footnotesize\bfseries,text=insarouge,anchor=south east] at ([xshift=-0.05cm,yshift=0.12cm]A.north west) {Toi};
      \node[font=\footnotesize\bfseries,text=insarouge,anchor=west] at ([xshift=0.28cm]E) {Ton contact};
      \node[pcli,anchor=north] at (0,-2.3) {Réseau maillé};
      \node[font=\scriptsize,text=insagris,align=center,anchor=north] at (0,-3.0)
        {chaque antenne est un relais~: le message contourne la panne};
    \end{scope}
  \end{tikzpicture}
  \srcnote{Principe~: chaque nœud retransmet les paquets qu'il reçoit (diffusion gérée) \cite{meshtastic_algo}}
\end{frame}
              % 1.  L'idée
%%%%%%%%%%%%%%%%%%%%%%%%%%%%%%%%%%%%%%%%%%%%%%%%%%%%%%%%%%%%%%%%%%%%%%%%%%%%%%
%  Partie 2 -- Fonctionnement : vue d'ensemble, partie physique, partie logicielle,
%  trajet d'un message (4 diapositives)
%%%%%%%%%%%%%%%%%%%%%%%%%%%%%%%%%%%%%%%%%%%%%%%%%%%%%%%%%%%%%%%%%%%%%%%%%%%%%%
\section{Fonctionnement}

%---------------------------------------------------------------------------
% Diapositive 6 : vue d'ensemble (physique + logiciel)
%---------------------------------------------------------------------------
\begin{frame}{Vue d'ensemble~: de ton téléphone à ton contact}
  \centering
  \begin{tikzpicture}[font=\footnotesize]
    \tikzset{
      bt/.style   = {{Stealth[length=2mm]}-{Stealth[length=2mm]},thick,insagris,dashed},
      lora/.style = {{Stealth[length=2mm]}-{Stealth[length=2mm]},thick,insarouge,decorate,
                     decoration={snake,amplitude=1.3pt,segment length=7pt,pre length=1.6mm,post length=1.6mm}},
      legende/.style = {font=\scriptsize,text=insagris,align=center,text width=1.9cm}}
    \pic[insagris,scale=1.4]   at (0,0.6)     {icone tel};
    \pic[insarouge,scale=1.4]  at (2.85,0.6)  {icone antenne};
    \pic[insaorange,scale=1.4] at (5.7,0.6)   {icone antenne};
    \pic[insarouge,scale=1.4]  at (8.55,0.6)  {icone antenne};
    \pic[insagris,scale=1.4]   at (11.4,0.6)  {icone tel};
    \draw[bt]   (0.5,0.6) -- node[etiqf,above=2pt] {Bluetooth} node[below=2pt,font=\tiny,text=insagris] {quelques mètres} (2.3,0.6);
    \draw[lora] (3.35,0.6) -- node[etiqf,above=3pt] {LoRa} node[below=3pt,font=\tiny,text=insagris] {plusieurs km} (5.2,0.6);
    \draw[lora] (6.2,0.6) -- node[etiqf,above=3pt] {LoRa} node[below=3pt,font=\tiny,text=insagris] {plusieurs km} (8.05,0.6);
    \draw[bt]   (9.05,0.6) -- node[etiqf,above=2pt] {Bluetooth} node[below=2pt,font=\tiny,text=insagris] {quelques mètres} (10.9,0.6);
    \node[pcli] at (0,-0.4)     {Toi};
    \node[pfoc] at (2.85,-0.4)  {Ton kit};
    \node[psvc] at (5.7,-0.4)   {Un voisin};
    \node[pfoc] at (8.55,-0.4)  {Kit du contact};
    \node[pcli] at (11.4,-0.4)  {Ton contact};
    \node[legende] at (0,-1.1)     {application mobile};
    \node[legende] at (2.85,-1.1)  {carte Heltec V3 + antenne};
    \node[legende] at (5.7,-1.1)   {relaie ton message};
    \node[legende] at (8.55,-1.1)  {reçoit et déchiffre};
    \node[legende] at (11.4,-1.1)  {application mobile};
  \end{tikzpicture}

  \vspace{0.2cm}
  \begin{columns}[T]
    \begin{column}{0.485\textwidth}
      \begin{encadre}{Partie physique~: ce que tu tiens en main}
        \begin{puces}
          \item Carte Heltec V3~: radio LoRa, Bluetooth, écran
          \item Antenne 868 MHz, batterie Li-Po (recharge USB-C)
          \item Boîtier de protection
        \end{puces}
      \end{encadre}
    \end{column}
    \begin{column}{0.485\textwidth}
      \begin{encadre}{Partie logicielle~: ce qui la fait vivre}
        \begin{puces}
          \item Firmware libre Meshtastic (GPL-3.0)~: maillage, chiffrement
          \item Application mobile gratuite~: messages, carte, positions
          \item Réglages préchargés, guides en français
        \end{puces}
      \end{encadre}
    \end{column}
  \end{columns}
  \srcnote{Sources~: \cite{heltec_v3,meshtastic_overview,meshtastic_android}}
\end{frame}

%---------------------------------------------------------------------------
% Diapositive 7 : partie physique (schéma-bloc de la carte + contenu du kit)
%---------------------------------------------------------------------------
\begin{frame}{Partie physique~: une carte de 5\,cm sur 2,5\,cm}
  \begin{columns}[T]
    \begin{column}{0.64\textwidth}
      \centering
      \begin{tikzpicture}[font=\scriptsize]
        \tikzset{lien/.style={{Stealth[length=1.8mm]}-{Stealth[length=1.8mm]},thick,insagris}}
        \node[svc,minimum width=2.5cm,minimum height=1.75cm,fill=insarose,line width=1.1pt] (mcu) at (0,0)
          {{\bfseries ESP32-S3}\\processeur 2 c\oe urs\\Bluetooth 5 LE, Wi-Fi\\8 Mo de flash};
        \node[svc,minimum width=1.9cm,minimum height=1.15cm] (lora) at (-3.15,0)
          {{\bfseries SX1262}\\radio LoRa\\863--928 MHz};
        \node[svcgris,minimum width=1.9cm,minimum height=1.15cm] (oled) at (3.15,0)
          {{\bfseries Écran OLED}\\0,96~pouce\\128$\times$64 points};
        \node[svcgris,minimum width=2.1cm,minimum height=0.95cm] (usb) at (-1.3,-2.0)
          {{\bfseries Port USB-C}\\charge, mises à jour};
        \node[svcgris,minimum width=2.1cm,minimum height=0.95cm] (bat) at (1.3,-2.0)
          {{\bfseries Batterie Li-Po}\\3,7\,V, puce de charge};
        \draw[lien] (lora) -- (mcu);
        \draw[lien] (mcu) -- (oled);
        \draw[lien] (usb.north) -- (usb.north |- mcu.south);
        \draw[lien] (bat.north) -- (bat.north |- mcu.south);
        % antenne : le signal radio sort par la puce SX1262
        \draw[thick,insagris] (lora.north) -- ++(0,0.45);
        \pic[insarouge,scale=1.2] at ([yshift=0.95cm]lora.north) {icone tour};
        \node[anchor=west,align=left,font=\scriptsize\bfseries,text=insarouge]
          at ([xshift=0.45cm,yshift=0.95cm]lora.north) {Antenne\\868 MHz};
        % téléphone : liaison Bluetooth
        \pic[insagris,scale=1.2] at (2.65,1.85) {icone tel};
        \node[anchor=west,font=\scriptsize\bfseries,text=insagris] at (2.98,1.85) {Téléphone};
        \draw[lien,dashed] ([xshift=0.8cm]mcu.north) -- (2.35,1.5)
          node[pos=0.45,etiqf,left=1pt] {Bluetooth};
        % portée annoncée par la communauté (fourchettes indicatives)
        \node[anchor=north,align=center,text width=1.75cm,font=\tiny,text=insagris]
          at ([yshift=-0.15cm]lora.south) {portée\\0,5 à 2\,km en ville\\5 à 15\,km en campagne};
      \end{tikzpicture}
    \end{column}
    \begin{column}{0.33\textwidth}
      \begin{encadre}{Dans la boîte}
        \begin{puces}
          \item Carte Heltec V3 préconfigurée
          \item Antenne 868 MHz, batterie Li-Po
          \item Boîtier imprimé en 3D (fichiers libres)
          \item Câble USB-C, guide de démarrage
        \end{puces}
      \end{encadre}
      \vspace{0.12cm}
      \begin{pointcle}\centering\footnotesize
        Composants standards du commerce, assemblage simple~: aucune pièce propriétaire.
      \end{pointcle}
    \end{column}
  \end{columns}
  \srcnote{Sources~: \cite{heltec_v3,heltec_docs,techconnect}}
\end{frame}

%---------------------------------------------------------------------------
% Diapositive 8 : partie logicielle (socle libre + notre contribution)
%---------------------------------------------------------------------------
\begin{frame}{Partie logicielle~: du code libre et vérifiable}
  \begin{columns}[T]
    \begin{column}{0.52\textwidth}
      \centering
      \begin{tikzpicture}[font=\scriptsize]
        \tikzset{couche/.style={minimum width=5.2cm,minimum height=0.8cm}}
        \node[psvc,couche]  (app) at (0,0.6)    {Application Meshtastic};
        \node[psvc,couche]  (fw)  at (0,-1.0)   {Firmware Meshtastic (GPL-3.0)};
        \node[pfoc,couche]  (cfg) at (0,-1.9)   {Réglages \nomproduit{} préchargés};
        \node[pgris,couche] (hw)  at (0,-2.8)   {Carte Heltec WiFi LoRa 32 V3};
        \draw[{Stealth[length=2mm]}-{Stealth[length=2mm]},thick,insagris,dashed] (app.south) -- (fw.north)
          node[midway,right=3pt,font=\scriptsize,text=insagris] {Bluetooth};
        \draw[decorate,decoration={brace,amplitude=4pt},insagris,thick] (-2.95,1.0) -- (-2.95,0.2)
          node[midway,left=5pt,rotate=90,anchor=south,font=\scriptsize\bfseries,text=insagris] {téléphone};
        \draw[decorate,decoration={brace,amplitude=4pt},insagris,thick] (-2.95,-0.6) -- (-2.95,-3.2)
          node[midway,left=5pt,rotate=90,anchor=south,font=\scriptsize\bfseries,text=insagris] {kit};
        \node[font=\tiny,text=insagris,anchor=north west] at (-2.6,-3.35)
          {\textcolor{insaorange}{$\blacksquare$} existant, libre\quad\textcolor{insarouge}{$\blacksquare$} notre apport\quad\textcolor{insagris}{$\blacksquare$} matériel};
      \end{tikzpicture}
    \end{column}
    \begin{column}{0.45\textwidth}
      \begin{encadre}{Ce que nous ajoutons}
        \begin{numeros}
          \item \textbf{Prêt à l'emploi}~: réglages Europe et canal privé à clé aléatoire (les clés par défaut sont publiques).
          \item \textbf{Prise en main simple}~: invitation d'un proche par QR code, guides en français.
          \item \textbf{Mises à jour testées} avant diffusion.
          \item \textbf{Contributions en amont}~: correctifs, traductions, documentation reversés au projet libre.
        \end{numeros}
      \end{encadre}
    \end{column}
  \end{columns}
  \vspace{-0.15cm}
  \begin{pointcle}\centering\footnotesize
    Aucun compte, aucun serveur central~: un code ouvert que chacun peut vérifier.
  \end{pointcle}
  \srcnote{Sources~: \cite{meshtastic_fw,meshtastic_android,meshtastic_crypto,meshtastic_channels,meshtastic_qr}}
\end{frame}

%---------------------------------------------------------------------------
% Diapositive 9 : trajet d'un message + limites assumées
%---------------------------------------------------------------------------
\begin{frame}{Trajet d'un message~: aucun opérateur sur la route}
  \centering
  \begin{tikzpicture}[font=\footnotesize]
    \node[acteur] (a0) at (0,0)    {Ton\\téléphone};
    \node[acteur,fill=insarouge!25] (a1) at (2.8,0)  {Ton\\kit};
    \node[acteur,fill=insapeche,draw=insaorange] (a2) at (5.6,0)  {Kit d'un\\voisin};
    \node[acteur,fill=insarouge!25] (a3) at (8.4,0)  {Kit du\\contact};
    \node[acteur] (a4) at (11.2,0) {Son\\téléphone};
    \foreach \i in {0,...,4} {\draw[vie] (a\i.south) -- ++(0,-3.1);}
    \seqmsg{0}{2.8}{-0.95}{1}{Bluetooth}
    \seqmsg{2.8}{5.6}{-1.5}{2}{LoRa, chiffré}
    \seqmsg{5.6}{8.4}{-2.05}{3}{Retransmis}
    \seqmsg{8.4}{11.2}{-2.6}{4}{Bluetooth}
    \seqrep{8.4}{2.8}{-3.15}{5}{accusé de réception, par le maillage}
  \end{tikzpicture}

  \vspace{0.15cm}
  \begin{columns}[T]
    \begin{column}{0.245\textwidth}
      \begin{carte}\centering\gros{239 kilooctets}\\[1pt]{\scriptsize par message~: du texte seulement}\end{carte} %239 octets
    \end{column}
    \begin{column}{0.245\textwidth}
      \begin{carte}\centering\gros{≈ 1 Mbit/s}\\[1pt]{\scriptsize débit du mode LongFast, portée maximale}\end{carte} % 1ko
    \end{column}
    \begin{column}{0.245\textwidth}
      \begin{carte}\centering\gros{10\,\%}\\[1pt]{\scriptsize du temps d'émission au plus, en Europe}\end{carte}
    \end{column}
    \begin{column}{0.245\textwidth}
      \begin{carte}\centering\gros{3 relais}\\[1pt]{\scriptsize par défaut (50 au maximum)}\end{carte} % 7 au max
    \end{column}
  \end{columns}
  \srcnote{Limites assumées. Sources~: \cite{meshtastic_overview,meshtastic_regions,meshtastic_lora,meshtastic_algo}}
\end{frame}
    % 2.  Fonctionnement : partie physique, partie logicielle
%%%%%%%%%%%%%%%%%%%%%%%%%%%%%%%%%%%%%%%%%%%%%%%%%%%%%%%%%%%%%%%%%%%%%%%%%%%%%%
%  Partie 4 -- Marché cible : trois premiers segments puis grand public, approche commerciale
%  (2 diapositives)
%%%%%%%%%%%%%%%%%%%%%%%%%%%%%%%%%%%%%%%%%%%%%%%%%%%%%%%%%%%%%%%%%%%%%%%%%%%%%%
\section{Marché cible}

%---------------------------------------------------------------------------
% Diapositive 13 : trois premiers segments, puis le grand public
%---------------------------------------------------------------------------
\begin{frame}{Trois segments d'abord, le grand public ensuite}
  \centering
  \hyphenpenalty=10000 \exhyphenpenalty=10000
  \begin{tikzpicture}[font=\footnotesize]
    \tikzset{seg/.style={fill=insabeige,rounded corners=3pt,text width=6.0cm,minimum height=1.45cm,
                         inner sep=6pt,anchor=west,align=left}}
    % trois segments de départ
    \node[seg] (s1) at (0.95,1.6)  {{\bfseries ONG et secours}~: guerre, séismes, inondations\\
                                   Coordonner des équipes quand les antennes-relais sont tombées ou coupées.};
    \node[seg] (s2) at (0.95,0)    {{\bfseries Randonneurs et aventuriers}\\
                                   Prévenir un proche, se retrouver, partager sa position en zone blanche.};
    \node[seg] (s3) at (0.95,-1.6) {{\bfseries Amateurs d'alternatives}~: high-tech, low-tech\\
                                   Posséder ses outils, se passer des plateformes~; déjà sensibles au libre.};
    \pic[insarouge,scale=1.2] at (0.45,1.6)  {icone tente};
    \pic[insarouge,scale=1.2] at (0.45,0)    {icone mont};
    \pic[insarouge,scale=1.2] at (0.45,-1.6) {icone code};
    % flèches vers le grand public
    \foreach \y in {1.6,0,-1.6} {\draw[fluxr] (7.05,\y) -- (8.55,0);}
    % grand public (phase 2)
    \node[fill=insarose,draw=insarouge,thick,rounded corners=4pt,text width=3.9cm,minimum height=4.7cm,
          inner sep=7pt,anchor=west,align=center] (gp) at (8.65,0)
      {{\sffamily\bfseries\color{insarouge}Phase 2\\Grand public}\\[4pt]
       Familles, voyageurs, habitants des zones rurales, sportifs d'extérieur.\\[5pt]
       {\gros{27 M}}\\{\scriptsize de Français pratiquent la randonnée}\\[5pt]
       {\scriptsize\bfseries\color{insarouge}Effet de réseau~: chaque kit renforce le maillage de tous}};
  \end{tikzpicture}
  \srcnote{Source~: \cite{ffrando}}
\end{frame}

%---------------------------------------------------------------------------
% Diapositive 14 : comment les atteindre (message, canaux) + programme solidaire
%---------------------------------------------------------------------------
\begin{frame}{Comment les atteindre~: un message par cible}
  \renewcommand{\arraystretch}{1.3}
  \footnotesize\hyphenpenalty=10000 \exhyphenpenalty=10000
  \noindent\begin{tabularx}{\textwidth}{@{}>{\bfseries}p{2.35cm}>{\raggedright\arraybackslash}X>{\raggedright\arraybackslash}X@{}}
    \rowcolor{insarose}
    \textcolor{insarouge}{Cible} & \textcolor{insarouge}{Message} & \textcolor{insarouge}{Canaux} \\
    ONG, secours & «\,Une messagerie qui survit aux infrastructures\,» &
      kits de démonstration, pilotes gratuits, réseaux de bénévoles, appels d'offres \\
    \arrayrulecolor{insagris!30}\hline
    Randonneurs & «\,Écrire à ses proches là où le réseau s'arrête\,» &
      clubs et fédérations, refuges, boutiques de sport, salons outdoor \\
    \hline
    Alternatives & «\,Ton réseau t'appartient\,» &
      forums du libre et des makers, financement participatif, ateliers (fablabs), vidéos \\
    \hline
    Grand public & «\,Un filet de sécurité à 69\,€, sans abonnement\,» &
      site marchand, distributeurs spécialisés, bouche-à-oreille \\
  \end{tabularx}

  \vspace{0.15cm}
  \begin{columns}[T]
    \begin{column}{0.485\textwidth}
      \begin{carte}[height=2.15cm,valign=top]
        \gros{1 kit offert pour 10 vendus}\\[2pt]
        {\footnotesize Programme solidaire pour les ONG~: un argument de marque et un moyen d'équiper des pilotes (coût déjà inclus dans les finances).}
      \end{carte}
    \end{column}
    \begin{column}{0.485\textwidth}
      \begin{carte}[height=2.15cm,valign=top]
        \gros{Précommandes d'abord}\\[2pt]
        {\footnotesize Un financement participatif valide la demande et finance le premier stock avant toute production en série.}
      \end{carte}
    \end{column}
  \end{columns}
\end{frame}
            % 3.  Marché cible : segments, taille, canaux
%%%%%%%%%%%%%%%%%%%%%%%%%%%%%%%%%%%%%%%%%%%%%%%%%%%%%%%%%%%%%%%%%%%%%%%%%%%%%%
%  Partie « Concurrence » (4 diapositives) -- à importer avec %%%%%%%%%%%%%%%%%%%%%%%%%%%%%%%%%%%%%%%%%%%%%%%%%%%%%%%%%%%%%%%%%%%%%%%%%%%%%%
%  Partie « Concurrence » (4 diapositives) -- à importer avec %%%%%%%%%%%%%%%%%%%%%%%%%%%%%%%%%%%%%%%%%%%%%%%%%%%%%%%%%%%%%%%%%%%%%%%%%%%%%%
%  Partie « Concurrence » (4 diapositives) -- à importer avec \include{chap_concurrence}
%  Fichier de contenu : ni préambule ni \begin{document}, commence par \section.
%
%  Contenu : (1) carte de positionnement, (2) comparatif détaillé, (3) coût sur trois ans,
%  (4) atouts et limites face à chaque famille de concurrents.
%  Cette partie REGROUPE les comparaisons de prix qui étaient dispersées (ancienne diapositive
%  « Contre l'abonnement » de la partie Marketing, tableau et cartes de l'étude documentaire).
%  Les prix sont des prix publics relevés du 5 au 8 octobre 2026 (voir sources.bib) : ils varient
%  selon le revendeur, d'où des fourchettes. Choix faits pour les fourchettes :
%   - talkies PMR446 : 65 à 80 € la paire, soit 33 à 40 € par personne ;
%   - kits Meshtastic prêts à l'emploi : de 59 € (SenseCAP T1000-E, 49 € HT) à 153 € (WisMesh Pocket V2) ;
%   - balise PLB : 300 à 400 € (pas d'abonnement) ;
%   - inReach Mini 2 : 329 à 349 € + abonnement ; 827 € sur trois ans (calcul sous la 3e diapositive).
%  Les coordonnées du graphique de la 1re diapositive sont en échelle logarithmique :
%  x = 3,1 + 4 × log10(coût / 10 €) en cm (10 € = origine, 100 € = 7,1 cm).
%%%%%%%%%%%%%%%%%%%%%%%%%%%%%%%%%%%%%%%%%%%%%%%%%%%%%%%%%%%%%%%%%%%%%%%%%%%%%%
\section{Concurrence}

%---------------------------------------------------------------------------
% Diapositive 1 : carte de positionnement (qui peut-on joindre, pour quel coût ?)
% Abscisse : coût sur 3 ans par personne équipée, échelle logarithmique.
% Ordonnée : trois niveaux de « portée du contact ».
%---------------------------------------------------------------------------
\begin{frame}{Positionnement~: qui peut-on joindre, pour quel coût~?}
  \centering
  \begin{tikzpicture}[font=\scriptsize,x=1cm,y=1cm]
    % bandes (une sur deux en beige) et libellés des trois niveaux
    \fill[insabeige] (3.0,0.4) rectangle (11.6,1.5);
    \fill[insabeige] (3.0,3.2) rectangle (11.6,4.5);
    \draw[insagris!40,thin] (3.0,0.4) rectangle (11.6,4.5);
    \node[anchor=east,align=right,font=\scriptsize\bfseries] at (2.9,3.85) {n'importe qui,\\les secours};
    \node[anchor=east,align=right,font=\scriptsize\bfseries] at (2.9,2.35) {tout le maillage\\(7 sauts au plus)};
    \node[anchor=east,align=right,font=\scriptsize\bfseries] at (2.9,0.95) {ceux qui sont\\à portée radio};
    % axe des coûts
    \draw[insagris,thick,-{Stealth[length=1.8mm]}] (3.0,0.4) -- (11.75,0.4);
    \foreach \x/\t in {3.1/0\,€,5.01/30\,€,7.1/100\,€,9.01/300\,€,11.1/1\,000\,€} {
      \draw[insagris] (\x,0.4) -- (\x,0.32);
      \node[anchor=north,font=\tiny,text=insagris] at (\x,0.32) {\t};}
    \node[anchor=north,font=\scriptsize,text=insagris] at (7.3,-0.1) {coût sur trois ans par personne équipée (échelle logarithmique)};

    % --- niveau « n'importe qui » : satellite (cercle vide = SOS seul, pas de messages)
    \draw[insaorange,thick,fill=white] (3.1,3.85) circle (0.14);
    \node[anchor=west,align=left] at (3.35,3.85) {SOS du smartphone~: iPhone 14 et plus\\(gratuit deux ans, déjà dans la poche)};
    \draw[insaorange,thick,fill=white] (9.28,3.85) circle (0.14);
    \node[anchor=east,align=right] at (9.1,3.85) {balise PLB\\300 à 400\,€};
    \fill[insaorange] (10.78,3.85) circle (0.17);
    \node[anchor=north,align=center] at (10.78,3.62) {inReach Mini 2\\≈ 830\,€};

    % --- niveau « tout le maillage » : Meshtastic
    \draw[insarouge!55,line width=4pt,line cap=round] (4.12,2.25) -- (5.01,2.25);
    \node[anchor=north,align=center] at (4.56,2.08) {carte nue\\18 à 30\,€};
    \draw[insarouge!55,line width=4pt,line cap=round] (6.19,2.25) -- (7.84,2.25);
    \node[anchor=north,align=center] at (7.0,2.08) {kits prêts à l'emploi\\59 à 153\,€};
    \fill[insarouge] (6.46,2.25) circle (0.17);
    \draw[insarouge,thick] (6.46,2.25) circle (0.27);
    \node[anchor=south,font=\footnotesize\bfseries,text=insarouge] at (6.46,2.58) {\nomproduit\ 69\,€};
    \node[callout,anchor=center,text width=2.6cm,align=left,font=\scriptsize] at (10.1,2.35)
      {Notre créneau~: un groupe à portée, sans abonnement, prêt à l'emploi.};

    % --- niveau « à portée radio » : talkies PMR446
    \draw[insagris,line width=4pt,line cap=round] (4.89,0.95) -- (5.51,0.95);
    \node[anchor=west,align=left] at (5.8,0.95) {talkies PMR446\\33 à 40\,€ par personne};

    % légende
    \begin{scope}[shift={(3.0,-0.85)},font=\tiny]
      \fill[insarouge] (0.1,0) circle (0.08);\node[anchor=west] at (0.22,0) {Meshtastic (carte, kits)};
      \fill[insagris] (3.0,0) circle (0.08);\node[anchor=west] at (3.12,0) {radio libre sans maillage};
      \fill[insaorange] (5.95,0) circle (0.08);\node[anchor=west] at (6.07,0) {satellite};
      \draw[insaorange,thick,fill=white] (7.4,0) circle (0.08);\node[anchor=west] at (7.52,0) {SOS seul (pas de messages)};
    \end{scope}
  \end{tikzpicture}
  \srcnote{Prix publics relevés en octobre 2026. Sources~: \cite{irun_mini2,garmin_plans,plb_bateaux,plb_toplicht,pmr_decathlon,pmr_officeeasy,t1000e_passion,wismesh_openelab,apple_sat,techconnect}}
\end{frame}

%---------------------------------------------------------------------------
% Diapositive 2 : comparatif détaillé
%---------------------------------------------------------------------------
\begin{frame}{Comparatif~: prix, abonnement, ce que chacun permet}
  \centering\scriptsize\renewcommand{\arraystretch}{1.22}\setlength{\tabcolsep}{3pt}
  \hyphenpenalty=10000 \exhyphenpenalty=10000
  \noindent\begin{tabular}{@{}>{\raggedright\arraybackslash}p{2.8cm}>{\raggedright\arraybackslash}p{1.5cm}>{\raggedright\arraybackslash}p{1.55cm}>{\centering\arraybackslash}p{1.3cm}>{\centering\arraybackslash}p{1.4cm}>{\centering\arraybackslash}p{1.4cm}>{\raggedright\arraybackslash}p{2.35cm}@{}}
    \rowcolor{insarose}\textcolor{insarouge}{\bfseries Offre} & \textcolor{insarouge}{\bfseries Achat} & \textcolor{insarouge}{\bfseries Abonnement} & \textcolor{insarouge}{\bfseries Écrire à son groupe} & \textcolor{insarouge}{\bfseries Joindre un proche hors zone} & \textcolor{insarouge}{\bfseries Appeler les secours} & \textcolor{insarouge}{\bfseries Portée utile}\\
    \rowcolor{insabeige}\textbf{\nomproduit} (kit prêt à l'emploi) & 69\,€ & non & \oui & \non & \non & 0,5 à 2\,km en ville, 5 à 15\,km en campagne\\
    Carte Heltec V3 nue (à régler soi-même) & 18 à 30\,€ & non & \oui & \non & \non & idem\\
    Meshtastic prêt à l'emploi (T1000-E, WisMesh) & 59 à 153\,€ & non & \oui & \non & \non & idem\\
    Talkies PMR446 (la paire) & 65 à 80\,€ & non & voix seule & \non & \non & 0,2 à 1\,km en ville, 3 à 6\,km dégagé\\
    Garmin inReach Mini 2 & 329 à 349\,€ & ≈ 13\,€ par mois & \oui & \oui & \oui, 24\,h/24 & mondiale, ciel dégagé\\
    Balise de détresse PLB 406\,MHz & 300 à 400\,€ & non & \non & \non & \oui, SOS seul & mondiale\\
    iPhone 14 ou plus (SOS par satellite) & 0\,€ si déjà équipé & gratuit deux ans & \non & \non & \oui, SOS seul & mondiale, ciel dégagé\\
  \end{tabular}
  \vspace{0.12cm}
  \begin{pcompact}\centering\footnotesize
    Le kit est un outil de groupe, pas un secours~: il ne remplace ni le 112, ni une balise, ni le téléphone.
  \end{pcompact}
  \srcnote{Prix relevés en octobre 2026, variables selon le revendeur. Les messages par satellite d'Apple ne sont pas disponibles en France. Sources~: \cite{irun_mini2,irun_msgplus,garmin_plans,plb_bateaux,plb_toplicht,pmr_decathlon,pmr_officeeasy,pmr_hellopro,t1000e_passion,wismesh_hexaspot,wismesh_openelab,apple_sat,apple_msg,techconnect}}
\end{frame}

%---------------------------------------------------------------------------
% Diapositive 3 : coût cumulé sur trois ans (ancienne diapositive « Contre l'abonnement »)
%---------------------------------------------------------------------------
\begin{frame}{Contre l'abonnement~: le coût sur trois ans}
  \begin{columns}[T]
    \begin{column}{0.62\textwidth}
      \centering
      \begin{tikzpicture}
        \begin{axis}[
            width=7.3cm,height=4.5cm,xmin=0,xmax=36,ymin=0,ymax=1000,
            xtick={0,12,24,36},xticklabels={départ,1 an,2 ans,3 ans},
            ytick={0,250,500,750,1000},
            tick label style={font=\scriptsize},
            ylabel={coût cumulé (€)},ylabel style={font=\scriptsize,yshift=-3pt},
            axis lines=left,grid=major,grid style={insagris!25,thin},
            legend style={at={(0.02,1.0)},anchor=north west,font=\scriptsize,draw=none,fill=none,
                          legend cell align=left},
            clip=false]
          \addplot[insagris,very thick] coordinates {(0,0) (36,643.68)};
          \addlegendentry{Forfait mobile moyen}
          \addplot[insaorange,very thick] coordinates {(0,363.3) (36,827.4)};
          \addlegendentry{Une balise + abonnement}
          \addplot[insarouge,line width=2.2pt,dashed] coordinates {(0,276) (36,276)};
          \addlegendentry{4 kits \nomproduit{}}
          \addplot[insarouge,line width=2.2pt] coordinates {(0,138) (36,138)};
          \addlegendentry{2 kits \nomproduit{}}
          \node[anchor=west,font=\scriptsize\bfseries,text=insagris] at (axis cs:36.6,644) {644\,€};
          \node[anchor=west,font=\scriptsize\bfseries,text=insaorange] at (axis cs:36.6,827) {827\,€};
          \node[anchor=west,font=\scriptsize\bfseries,text=insarouge] at (axis cs:36.6,290) {276\,€};
          \node[anchor=west,font=\scriptsize\bfseries,text=insarouge] at (axis cs:36.6,120) {138\,€};
        \end{axis}
      \end{tikzpicture}
    \end{column}
    \begin{column}{0.35\textwidth}
      \begin{remarque}[Ce que le graphique ne dit pas]
        \begin{puces}
          \item Le kit ne remplace pas le téléphone~: ni appels, ni internet.
          \item Aucun centre de secours n'est relié~: ce n'est pas une balise de détresse.
          \item Les contacts doivent être équipés aussi~: d'où les lots de 2 à 4 kits.
        \end{puces}
      \end{remarque}
    \end{column}
  \end{columns}
  \vspace{0.05cm}
  \begin{pointcle}\centering\footnotesize
    138\,€ pour deux kits, une fois~; 827\,€ sur trois ans pour une seule balise, qui elle relie aux secours.
  \end{pointcle}
  \srcnote{Hypothèses~: 36 mois, 1\,\$ = 0,86\,€, hors remplacement de batterie. Balise~: inReach Mini 2 (328,90\,€), forfait 14,99\,\$ par mois, activation 39,99\,\$. Forfait mobile moyen~: 17,88\,€ par mois. Sources~: \cite{arcep_prix,garmin_mini2,garmin_plans}}
\end{frame}

%---------------------------------------------------------------------------
% Diapositive 4 : notre avantage face à chaque famille, et ce qu'il faut leur reconnaître
%---------------------------------------------------------------------------
\begin{frame}{Face à chaque concurrent~: atouts et limites}
  \footnotesize\renewcommand{\arraystretch}{1.25}\setlength{\tabcolsep}{4pt}
  \hyphenpenalty=10000 \exhyphenpenalty=10000
  \noindent\begin{tabularx}{\textwidth}{@{}>{\raggedright\arraybackslash\bfseries}p{2.9cm}>{\raggedright\arraybackslash}X>{\raggedright\arraybackslash}X@{}}
    \rowcolor{insarose}\textcolor{insarouge}{Face à…} & \textcolor{insarouge}{Notre avantage} & \textcolor{insarouge}{Leur avantage}\\
    le fait-maison (carte nue, 18 à 30\,€) & Prêt à l'emploi~: réglages Europe, clé privée, boîtier, guide en français, garantie. & Prix plus bas, matériel au choix, aucune attente de livraison.\\
    \arrayrulecolor{insagris!30}\hline
    les kits prêts à l'emploi (59 à 153\,€) & Lots de 2 à 4 kits préconfigurés ensemble, guides et support en français, assemblage en France. & Marque établie, GPS, boîtier étanche (IP65 pour le T1000-E), parfois moins cher (59\,€).\\
    \hline
    les talkies PMR446 (65 à 80\,€ la paire) & Messages et positions sur carte, relais par les voisins, chiffrement. & Voix immédiate, sans smartphone ni réglage, usage libre.\\
    \hline
    le satellite (inReach, PLB, SOS iPhone) & Aucun abonnement~; marche sous les arbres et en ville, sans ciel dégagé. & SOS relié aux secours 24\,h/24, couverture mondiale, joignent tout le monde.\\
  \end{tabularx}
  \vspace{0.15cm}
  \begin{pcompact}\centering\footnotesize
    À mesurer dans l'enquête~: ce que les futurs clients emportent déjà (talkie, balise, rien) et ce qu'ils jugeraient suffisant.
  \end{pcompact}
  \srcnote{Sources~: \cite{pmr_hellopro,apple_sat,irun_msgplus,seeed_t1000e,wismesh_openelab}. Le satellite exige une vue dégagée du ciel (notices Apple et Garmin).}
\end{frame}

%  Fichier de contenu : ni préambule ni \begin{document}, commence par \section.
%
%  Contenu : (1) carte de positionnement, (2) comparatif détaillé, (3) coût sur trois ans,
%  (4) atouts et limites face à chaque famille de concurrents.
%  Cette partie REGROUPE les comparaisons de prix qui étaient dispersées (ancienne diapositive
%  « Contre l'abonnement » de la partie Marketing, tableau et cartes de l'étude documentaire).
%  Les prix sont des prix publics relevés du 5 au 8 octobre 2026 (voir sources.bib) : ils varient
%  selon le revendeur, d'où des fourchettes. Choix faits pour les fourchettes :
%   - talkies PMR446 : 65 à 80 € la paire, soit 33 à 40 € par personne ;
%   - kits Meshtastic prêts à l'emploi : de 59 € (SenseCAP T1000-E, 49 € HT) à 153 € (WisMesh Pocket V2) ;
%   - balise PLB : 300 à 400 € (pas d'abonnement) ;
%   - inReach Mini 2 : 329 à 349 € + abonnement ; 827 € sur trois ans (calcul sous la 3e diapositive).
%  Les coordonnées du graphique de la 1re diapositive sont en échelle logarithmique :
%  x = 3,1 + 4 × log10(coût / 10 €) en cm (10 € = origine, 100 € = 7,1 cm).
%%%%%%%%%%%%%%%%%%%%%%%%%%%%%%%%%%%%%%%%%%%%%%%%%%%%%%%%%%%%%%%%%%%%%%%%%%%%%%
\section{Concurrence}

%---------------------------------------------------------------------------
% Diapositive 1 : carte de positionnement (qui peut-on joindre, pour quel coût ?)
% Abscisse : coût sur 3 ans par personne équipée, échelle logarithmique.
% Ordonnée : trois niveaux de « portée du contact ».
%---------------------------------------------------------------------------
\begin{frame}{Positionnement~: qui peut-on joindre, pour quel coût~?}
  \centering
  \begin{tikzpicture}[font=\scriptsize,x=1cm,y=1cm]
    % bandes (une sur deux en beige) et libellés des trois niveaux
    \fill[insabeige] (3.0,0.4) rectangle (11.6,1.5);
    \fill[insabeige] (3.0,3.2) rectangle (11.6,4.5);
    \draw[insagris!40,thin] (3.0,0.4) rectangle (11.6,4.5);
    \node[anchor=east,align=right,font=\scriptsize\bfseries] at (2.9,3.85) {n'importe qui,\\les secours};
    \node[anchor=east,align=right,font=\scriptsize\bfseries] at (2.9,2.35) {tout le maillage\\(7 sauts au plus)};
    \node[anchor=east,align=right,font=\scriptsize\bfseries] at (2.9,0.95) {ceux qui sont\\à portée radio};
    % axe des coûts
    \draw[insagris,thick,-{Stealth[length=1.8mm]}] (3.0,0.4) -- (11.75,0.4);
    \foreach \x/\t in {3.1/0\,€,5.01/30\,€,7.1/100\,€,9.01/300\,€,11.1/1\,000\,€} {
      \draw[insagris] (\x,0.4) -- (\x,0.32);
      \node[anchor=north,font=\tiny,text=insagris] at (\x,0.32) {\t};}
    \node[anchor=north,font=\scriptsize,text=insagris] at (7.3,-0.1) {coût sur trois ans par personne équipée (échelle logarithmique)};

    % --- niveau « n'importe qui » : satellite (cercle vide = SOS seul, pas de messages)
    \draw[insaorange,thick,fill=white] (3.1,3.85) circle (0.14);
    \node[anchor=west,align=left] at (3.35,3.85) {SOS du smartphone~: iPhone 14 et plus\\(gratuit deux ans, déjà dans la poche)};
    \draw[insaorange,thick,fill=white] (9.28,3.85) circle (0.14);
    \node[anchor=east,align=right] at (9.1,3.85) {balise PLB\\300 à 400\,€};
    \fill[insaorange] (10.78,3.85) circle (0.17);
    \node[anchor=north,align=center] at (10.78,3.62) {inReach Mini 2\\≈ 830\,€};

    % --- niveau « tout le maillage » : Meshtastic
    \draw[insarouge!55,line width=4pt,line cap=round] (4.12,2.25) -- (5.01,2.25);
    \node[anchor=north,align=center] at (4.56,2.08) {carte nue\\18 à 30\,€};
    \draw[insarouge!55,line width=4pt,line cap=round] (6.19,2.25) -- (7.84,2.25);
    \node[anchor=north,align=center] at (7.0,2.08) {kits prêts à l'emploi\\59 à 153\,€};
    \fill[insarouge] (6.46,2.25) circle (0.17);
    \draw[insarouge,thick] (6.46,2.25) circle (0.27);
    \node[anchor=south,font=\footnotesize\bfseries,text=insarouge] at (6.46,2.58) {\nomproduit\ 69\,€};
    \node[callout,anchor=center,text width=2.6cm,align=left,font=\scriptsize] at (10.1,2.35)
      {Notre créneau~: un groupe à portée, sans abonnement, prêt à l'emploi.};

    % --- niveau « à portée radio » : talkies PMR446
    \draw[insagris,line width=4pt,line cap=round] (4.89,0.95) -- (5.51,0.95);
    \node[anchor=west,align=left] at (5.8,0.95) {talkies PMR446\\33 à 40\,€ par personne};

    % légende
    \begin{scope}[shift={(3.0,-0.85)},font=\tiny]
      \fill[insarouge] (0.1,0) circle (0.08);\node[anchor=west] at (0.22,0) {Meshtastic (carte, kits)};
      \fill[insagris] (3.0,0) circle (0.08);\node[anchor=west] at (3.12,0) {radio libre sans maillage};
      \fill[insaorange] (5.95,0) circle (0.08);\node[anchor=west] at (6.07,0) {satellite};
      \draw[insaorange,thick,fill=white] (7.4,0) circle (0.08);\node[anchor=west] at (7.52,0) {SOS seul (pas de messages)};
    \end{scope}
  \end{tikzpicture}
  \srcnote{Prix publics relevés en octobre 2026. Sources~: \cite{irun_mini2,garmin_plans,plb_bateaux,plb_toplicht,pmr_decathlon,pmr_officeeasy,t1000e_passion,wismesh_openelab,apple_sat,techconnect}}
\end{frame}

%---------------------------------------------------------------------------
% Diapositive 2 : comparatif détaillé
%---------------------------------------------------------------------------
\begin{frame}{Comparatif~: prix, abonnement, ce que chacun permet}
  \centering\scriptsize\renewcommand{\arraystretch}{1.22}\setlength{\tabcolsep}{3pt}
  \hyphenpenalty=10000 \exhyphenpenalty=10000
  \noindent\begin{tabular}{@{}>{\raggedright\arraybackslash}p{2.8cm}>{\raggedright\arraybackslash}p{1.5cm}>{\raggedright\arraybackslash}p{1.55cm}>{\centering\arraybackslash}p{1.3cm}>{\centering\arraybackslash}p{1.4cm}>{\centering\arraybackslash}p{1.4cm}>{\raggedright\arraybackslash}p{2.35cm}@{}}
    \rowcolor{insarose}\textcolor{insarouge}{\bfseries Offre} & \textcolor{insarouge}{\bfseries Achat} & \textcolor{insarouge}{\bfseries Abonnement} & \textcolor{insarouge}{\bfseries Écrire à son groupe} & \textcolor{insarouge}{\bfseries Joindre un proche hors zone} & \textcolor{insarouge}{\bfseries Appeler les secours} & \textcolor{insarouge}{\bfseries Portée utile}\\
    \rowcolor{insabeige}\textbf{\nomproduit} (kit prêt à l'emploi) & 69\,€ & non & \oui & \non & \non & 0,5 à 2\,km en ville, 5 à 15\,km en campagne\\
    Carte Heltec V3 nue (à régler soi-même) & 18 à 30\,€ & non & \oui & \non & \non & idem\\
    Meshtastic prêt à l'emploi (T1000-E, WisMesh) & 59 à 153\,€ & non & \oui & \non & \non & idem\\
    Talkies PMR446 (la paire) & 65 à 80\,€ & non & voix seule & \non & \non & 0,2 à 1\,km en ville, 3 à 6\,km dégagé\\
    Garmin inReach Mini 2 & 329 à 349\,€ & ≈ 13\,€ par mois & \oui & \oui & \oui, 24\,h/24 & mondiale, ciel dégagé\\
    Balise de détresse PLB 406\,MHz & 300 à 400\,€ & non & \non & \non & \oui, SOS seul & mondiale\\
    iPhone 14 ou plus (SOS par satellite) & 0\,€ si déjà équipé & gratuit deux ans & \non & \non & \oui, SOS seul & mondiale, ciel dégagé\\
  \end{tabular}
  \vspace{0.12cm}
  \begin{pcompact}\centering\footnotesize
    Le kit est un outil de groupe, pas un secours~: il ne remplace ni le 112, ni une balise, ni le téléphone.
  \end{pcompact}
  \srcnote{Prix relevés en octobre 2026, variables selon le revendeur. Les messages par satellite d'Apple ne sont pas disponibles en France. Sources~: \cite{irun_mini2,irun_msgplus,garmin_plans,plb_bateaux,plb_toplicht,pmr_decathlon,pmr_officeeasy,pmr_hellopro,t1000e_passion,wismesh_hexaspot,wismesh_openelab,apple_sat,apple_msg,techconnect}}
\end{frame}

%---------------------------------------------------------------------------
% Diapositive 3 : coût cumulé sur trois ans (ancienne diapositive « Contre l'abonnement »)
%---------------------------------------------------------------------------
\begin{frame}{Contre l'abonnement~: le coût sur trois ans}
  \begin{columns}[T]
    \begin{column}{0.62\textwidth}
      \centering
      \begin{tikzpicture}
        \begin{axis}[
            width=7.3cm,height=4.5cm,xmin=0,xmax=36,ymin=0,ymax=1000,
            xtick={0,12,24,36},xticklabels={départ,1 an,2 ans,3 ans},
            ytick={0,250,500,750,1000},
            tick label style={font=\scriptsize},
            ylabel={coût cumulé (€)},ylabel style={font=\scriptsize,yshift=-3pt},
            axis lines=left,grid=major,grid style={insagris!25,thin},
            legend style={at={(0.02,1.0)},anchor=north west,font=\scriptsize,draw=none,fill=none,
                          legend cell align=left},
            clip=false]
          \addplot[insagris,very thick] coordinates {(0,0) (36,643.68)};
          \addlegendentry{Forfait mobile moyen}
          \addplot[insaorange,very thick] coordinates {(0,363.3) (36,827.4)};
          \addlegendentry{Une balise + abonnement}
          \addplot[insarouge,line width=2.2pt,dashed] coordinates {(0,276) (36,276)};
          \addlegendentry{4 kits \nomproduit{}}
          \addplot[insarouge,line width=2.2pt] coordinates {(0,138) (36,138)};
          \addlegendentry{2 kits \nomproduit{}}
          \node[anchor=west,font=\scriptsize\bfseries,text=insagris] at (axis cs:36.6,644) {644\,€};
          \node[anchor=west,font=\scriptsize\bfseries,text=insaorange] at (axis cs:36.6,827) {827\,€};
          \node[anchor=west,font=\scriptsize\bfseries,text=insarouge] at (axis cs:36.6,290) {276\,€};
          \node[anchor=west,font=\scriptsize\bfseries,text=insarouge] at (axis cs:36.6,120) {138\,€};
        \end{axis}
      \end{tikzpicture}
    \end{column}
    \begin{column}{0.35\textwidth}
      \begin{remarque}[Ce que le graphique ne dit pas]
        \begin{puces}
          \item Le kit ne remplace pas le téléphone~: ni appels, ni internet.
          \item Aucun centre de secours n'est relié~: ce n'est pas une balise de détresse.
          \item Les contacts doivent être équipés aussi~: d'où les lots de 2 à 4 kits.
        \end{puces}
      \end{remarque}
    \end{column}
  \end{columns}
  \vspace{0.05cm}
  \begin{pointcle}\centering\footnotesize
    138\,€ pour deux kits, une fois~; 827\,€ sur trois ans pour une seule balise, qui elle relie aux secours.
  \end{pointcle}
  \srcnote{Hypothèses~: 36 mois, 1\,\$ = 0,86\,€, hors remplacement de batterie. Balise~: inReach Mini 2 (328,90\,€), forfait 14,99\,\$ par mois, activation 39,99\,\$. Forfait mobile moyen~: 17,88\,€ par mois. Sources~: \cite{arcep_prix,garmin_mini2,garmin_plans}}
\end{frame}

%---------------------------------------------------------------------------
% Diapositive 4 : notre avantage face à chaque famille, et ce qu'il faut leur reconnaître
%---------------------------------------------------------------------------
\begin{frame}{Face à chaque concurrent~: atouts et limites}
  \footnotesize\renewcommand{\arraystretch}{1.25}\setlength{\tabcolsep}{4pt}
  \hyphenpenalty=10000 \exhyphenpenalty=10000
  \noindent\begin{tabularx}{\textwidth}{@{}>{\raggedright\arraybackslash\bfseries}p{2.9cm}>{\raggedright\arraybackslash}X>{\raggedright\arraybackslash}X@{}}
    \rowcolor{insarose}\textcolor{insarouge}{Face à…} & \textcolor{insarouge}{Notre avantage} & \textcolor{insarouge}{Leur avantage}\\
    le fait-maison (carte nue, 18 à 30\,€) & Prêt à l'emploi~: réglages Europe, clé privée, boîtier, guide en français, garantie. & Prix plus bas, matériel au choix, aucune attente de livraison.\\
    \arrayrulecolor{insagris!30}\hline
    les kits prêts à l'emploi (59 à 153\,€) & Lots de 2 à 4 kits préconfigurés ensemble, guides et support en français, assemblage en France. & Marque établie, GPS, boîtier étanche (IP65 pour le T1000-E), parfois moins cher (59\,€).\\
    \hline
    les talkies PMR446 (65 à 80\,€ la paire) & Messages et positions sur carte, relais par les voisins, chiffrement. & Voix immédiate, sans smartphone ni réglage, usage libre.\\
    \hline
    le satellite (inReach, PLB, SOS iPhone) & Aucun abonnement~; marche sous les arbres et en ville, sans ciel dégagé. & SOS relié aux secours 24\,h/24, couverture mondiale, joignent tout le monde.\\
  \end{tabularx}
  \vspace{0.15cm}
  \begin{pcompact}\centering\footnotesize
    À mesurer dans l'enquête~: ce que les futurs clients emportent déjà (talkie, balise, rien) et ce qu'ils jugeraient suffisant.
  \end{pcompact}
  \srcnote{Sources~: \cite{pmr_hellopro,apple_sat,irun_msgplus,seeed_t1000e,wismesh_openelab}. Le satellite exige une vue dégagée du ciel (notices Apple et Garmin).}
\end{frame}

%  Fichier de contenu : ni préambule ni \begin{document}, commence par \section.
%
%  Contenu : (1) carte de positionnement, (2) comparatif détaillé, (3) coût sur trois ans,
%  (4) atouts et limites face à chaque famille de concurrents.
%  Cette partie REGROUPE les comparaisons de prix qui étaient dispersées (ancienne diapositive
%  « Contre l'abonnement » de la partie Marketing, tableau et cartes de l'étude documentaire).
%  Les prix sont des prix publics relevés du 5 au 8 octobre 2026 (voir sources.bib) : ils varient
%  selon le revendeur, d'où des fourchettes. Choix faits pour les fourchettes :
%   - talkies PMR446 : 65 à 80 € la paire, soit 33 à 40 € par personne ;
%   - kits Meshtastic prêts à l'emploi : de 59 € (SenseCAP T1000-E, 49 € HT) à 153 € (WisMesh Pocket V2) ;
%   - balise PLB : 300 à 400 € (pas d'abonnement) ;
%   - inReach Mini 2 : 329 à 349 € + abonnement ; 827 € sur trois ans (calcul sous la 3e diapositive).
%  Les coordonnées du graphique de la 1re diapositive sont en échelle logarithmique :
%  x = 3,1 + 4 × log10(coût / 10 €) en cm (10 € = origine, 100 € = 7,1 cm).
%%%%%%%%%%%%%%%%%%%%%%%%%%%%%%%%%%%%%%%%%%%%%%%%%%%%%%%%%%%%%%%%%%%%%%%%%%%%%%
\section{Concurrence}

%---------------------------------------------------------------------------
% Diapositive 1 : carte de positionnement (qui peut-on joindre, pour quel coût ?)
% Abscisse : coût sur 3 ans par personne équipée, échelle logarithmique.
% Ordonnée : trois niveaux de « portée du contact ».
%---------------------------------------------------------------------------
\begin{frame}{Positionnement~: qui peut-on joindre, pour quel coût~?}
  \centering
  \begin{tikzpicture}[font=\scriptsize,x=1cm,y=1cm]
    % bandes (une sur deux en beige) et libellés des trois niveaux
    \fill[insabeige] (3.0,0.4) rectangle (11.6,1.5);
    \fill[insabeige] (3.0,3.2) rectangle (11.6,4.5);
    \draw[insagris!40,thin] (3.0,0.4) rectangle (11.6,4.5);
    \node[anchor=east,align=right,font=\scriptsize\bfseries] at (2.9,3.85) {n'importe qui,\\les secours};
    \node[anchor=east,align=right,font=\scriptsize\bfseries] at (2.9,2.35) {tout le maillage\\(7 sauts au plus)};
    \node[anchor=east,align=right,font=\scriptsize\bfseries] at (2.9,0.95) {ceux qui sont\\à portée radio};
    % axe des coûts
    \draw[insagris,thick,-{Stealth[length=1.8mm]}] (3.0,0.4) -- (11.75,0.4);
    \foreach \x/\t in {3.1/0\,€,5.01/30\,€,7.1/100\,€,9.01/300\,€,11.1/1\,000\,€} {
      \draw[insagris] (\x,0.4) -- (\x,0.32);
      \node[anchor=north,font=\tiny,text=insagris] at (\x,0.32) {\t};}
    \node[anchor=north,font=\scriptsize,text=insagris] at (7.3,-0.1) {coût sur trois ans par personne équipée (échelle logarithmique)};

    % --- niveau « n'importe qui » : satellite (cercle vide = SOS seul, pas de messages)
    \draw[insaorange,thick,fill=white] (3.1,3.85) circle (0.14);
    \node[anchor=west,align=left] at (3.35,3.85) {SOS du smartphone~: iPhone 14 et plus\\(gratuit deux ans, déjà dans la poche)};
    \draw[insaorange,thick,fill=white] (9.28,3.85) circle (0.14);
    \node[anchor=east,align=right] at (9.1,3.85) {balise PLB\\300 à 400\,€};
    \fill[insaorange] (10.78,3.85) circle (0.17);
    \node[anchor=north,align=center] at (10.78,3.62) {inReach Mini 2\\≈ 830\,€};

    % --- niveau « tout le maillage » : Meshtastic
    \draw[insarouge!55,line width=4pt,line cap=round] (4.12,2.25) -- (5.01,2.25);
    \node[anchor=north,align=center] at (4.56,2.08) {carte nue\\18 à 30\,€};
    \draw[insarouge!55,line width=4pt,line cap=round] (6.19,2.25) -- (7.84,2.25);
    \node[anchor=north,align=center] at (7.0,2.08) {kits prêts à l'emploi\\59 à 153\,€};
    \fill[insarouge] (6.46,2.25) circle (0.17);
    \draw[insarouge,thick] (6.46,2.25) circle (0.27);
    \node[anchor=south,font=\footnotesize\bfseries,text=insarouge] at (6.46,2.58) {\nomproduit\ 69\,€};
    \node[callout,anchor=center,text width=2.6cm,align=left,font=\scriptsize] at (10.1,2.35)
      {Notre créneau~: un groupe à portée, sans abonnement, prêt à l'emploi.};

    % --- niveau « à portée radio » : talkies PMR446
    \draw[insagris,line width=4pt,line cap=round] (4.89,0.95) -- (5.51,0.95);
    \node[anchor=west,align=left] at (5.8,0.95) {talkies PMR446\\33 à 40\,€ par personne};

    % légende
    \begin{scope}[shift={(3.0,-0.85)},font=\tiny]
      \fill[insarouge] (0.1,0) circle (0.08);\node[anchor=west] at (0.22,0) {Meshtastic (carte, kits)};
      \fill[insagris] (3.0,0) circle (0.08);\node[anchor=west] at (3.12,0) {radio libre sans maillage};
      \fill[insaorange] (5.95,0) circle (0.08);\node[anchor=west] at (6.07,0) {satellite};
      \draw[insaorange,thick,fill=white] (7.4,0) circle (0.08);\node[anchor=west] at (7.52,0) {SOS seul (pas de messages)};
    \end{scope}
  \end{tikzpicture}
  \srcnote{Prix publics relevés en octobre 2026. Sources~: \cite{irun_mini2,garmin_plans,plb_bateaux,plb_toplicht,pmr_decathlon,pmr_officeeasy,t1000e_passion,wismesh_openelab,apple_sat,techconnect}}
\end{frame}

%---------------------------------------------------------------------------
% Diapositive 2 : comparatif détaillé
%---------------------------------------------------------------------------
\begin{frame}{Comparatif~: prix, abonnement, ce que chacun permet}
  \centering\scriptsize\renewcommand{\arraystretch}{1.22}\setlength{\tabcolsep}{3pt}
  \hyphenpenalty=10000 \exhyphenpenalty=10000
  \noindent\begin{tabular}{@{}>{\raggedright\arraybackslash}p{2.8cm}>{\raggedright\arraybackslash}p{1.5cm}>{\raggedright\arraybackslash}p{1.55cm}>{\centering\arraybackslash}p{1.3cm}>{\centering\arraybackslash}p{1.4cm}>{\centering\arraybackslash}p{1.4cm}>{\raggedright\arraybackslash}p{2.35cm}@{}}
    \rowcolor{insarose}\textcolor{insarouge}{\bfseries Offre} & \textcolor{insarouge}{\bfseries Achat} & \textcolor{insarouge}{\bfseries Abonnement} & \textcolor{insarouge}{\bfseries Écrire à son groupe} & \textcolor{insarouge}{\bfseries Joindre un proche hors zone} & \textcolor{insarouge}{\bfseries Appeler les secours} & \textcolor{insarouge}{\bfseries Portée utile}\\
    \rowcolor{insabeige}\textbf{\nomproduit} (kit prêt à l'emploi) & 69\,€ & non & \oui & \non & \non & 0,5 à 2\,km en ville, 5 à 15\,km en campagne\\
    Carte Heltec V3 nue (à régler soi-même) & 18 à 30\,€ & non & \oui & \non & \non & idem\\
    Meshtastic prêt à l'emploi (T1000-E, WisMesh) & 59 à 153\,€ & non & \oui & \non & \non & idem\\
    Talkies PMR446 (la paire) & 65 à 80\,€ & non & voix seule & \non & \non & 0,2 à 1\,km en ville, 3 à 6\,km dégagé\\
    Garmin inReach Mini 2 & 329 à 349\,€ & ≈ 13\,€ par mois & \oui & \oui & \oui, 24\,h/24 & mondiale, ciel dégagé\\
    Balise de détresse PLB 406\,MHz & 300 à 400\,€ & non & \non & \non & \oui, SOS seul & mondiale\\
    iPhone 14 ou plus (SOS par satellite) & 0\,€ si déjà équipé & gratuit deux ans & \non & \non & \oui, SOS seul & mondiale, ciel dégagé\\
  \end{tabular}
  \vspace{0.12cm}
  \begin{pcompact}\centering\footnotesize
    Le kit est un outil de groupe, pas un secours~: il ne remplace ni le 112, ni une balise, ni le téléphone.
  \end{pcompact}
  \srcnote{Prix relevés en octobre 2026, variables selon le revendeur. Les messages par satellite d'Apple ne sont pas disponibles en France. Sources~: \cite{irun_mini2,irun_msgplus,garmin_plans,plb_bateaux,plb_toplicht,pmr_decathlon,pmr_officeeasy,pmr_hellopro,t1000e_passion,wismesh_hexaspot,wismesh_openelab,apple_sat,apple_msg,techconnect}}
\end{frame}

%---------------------------------------------------------------------------
% Diapositive 3 : coût cumulé sur trois ans (ancienne diapositive « Contre l'abonnement »)
%---------------------------------------------------------------------------
\begin{frame}{Contre l'abonnement~: le coût sur trois ans}
  \begin{columns}[T]
    \begin{column}{0.62\textwidth}
      \centering
      \begin{tikzpicture}
        \begin{axis}[
            width=7.3cm,height=4.5cm,xmin=0,xmax=36,ymin=0,ymax=1000,
            xtick={0,12,24,36},xticklabels={départ,1 an,2 ans,3 ans},
            ytick={0,250,500,750,1000},
            tick label style={font=\scriptsize},
            ylabel={coût cumulé (€)},ylabel style={font=\scriptsize,yshift=-3pt},
            axis lines=left,grid=major,grid style={insagris!25,thin},
            legend style={at={(0.02,1.0)},anchor=north west,font=\scriptsize,draw=none,fill=none,
                          legend cell align=left},
            clip=false]
          \addplot[insagris,very thick] coordinates {(0,0) (36,643.68)};
          \addlegendentry{Forfait mobile moyen}
          \addplot[insaorange,very thick] coordinates {(0,363.3) (36,827.4)};
          \addlegendentry{Une balise + abonnement}
          \addplot[insarouge,line width=2.2pt,dashed] coordinates {(0,276) (36,276)};
          \addlegendentry{4 kits \nomproduit{}}
          \addplot[insarouge,line width=2.2pt] coordinates {(0,138) (36,138)};
          \addlegendentry{2 kits \nomproduit{}}
          \node[anchor=west,font=\scriptsize\bfseries,text=insagris] at (axis cs:36.6,644) {644\,€};
          \node[anchor=west,font=\scriptsize\bfseries,text=insaorange] at (axis cs:36.6,827) {827\,€};
          \node[anchor=west,font=\scriptsize\bfseries,text=insarouge] at (axis cs:36.6,290) {276\,€};
          \node[anchor=west,font=\scriptsize\bfseries,text=insarouge] at (axis cs:36.6,120) {138\,€};
        \end{axis}
      \end{tikzpicture}
    \end{column}
    \begin{column}{0.35\textwidth}
      \begin{remarque}[Ce que le graphique ne dit pas]
        \begin{puces}
          \item Le kit ne remplace pas le téléphone~: ni appels, ni internet.
          \item Aucun centre de secours n'est relié~: ce n'est pas une balise de détresse.
          \item Les contacts doivent être équipés aussi~: d'où les lots de 2 à 4 kits.
        \end{puces}
      \end{remarque}
    \end{column}
  \end{columns}
  \vspace{0.05cm}
  \begin{pointcle}\centering\footnotesize
    138\,€ pour deux kits, une fois~; 827\,€ sur trois ans pour une seule balise, qui elle relie aux secours.
  \end{pointcle}
  \srcnote{Hypothèses~: 36 mois, 1\,\$ = 0,86\,€, hors remplacement de batterie. Balise~: inReach Mini 2 (328,90\,€), forfait 14,99\,\$ par mois, activation 39,99\,\$. Forfait mobile moyen~: 17,88\,€ par mois. Sources~: \cite{arcep_prix,garmin_mini2,garmin_plans}}
\end{frame}

%---------------------------------------------------------------------------
% Diapositive 4 : notre avantage face à chaque famille, et ce qu'il faut leur reconnaître
%---------------------------------------------------------------------------
\begin{frame}{Face à chaque concurrent~: atouts et limites}
  \footnotesize\renewcommand{\arraystretch}{1.25}\setlength{\tabcolsep}{4pt}
  \hyphenpenalty=10000 \exhyphenpenalty=10000
  \noindent\begin{tabularx}{\textwidth}{@{}>{\raggedright\arraybackslash\bfseries}p{2.9cm}>{\raggedright\arraybackslash}X>{\raggedright\arraybackslash}X@{}}
    \rowcolor{insarose}\textcolor{insarouge}{Face à…} & \textcolor{insarouge}{Notre avantage} & \textcolor{insarouge}{Leur avantage}\\
    le fait-maison (carte nue, 18 à 30\,€) & Prêt à l'emploi~: réglages Europe, clé privée, boîtier, guide en français, garantie. & Prix plus bas, matériel au choix, aucune attente de livraison.\\
    \arrayrulecolor{insagris!30}\hline
    les kits prêts à l'emploi (59 à 153\,€) & Lots de 2 à 4 kits préconfigurés ensemble, guides et support en français, assemblage en France. & Marque établie, GPS, boîtier étanche (IP65 pour le T1000-E), parfois moins cher (59\,€).\\
    \hline
    les talkies PMR446 (65 à 80\,€ la paire) & Messages et positions sur carte, relais par les voisins, chiffrement. & Voix immédiate, sans smartphone ni réglage, usage libre.\\
    \hline
    le satellite (inReach, PLB, SOS iPhone) & Aucun abonnement~; marche sous les arbres et en ville, sans ciel dégagé. & SOS relié aux secours 24\,h/24, couverture mondiale, joignent tout le monde.\\
  \end{tabularx}
  \vspace{0.15cm}
  \begin{pcompact}\centering\footnotesize
    À mesurer dans l'enquête~: ce que les futurs clients emportent déjà (talkie, balise, rien) et ce qu'ils jugeraient suffisant.
  \end{pcompact}
  \srcnote{Sources~: \cite{pmr_hellopro,apple_sat,irun_msgplus,seeed_t1000e,wismesh_openelab}. Le satellite exige une vue dégagée du ciel (notices Apple et Garmin).}
\end{frame}
       % 4.  Concurrence : positionnement, comparatif, coût sur trois ans
%%%%%%%%%%%%%%%%%%%%%%%%%%%%%%%%%%%%%%%%%%%%%%%%%%%%%%%%%%%%%%%%%%%%%%%%%%%%%%
%  Partie 3 -- Marketing : promesses, coût face à l'abonnement, terrains d'usage
%  (3 diapositives)
%  Astuce de mise en page : les cartes à hauteur fixe se règlent avec carte[height=...,valign=top] ;
%  chaque icône est centrée dans un cadre invisible de 1,2 cm de haut pour aligner les titres.
%%%%%%%%%%%%%%%%%%%%%%%%%%%%%%%%%%%%%%%%%%%%%%%%%%%%%%%%%%%%%%%%%%%%%%%%%%%%%%
\section{Marketing}

%---------------------------------------------------------------------------
% Diapositive 10 : les trois promesses + slogans
%---------------------------------------------------------------------------
\begin{frame}{Trois promesses~: prix, autonomie, zones blanches}
  \begin{columns}[T]
    \begin{column}{0.325\textwidth}
      \begin{carte}[height=4.05cm,valign=top]
        \centering
        \begin{tikzpicture}\path (-0.7,-0.6) rectangle (0.7,0.6);\pic[insarouge,scale=1.7] at (0,0) {icone euro};\end{tikzpicture}\\[3pt]
        \gros{Petit prix}\\[3pt]
        {\footnotesize ≈ 69\,€ une seule fois. Aucun abonnement, aucune facture mensuelle, aucun frais caché.}
      \end{carte}
    \end{column}
    \begin{column}{0.325\textwidth}
      \begin{carte}[height=4.05cm,valign=top]
        \centering
        \begin{tikzpicture}\path (-0.7,-0.6) rectangle (0.7,0.6);\pic[insarouge,scale=1.9] at (0,0) {icone cadenas};\end{tikzpicture}\\[3pt]
        \gros{Autonomie}\\[3pt]
        {\footnotesize Ton antenne, ton logiciel libre, ta clé~: ni Orange, ni Free, ni SFR, ni réseau social.}
      \end{carte}
    \end{column}
    \begin{column}{0.325\textwidth}
      \begin{carte}[height=4.05cm,valign=top]
        \centering
        \begin{tikzpicture}\path (-0.7,-0.6) rectangle (0.7,0.6);\pic[scale=1.6] at (0,0) {icone mont};\end{tikzpicture}\\[3pt]
        \gros{Zones blanches}\\[3pt]
        {\footnotesize Campagne, montagne, zones de guerre, catastrophes naturelles~: là où le réseau manque ou tombe.}
      \end{carte}
    \end{column}
  \end{columns}
  \vspace{0.15cm}
  \begin{idee}\centering
    {\sffamily\bfseries\color{insarouge}Pas d'abonnement. Pas de compte. Pas de permission.}\\[3pt]
    {\sffamily\bfseries\color{insarouge}Là où le réseau s'arrête, le tien commence.}
  \end{idee}
\end{frame}

%---------------------------------------------------------------------------
% Diapositive 11 : coût cumulé face à l'abonnement
%---------------------------------------------------------------------------
\begin{frame}{Contre l'abonnement~: le coût sur trois ans}
  \begin{columns}[T]
    \begin{column}{0.62\textwidth}
      \centering
      \begin{tikzpicture}
        \begin{axis}[
            width=7.3cm,height=4.9cm,xmin=0,xmax=36,ymin=0,ymax=1180,
            xtick={0,12,24,36},xticklabels={départ,1 an,2 ans,3 ans},
            ytick={0,250,500,750,1000},
            tick label style={font=\scriptsize},
            ylabel={coût cumulé (€)},ylabel style={font=\scriptsize,yshift=-3pt},
            axis lines=left,grid=major,grid style={insagris!25,thin},
            legend style={at={(0.02,1.0)},anchor=north west,font=\scriptsize,draw=none,fill=none,
                          legend cell align=left},
            clip=false]
          \addplot[insagris,very thick] coordinates {(0,0) (36,643.68)};
          \addlegendentry{Forfait mobile moyen}
          \addplot[insaorange,very thick] coordinates {(0,363.3) (36,827.4)};
          \addlegendentry{Balise + abonnement}
          \addplot[insarouge,line width=2.2pt] coordinates {(0,69) (36,69)};
          \addlegendentry{Kit \nomproduit{}}
          \node[anchor=west,font=\scriptsize\bfseries,text=insagris] at (axis cs:36.6,644) {644\,€};
          \node[anchor=west,font=\scriptsize\bfseries,text=insaorange] at (axis cs:36.6,827) {827\,€};
          \node[anchor=west,font=\scriptsize\bfseries,text=insarouge] at (axis cs:36.6,69) {69\,€};
        \end{axis}
      \end{tikzpicture}
    \end{column}
    \begin{column}{0.35\textwidth}
      \begin{remarque}[Un complément]
        \begin{puces}
          \item Il ne remplace pas ton téléphone~: appels, internet.
          \item Aucun centre de secours n'est relié~: ce n'est pas une balise de détresse.
          \item Tes contacts doivent être équipés aussi.
        \end{puces}
      \end{remarque}
    \end{column}
  \end{columns}
  \vspace{0.1cm}
  \begin{pointcle}\centering\footnotesize
    69\,€, c'est moins de quatre mois d'un forfait mobile moyen (17,88\,€ par mois).
  \end{pointcle}
  \srcnote{Hypothèses~: 36 mois, 1\,\$ = 0,86\,€, hors remplacement de batterie. Balise~: inReach Mini 2 (328,90\,€), forfait 14,99\,\$ par mois, activation 39,99\,\$.
           Sources~: \cite{arcep_prix,garmin_mini2,garmin_plans}}
\end{frame}

%---------------------------------------------------------------------------
% Diapositive 12 : quatre terrains d'usage
%---------------------------------------------------------------------------
\begin{frame}{Là où le réseau s'arrête~: quatre terrains}
  \begin{columns}[T]
    \begin{column}{0.245\textwidth}
      \begin{carte}[height=4.0cm,valign=top]
        \centering
        \begin{tikzpicture}\path (-0.7,-0.6) rectangle (0.7,0.6);
          \pic[insagris,scale=1.05] at (-0.12,-0.05) {icone tour};\pic[scale=0.95] at (0.42,0.32) {icone croix};\end{tikzpicture}\\[3pt]
        \gros{Campagne}\\[3pt]
        {\footnotesize Il reste des zones blanches malgré 99\,\% de couverture 4G.}
      \end{carte}
    \end{column}
    \begin{column}{0.245\textwidth}
      \begin{carte}[height=4.0cm,valign=top]
        \centering
        \begin{tikzpicture}\path (-0.7,-0.6) rectangle (0.7,0.6);\pic[scale=1.5] at (0,0) {icone mont};\end{tikzpicture}\\[3pt]
        \gros{Montagne}\\[3pt]
        {\footnotesize 27 millions de Français pratiquent la randonnée~; le relief coupe le signal.}
      \end{carte}
    \end{column}
    \begin{column}{0.245\textwidth}
      \begin{carte}[height=4.0cm,valign=top]
        \centering
        \begin{tikzpicture}\path (-0.7,-0.6) rectangle (0.7,0.6);\pic[insagris,scale=1.5] at (0,0) {icone immeuble};\end{tikzpicture}\\[3pt]
        \gros{Guerre}\\[3pt]
        {\footnotesize Antennes détruites, coupures imposées~: un maillage n'a pas de centre.}
      \end{carte}
    \end{column}
    \begin{column}{0.245\textwidth}
      \begin{carte}[height=4.0cm,valign=top]
        \centering
        \begin{tikzpicture}\path (-0.7,-0.6) rectangle (0.7,0.6);\pic[insarouge,scale=1.8] at (0,0) {icone vague};\end{tikzpicture}\\[3pt]
        \gros{Catastrophes}\\[3pt]
        {\footnotesize Tempête, séisme, inondation~: les relais tombent, les kits sur batterie restent.}
      \end{carte}
    \end{column}
  \end{columns}
  \vspace{0.15cm}
  \begin{pointcle}\centering\footnotesize
    Un seul produit pour quatre terrains~: la même antenne, le même logiciel.
  \end{pointcle}
  \srcnote{Sources~: \cite{arcep_mobile,ffrando,laresilience}}
\end{frame}
         % 5.  Marketing : promesses, terrains, ce qu'on ne promet pas
%%%%%%%%%%%%%%%%%%%%%%%%%%%%%%%%%%%%%%%%%%%%%%%%%%%%%%%%%%%%%%%%%%%%%%%%%%%%%%
%  Partie « Finances » (2 diapositives) -- à importer avec %%%%%%%%%%%%%%%%%%%%%%%%%%%%%%%%%%%%%%%%%%%%%%%%%%%%%%%%%%%%%%%%%%%%%%%%%%%%%%
%  Partie « Finances » (2 diapositives) -- à importer avec %%%%%%%%%%%%%%%%%%%%%%%%%%%%%%%%%%%%%%%%%%%%%%%%%%%%%%%%%%%%%%%%%%%%%%%%%%%%%%
%  Partie « Finances » (2 diapositives) -- à importer avec \include{chap_finances}
%  Fichier de contenu : ni préambule ni \begin{document}, commence par \section.
%  Décomposition du prix d'un kit, compte de résultat sur trois ans.
%  Source unique du modèle : les parties SWOT (marges par offre, sensibilités) et Étude quantitative
%  (seuils) reprennent ces hypothèses sans les redéfinir.
%  Tous les chiffres sont des HYPOTHÈSES de travail (scénario de référence), cohérentes entre
%  elles : prix 69 € TTC (57,50 € HT), matériel 30,30 €, vente 4 €, kits solidaires 1 pour 10,
%  volumes 2 000 / 8 000 / 22 000 kits, charges fixes 200 / 330 / 430 k€.
%%%%%%%%%%%%%%%%%%%%%%%%%%%%%%%%%%%%%%%%%%%%%%%%%%%%%%%%%%%%%%%%%%%%%%%%%%%%%%
\section{Finances}

%---------------------------------------------------------------------------
% où vont les 69 € d'un kit
%---------------------------------------------------------------------------
\begin{frame}{Où vont les 69\,€ d'un kit (année 3)~?}
  \centering
  % échelle : 1 unité = 1 € = 0,1884 cm, donc 69 € = 13 cm
  \begin{tikzpicture}[font=\scriptsize,x=0.18841cm,y=1cm]
    \def\h{1.25}
    \fill[insagris!25] (0,0)       rectangle (11.5,\h);
    \fill[insarouge]   (11.5,0)    rectangle (41.8,\h);
    \fill[insasafran]  (41.8,0)    rectangle (48.83,\h);
    \fill[insaorange]  (48.83,0)   rectangle (59.74,\h);
    \fill[insarose]    (59.74,0)   rectangle (68.38,\h);
    \fill[black]       (68.38,0)   rectangle (69,\h);
    \draw[white,line width=0.9pt] (11.5,0) -- (11.5,\h) (41.8,0) -- (41.8,\h) (48.83,0) -- (48.83,\h)
                                  (59.74,0) -- (59.74,\h) (68.38,0) -- (68.38,\h);
    \node[align=center,text=black] at (5.75,0.625)  {TVA\\{\bfseries 11,50\,€}};
    \node[align=center,text=white] at (26.65,0.625) {{\small\bfseries L'antenne}~: matériel et assemblage\\{\large\bfseries\sffamily 30,30\,€}};
    \node[align=center,text=black] at (45.3,0.625)  {Vente,\\kits\\offerts\\{\bfseries 7,03\,€}};
    \node[align=center,text=white] at (54.3,0.625)  {{\bfseries Développeurs}\\{\bfseries 10,91\,€}};
    \node[align=center,text=black] at (64.06,0.625) {Autres\\charges\\{\bfseries 8,64\,€}};
    \node[anchor=south east,font=\tiny,text=black,align=right] at (69,\h+0.04) {résultat 0,62\,€};
    \draw[decorate,decoration={brace,mirror,amplitude=4pt},insagris,thick] (48.83,-0.12) -- (68.38,-0.12)
      node[midway,below=6pt,font=\scriptsize,text=insagris] {charges fixes~: 19,55\,€ par kit};
  \end{tikzpicture}

  \vspace{-0.3cm}
  \begin{columns}[T]
    \begin{column}{0.56\textwidth}
      \begin{encadre}{Dans les 30,30\,€ de l'antenne}
        \footnotesize\renewcommand{\arraystretch}{0.95}
        \begin{tabularx}{\linewidth}{@{}Xr@{}}
          Carte Heltec V3 + antenne LoRa   & 17,00\,€ \\
          Batterie Li-Po                   & 4,50\,€ \\
          Assemblage, flashage, test (10 min) & 3,50\,€ \\
          Boîtier imprimé en 3D            & 2,50\,€ \\
          Emballage, carte de démarrage    & 1,80\,€ \\
          Câble USB-C                      & 1,00\,€ \\
        \end{tabularx}
      \end{encadre}
    \end{column}
    \begin{column}{0.41\textwidth}
      \begin{carte}[height=1.58cm,valign=top]
        \raggedright
        \gros{≈ 11\,€ par kit}\\[1pt]
        {\footnotesize financent quatre développeurs du logiciel libre.}
      \end{carte}
      \vspace{0.1cm}
      \begin{carte}[height=1.58cm,valign=top]
        \raggedright
        \gros{17,90 à 19,90\,\$}\\[1pt]
        {\footnotesize prix public d'une carte V3~: nous comptons 17\,€.}
      \end{carte}
    \end{column}
  \end{columns}
  \srcnote{Hypothèses de travail~: 22\,000 kits (année 3), TVA 20\,\%, 1 kit offert pour 10 vendus. Source~: \cite{heltec_v3}}
\end{frame}

%---------------------------------------------------------------------------
% trois ans de compte de résultat
%---------------------------------------------------------------------------
\begin{frame}{Compte de résultat sur trois ans}
  \begin{columns}[T]
    \begin{column}{0.57\textwidth}
      \renewcommand{\arraystretch}{1.17}
      \footnotesize
      \begin{tabular}{@{}l@{\hspace{4pt}}r@{\hspace{9pt}}r@{\hspace{9pt}}r@{}}
        \rowcolor{insarose}
        \textcolor{insarouge}{\bfseries k€ sauf mention} & \textcolor{insarouge}{\bfseries An 1} & \textcolor{insarouge}{\bfseries An 2} & \textcolor{insarouge}{\bfseries An 3} \\
        Kits vendus                 & 2\,000    & 8\,000    & 22\,000 \\
        Chiffre d'affaires HT       & 115,0     & 460,0     & 1\,265,0 \\
        Coûts variables             & −74,7     & −298,6    & −821,3 \\
        Marge de contribution       & 40,3      & 161,4     & 443,7 \\
        Charges fixes               & −200,0    & −330,0    & −430,0 \\
        \quad dont développeurs     & 120,0     & 180,0     & 240,0 \\
        \rowcolor{insabeige}
        \bfseries Résultat          & \bfseries −159,7 & \bfseries −168,6 & \bfseries +13,7 \\
        Résultat cumulé             & −159,7    & −328,3    & −314,6 \\
      \end{tabular}
    \end{column}
    \begin{column}{0.40\textwidth}
      \centering
      \begin{tikzpicture}
        \begin{axis}[width=5.3cm,height=3.8cm,xmin=0,xmax=3,ymin=-520,ymax=60,
            xtick={0,1,2,3},xticklabels={départ,An 1,An 2,An 3},ytick={-400,-200,0},
            tick label style={font=\scriptsize},
            title={Résultat cumulé (k€)},title style={font=\scriptsize\bfseries,yshift=-3pt},
            axis x line=bottom,axis y line=left,grid=major,grid style={insagris!25,thin},clip=false]
          \addplot[insarouge,very thick,mark=*,mark size=2pt] coordinates {(0,0) (1,-159.7) (2,-328.3) (3,-314.6)};
          \node[font=\tiny,anchor=north] at (axis cs:1,-178) {−160};
          \node[font=\tiny,anchor=north] at (axis cs:2,-346) {−328};
          \node[font=\tiny,anchor=north] at (axis cs:3,-332) {−315};
          \draw[dashed,insaorange,thick] (axis cs:0,-450) -- (axis cs:3,-450);
          \node[font=\tiny,text=insaorange,anchor=south west] at (axis cs:0.05,-445) {besoin ≈ 450};
        \end{axis}
      \end{tikzpicture}
    \end{column}
  \end{columns}
  \vspace{0.15cm}
  \begin{columns}[T]
    \begin{column}{0.325\textwidth}
      \begin{carte}[height=1.75cm,valign=top]
        \raggedright\gros{≈ 21\,300 kits}\\[1pt]
        {\footnotesize par an~: le point mort, avec les charges fixes de l'année 3.}
      \end{carte}
    \end{column}
    \begin{column}{0.325\textwidth}
      \begin{carte}[height=1.75cm,valign=top]
        \raggedright\gros{≈ 450\,k€}\\[1pt]
        {\footnotesize à financer~: 328\,k€ de pertes, 100\,k€ de stock, une marge.}
      \end{carte}
    \end{column}
    \begin{column}{0.325\textwidth}
      \begin{carte}[height=1.75cm,valign=top]
        \raggedright\gros{−127\,k€}\\[1pt]
        {\footnotesize résultat de l'année 3 si seuls 15\,000 kits sont vendus.}
      \end{carte}
    \end{column}
  \end{columns}
  \srcnote{Scénario de référence, hypothèses à valider~: aucun salaire fondateur, aucune aide publique. Part de marché visée~: voir la partie Marché cible.}
\end{frame}

%  Fichier de contenu : ni préambule ni \begin{document}, commence par \section.
%  Décomposition du prix d'un kit, compte de résultat sur trois ans.
%  Source unique du modèle : les parties SWOT (marges par offre, sensibilités) et Étude quantitative
%  (seuils) reprennent ces hypothèses sans les redéfinir.
%  Tous les chiffres sont des HYPOTHÈSES de travail (scénario de référence), cohérentes entre
%  elles : prix 69 € TTC (57,50 € HT), matériel 30,30 €, vente 4 €, kits solidaires 1 pour 10,
%  volumes 2 000 / 8 000 / 22 000 kits, charges fixes 200 / 330 / 430 k€.
%%%%%%%%%%%%%%%%%%%%%%%%%%%%%%%%%%%%%%%%%%%%%%%%%%%%%%%%%%%%%%%%%%%%%%%%%%%%%%
\section{Finances}

%---------------------------------------------------------------------------
% où vont les 69 € d'un kit
%---------------------------------------------------------------------------
\begin{frame}{Où vont les 69\,€ d'un kit (année 3)~?}
  \centering
  % échelle : 1 unité = 1 € = 0,1884 cm, donc 69 € = 13 cm
  \begin{tikzpicture}[font=\scriptsize,x=0.18841cm,y=1cm]
    \def\h{1.25}
    \fill[insagris!25] (0,0)       rectangle (11.5,\h);
    \fill[insarouge]   (11.5,0)    rectangle (41.8,\h);
    \fill[insasafran]  (41.8,0)    rectangle (48.83,\h);
    \fill[insaorange]  (48.83,0)   rectangle (59.74,\h);
    \fill[insarose]    (59.74,0)   rectangle (68.38,\h);
    \fill[black]       (68.38,0)   rectangle (69,\h);
    \draw[white,line width=0.9pt] (11.5,0) -- (11.5,\h) (41.8,0) -- (41.8,\h) (48.83,0) -- (48.83,\h)
                                  (59.74,0) -- (59.74,\h) (68.38,0) -- (68.38,\h);
    \node[align=center,text=black] at (5.75,0.625)  {TVA\\{\bfseries 11,50\,€}};
    \node[align=center,text=white] at (26.65,0.625) {{\small\bfseries L'antenne}~: matériel et assemblage\\{\large\bfseries\sffamily 30,30\,€}};
    \node[align=center,text=black] at (45.3,0.625)  {Vente,\\kits\\offerts\\{\bfseries 7,03\,€}};
    \node[align=center,text=white] at (54.3,0.625)  {{\bfseries Développeurs}\\{\bfseries 10,91\,€}};
    \node[align=center,text=black] at (64.06,0.625) {Autres\\charges\\{\bfseries 8,64\,€}};
    \node[anchor=south east,font=\tiny,text=black,align=right] at (69,\h+0.04) {résultat 0,62\,€};
    \draw[decorate,decoration={brace,mirror,amplitude=4pt},insagris,thick] (48.83,-0.12) -- (68.38,-0.12)
      node[midway,below=6pt,font=\scriptsize,text=insagris] {charges fixes~: 19,55\,€ par kit};
  \end{tikzpicture}

  \vspace{-0.3cm}
  \begin{columns}[T]
    \begin{column}{0.56\textwidth}
      \begin{encadre}{Dans les 30,30\,€ de l'antenne}
        \footnotesize\renewcommand{\arraystretch}{0.95}
        \begin{tabularx}{\linewidth}{@{}Xr@{}}
          Carte Heltec V3 + antenne LoRa   & 17,00\,€ \\
          Batterie Li-Po                   & 4,50\,€ \\
          Assemblage, flashage, test (10 min) & 3,50\,€ \\
          Boîtier imprimé en 3D            & 2,50\,€ \\
          Emballage, carte de démarrage    & 1,80\,€ \\
          Câble USB-C                      & 1,00\,€ \\
        \end{tabularx}
      \end{encadre}
    \end{column}
    \begin{column}{0.41\textwidth}
      \begin{carte}[height=1.58cm,valign=top]
        \raggedright
        \gros{≈ 11\,€ par kit}\\[1pt]
        {\footnotesize financent quatre développeurs du logiciel libre.}
      \end{carte}
      \vspace{0.1cm}
      \begin{carte}[height=1.58cm,valign=top]
        \raggedright
        \gros{17,90 à 19,90\,\$}\\[1pt]
        {\footnotesize prix public d'une carte V3~: nous comptons 17\,€.}
      \end{carte}
    \end{column}
  \end{columns}
  \srcnote{Hypothèses de travail~: 22\,000 kits (année 3), TVA 20\,\%, 1 kit offert pour 10 vendus. Source~: \cite{heltec_v3}}
\end{frame}

%---------------------------------------------------------------------------
% trois ans de compte de résultat
%---------------------------------------------------------------------------
\begin{frame}{Compte de résultat sur trois ans}
  \begin{columns}[T]
    \begin{column}{0.57\textwidth}
      \renewcommand{\arraystretch}{1.17}
      \footnotesize
      \begin{tabular}{@{}l@{\hspace{4pt}}r@{\hspace{9pt}}r@{\hspace{9pt}}r@{}}
        \rowcolor{insarose}
        \textcolor{insarouge}{\bfseries k€ sauf mention} & \textcolor{insarouge}{\bfseries An 1} & \textcolor{insarouge}{\bfseries An 2} & \textcolor{insarouge}{\bfseries An 3} \\
        Kits vendus                 & 2\,000    & 8\,000    & 22\,000 \\
        Chiffre d'affaires HT       & 115,0     & 460,0     & 1\,265,0 \\
        Coûts variables             & −74,7     & −298,6    & −821,3 \\
        Marge de contribution       & 40,3      & 161,4     & 443,7 \\
        Charges fixes               & −200,0    & −330,0    & −430,0 \\
        \quad dont développeurs     & 120,0     & 180,0     & 240,0 \\
        \rowcolor{insabeige}
        \bfseries Résultat          & \bfseries −159,7 & \bfseries −168,6 & \bfseries +13,7 \\
        Résultat cumulé             & −159,7    & −328,3    & −314,6 \\
      \end{tabular}
    \end{column}
    \begin{column}{0.40\textwidth}
      \centering
      \begin{tikzpicture}
        \begin{axis}[width=5.3cm,height=3.8cm,xmin=0,xmax=3,ymin=-520,ymax=60,
            xtick={0,1,2,3},xticklabels={départ,An 1,An 2,An 3},ytick={-400,-200,0},
            tick label style={font=\scriptsize},
            title={Résultat cumulé (k€)},title style={font=\scriptsize\bfseries,yshift=-3pt},
            axis x line=bottom,axis y line=left,grid=major,grid style={insagris!25,thin},clip=false]
          \addplot[insarouge,very thick,mark=*,mark size=2pt] coordinates {(0,0) (1,-159.7) (2,-328.3) (3,-314.6)};
          \node[font=\tiny,anchor=north] at (axis cs:1,-178) {−160};
          \node[font=\tiny,anchor=north] at (axis cs:2,-346) {−328};
          \node[font=\tiny,anchor=north] at (axis cs:3,-332) {−315};
          \draw[dashed,insaorange,thick] (axis cs:0,-450) -- (axis cs:3,-450);
          \node[font=\tiny,text=insaorange,anchor=south west] at (axis cs:0.05,-445) {besoin ≈ 450};
        \end{axis}
      \end{tikzpicture}
    \end{column}
  \end{columns}
  \vspace{0.15cm}
  \begin{columns}[T]
    \begin{column}{0.325\textwidth}
      \begin{carte}[height=1.75cm,valign=top]
        \raggedright\gros{≈ 21\,300 kits}\\[1pt]
        {\footnotesize par an~: le point mort, avec les charges fixes de l'année 3.}
      \end{carte}
    \end{column}
    \begin{column}{0.325\textwidth}
      \begin{carte}[height=1.75cm,valign=top]
        \raggedright\gros{≈ 450\,k€}\\[1pt]
        {\footnotesize à financer~: 328\,k€ de pertes, 100\,k€ de stock, une marge.}
      \end{carte}
    \end{column}
    \begin{column}{0.325\textwidth}
      \begin{carte}[height=1.75cm,valign=top]
        \raggedright\gros{−127\,k€}\\[1pt]
        {\footnotesize résultat de l'année 3 si seuls 15\,000 kits sont vendus.}
      \end{carte}
    \end{column}
  \end{columns}
  \srcnote{Scénario de référence, hypothèses à valider~: aucun salaire fondateur, aucune aide publique. Part de marché visée~: voir la partie Marché cible.}
\end{frame}

%  Fichier de contenu : ni préambule ni \begin{document}, commence par \section.
%  Décomposition du prix d'un kit, compte de résultat sur trois ans.
%  Source unique du modèle : les parties SWOT (marges par offre, sensibilités) et Étude quantitative
%  (seuils) reprennent ces hypothèses sans les redéfinir.
%  Tous les chiffres sont des HYPOTHÈSES de travail (scénario de référence), cohérentes entre
%  elles : prix 69 € TTC (57,50 € HT), matériel 30,30 €, vente 4 €, kits solidaires 1 pour 10,
%  volumes 2 000 / 8 000 / 22 000 kits, charges fixes 200 / 330 / 430 k€.
%%%%%%%%%%%%%%%%%%%%%%%%%%%%%%%%%%%%%%%%%%%%%%%%%%%%%%%%%%%%%%%%%%%%%%%%%%%%%%
\section{Finances}

%---------------------------------------------------------------------------
% où vont les 69 € d'un kit
%---------------------------------------------------------------------------
\begin{frame}{Où vont les 69\,€ d'un kit (année 3)~?}
  \centering
  % échelle : 1 unité = 1 € = 0,1884 cm, donc 69 € = 13 cm
  \begin{tikzpicture}[font=\scriptsize,x=0.18841cm,y=1cm]
    \def\h{1.25}
    \fill[insagris!25] (0,0)       rectangle (11.5,\h);
    \fill[insarouge]   (11.5,0)    rectangle (41.8,\h);
    \fill[insasafran]  (41.8,0)    rectangle (48.83,\h);
    \fill[insaorange]  (48.83,0)   rectangle (59.74,\h);
    \fill[insarose]    (59.74,0)   rectangle (68.38,\h);
    \fill[black]       (68.38,0)   rectangle (69,\h);
    \draw[white,line width=0.9pt] (11.5,0) -- (11.5,\h) (41.8,0) -- (41.8,\h) (48.83,0) -- (48.83,\h)
                                  (59.74,0) -- (59.74,\h) (68.38,0) -- (68.38,\h);
    \node[align=center,text=black] at (5.75,0.625)  {TVA\\{\bfseries 11,50\,€}};
    \node[align=center,text=white] at (26.65,0.625) {{\small\bfseries L'antenne}~: matériel et assemblage\\{\large\bfseries\sffamily 30,30\,€}};
    \node[align=center,text=black] at (45.3,0.625)  {Vente,\\kits\\offerts\\{\bfseries 7,03\,€}};
    \node[align=center,text=white] at (54.3,0.625)  {{\bfseries Développeurs}\\{\bfseries 10,91\,€}};
    \node[align=center,text=black] at (64.06,0.625) {Autres\\charges\\{\bfseries 8,64\,€}};
    \node[anchor=south east,font=\tiny,text=black,align=right] at (69,\h+0.04) {résultat 0,62\,€};
    \draw[decorate,decoration={brace,mirror,amplitude=4pt},insagris,thick] (48.83,-0.12) -- (68.38,-0.12)
      node[midway,below=6pt,font=\scriptsize,text=insagris] {charges fixes~: 19,55\,€ par kit};
  \end{tikzpicture}

  \vspace{-0.3cm}
  \begin{columns}[T]
    \begin{column}{0.56\textwidth}
      \begin{encadre}{Dans les 30,30\,€ de l'antenne}
        \footnotesize\renewcommand{\arraystretch}{0.95}
        \begin{tabularx}{\linewidth}{@{}Xr@{}}
          Carte Heltec V3 + antenne LoRa   & 17,00\,€ \\
          Batterie Li-Po                   & 4,50\,€ \\
          Assemblage, flashage, test (10 min) & 3,50\,€ \\
          Boîtier imprimé en 3D            & 2,50\,€ \\
          Emballage, carte de démarrage    & 1,80\,€ \\
          Câble USB-C                      & 1,00\,€ \\
        \end{tabularx}
      \end{encadre}
    \end{column}
    \begin{column}{0.41\textwidth}
      \begin{carte}[height=1.58cm,valign=top]
        \raggedright
        \gros{≈ 11\,€ par kit}\\[1pt]
        {\footnotesize financent quatre développeurs du logiciel libre.}
      \end{carte}
      \vspace{0.1cm}
      \begin{carte}[height=1.58cm,valign=top]
        \raggedright
        \gros{17,90 à 19,90\,\$}\\[1pt]
        {\footnotesize prix public d'une carte V3~: nous comptons 17\,€.}
      \end{carte}
    \end{column}
  \end{columns}
  \srcnote{Hypothèses de travail~: 22\,000 kits (année 3), TVA 20\,\%, 1 kit offert pour 10 vendus. Source~: \cite{heltec_v3}}
\end{frame}

%---------------------------------------------------------------------------
% trois ans de compte de résultat
%---------------------------------------------------------------------------
\begin{frame}{Compte de résultat sur trois ans}
  \begin{columns}[T]
    \begin{column}{0.57\textwidth}
      \renewcommand{\arraystretch}{1.17}
      \footnotesize
      \begin{tabular}{@{}l@{\hspace{4pt}}r@{\hspace{9pt}}r@{\hspace{9pt}}r@{}}
        \rowcolor{insarose}
        \textcolor{insarouge}{\bfseries k€ sauf mention} & \textcolor{insarouge}{\bfseries An 1} & \textcolor{insarouge}{\bfseries An 2} & \textcolor{insarouge}{\bfseries An 3} \\
        Kits vendus                 & 2\,000    & 8\,000    & 22\,000 \\
        Chiffre d'affaires HT       & 115,0     & 460,0     & 1\,265,0 \\
        Coûts variables             & −74,7     & −298,6    & −821,3 \\
        Marge de contribution       & 40,3      & 161,4     & 443,7 \\
        Charges fixes               & −200,0    & −330,0    & −430,0 \\
        \quad dont développeurs     & 120,0     & 180,0     & 240,0 \\
        \rowcolor{insabeige}
        \bfseries Résultat          & \bfseries −159,7 & \bfseries −168,6 & \bfseries +13,7 \\
        Résultat cumulé             & −159,7    & −328,3    & −314,6 \\
      \end{tabular}
    \end{column}
    \begin{column}{0.40\textwidth}
      \centering
      \begin{tikzpicture}
        \begin{axis}[width=5.3cm,height=3.8cm,xmin=0,xmax=3,ymin=-520,ymax=60,
            xtick={0,1,2,3},xticklabels={départ,An 1,An 2,An 3},ytick={-400,-200,0},
            tick label style={font=\scriptsize},
            title={Résultat cumulé (k€)},title style={font=\scriptsize\bfseries,yshift=-3pt},
            axis x line=bottom,axis y line=left,grid=major,grid style={insagris!25,thin},clip=false]
          \addplot[insarouge,very thick,mark=*,mark size=2pt] coordinates {(0,0) (1,-159.7) (2,-328.3) (3,-314.6)};
          \node[font=\tiny,anchor=north] at (axis cs:1,-178) {−160};
          \node[font=\tiny,anchor=north] at (axis cs:2,-346) {−328};
          \node[font=\tiny,anchor=north] at (axis cs:3,-332) {−315};
          \draw[dashed,insaorange,thick] (axis cs:0,-450) -- (axis cs:3,-450);
          \node[font=\tiny,text=insaorange,anchor=south west] at (axis cs:0.05,-445) {besoin ≈ 450};
        \end{axis}
      \end{tikzpicture}
    \end{column}
  \end{columns}
  \vspace{0.15cm}
  \begin{columns}[T]
    \begin{column}{0.325\textwidth}
      \begin{carte}[height=1.75cm,valign=top]
        \raggedright\gros{≈ 21\,300 kits}\\[1pt]
        {\footnotesize par an~: le point mort, avec les charges fixes de l'année 3.}
      \end{carte}
    \end{column}
    \begin{column}{0.325\textwidth}
      \begin{carte}[height=1.75cm,valign=top]
        \raggedright\gros{≈ 450\,k€}\\[1pt]
        {\footnotesize à financer~: 328\,k€ de pertes, 100\,k€ de stock, une marge.}
      \end{carte}
    \end{column}
    \begin{column}{0.325\textwidth}
      \begin{carte}[height=1.75cm,valign=top]
        \raggedright\gros{−127\,k€}\\[1pt]
        {\footnotesize résultat de l'année 3 si seuls 15\,000 kits sont vendus.}
      \end{carte}
    \end{column}
  \end{columns}
  \srcnote{Scénario de référence, hypothèses à valider~: aucun salaire fondateur, aucune aide publique. Part de marché visée~: voir la partie Marché cible.}
\end{frame}
          % 6.  Finances
%%%%%%%%%%%%%%%%%%%%%%%%%%%%%%%%%%%%%%%%%%%%%%%%%%%%%%%%%%%%%%%%%%%%%%%%%%%%%%
%  Partie 6 -- Produire en France : la technologie est simple, la relocalisation possible par étapes
%  (2 diapositives : 17 = matrice des trois phases ; 18 = avec qui produire : injection,
%  puces de Grenoble, acier français)
%
%  Hypothèses chiffrées (à confirmer par devis) :
%  - carte française = 22 € (composants 14,5 + circuit 1,5 + assemblage 6), soit un coût de kit de
%    35,30 € (phase 2) puis 35,80 € (phase 3) ;
%  - phase 3 : boîtier injecté = pièce 1,50 € + moule amorti 1,50 € (≈ 33 k€ pour 22 000 kits,
%    moule en deux pièces : coque et capot) = 3,00 €, contre 2,50 € pour le boîtier imprimé en 3D ;
%  - les prix de 75 et 79 € ne contiennent PAS le surcoût éventuel de puces ou d'antennes françaises
%    (non chiffré) ;
%  - un fouet en acier de 3 mm sur 8 cm pèse ≈ 5 g : 22 000 kits ≈ 110 kg d'acier par an.
%%%%%%%%%%%%%%%%%%%%%%%%%%%%%%%%%%%%%%%%%%%%%%%%%%%%%%%%%%%%%%%%%%%%%%%%%%%%%%
\section{Produire en France}

% cellule « fabriqué en France » (coche rouge) et cellule « importé » (rond gris) : \cellok{x}{y}{texte}
\newcommand{\cellok}[3]{\pic[insarouge,scale=1.0] at ({#1-1.0},#2) {icone coche};
                        \node[anchor=west,font=\scriptsize] at ({#1-0.76},#2) {#3};}
\newcommand{\cellimp}[3]{\draw[insagris,thick] ({#1-1.0},#2) circle (0.13);
                         \node[anchor=west,font=\scriptsize,text=insagris] at ({#1-0.76},#2) {#3};}

%---------------------------------------------------------------------------
% Diapositive 17 : trois phases de relocalisation
%---------------------------------------------------------------------------
\begin{frame}{Technologie simple~: produisons en France}
  \centering
  \begin{tikzpicture}[font=\footnotesize]
    % en-têtes : trois phases, avec le prix public visé
    \node[pastille,draw=insarouge,fill=insarose,text=black,minimum width=2.8cm,minimum height=0.9cm] at (5.2,0)
      {An 1\\{\scriptsize\mdseries assemblé en France}\\{\scriptsize\color{insarouge}69\,€}};
    \node[pastille,draw=insarouge,fill=insarose,text=black,minimum width=2.8cm,minimum height=0.9cm] at (8.2,0)
      {Ans 2 et 3\\{\scriptsize\mdseries carte conçue ici}\\{\scriptsize\color{insarouge}75\,€}};
    \node[pastille,draw=insarouge,fill=insarose,text=black,minimum width=2.8cm,minimum height=0.9cm] at (11.2,0)
      {An 4 et après\\{\scriptsize\mdseries quasi tout ici}\\{\scriptsize\color{insarouge}79\,€}};
    % bandes de lecture (une ligne sur deux)
    \fill[insabeige] (0,-1.25) rectangle (12.7,-0.85);
    \fill[insabeige] (0,-2.05) rectangle (12.7,-1.65);
    \fill[insabeige] (0,-2.85) rectangle (12.7,-2.45);
    % libellés des lignes
    \node[anchor=west,font=\footnotesize\bfseries] at (0.1,-1.05) {Puce radio LoRa};
    \node[anchor=west,font=\footnotesize\bfseries] at (0.1,-1.45) {Processeur};
    \node[anchor=west,font=\footnotesize\bfseries] at (0.1,-1.85) {Carte électronique};
    \node[anchor=west,font=\footnotesize\bfseries] at (0.1,-2.25) {Boîtier};
    \node[anchor=west,font=\footnotesize\bfseries] at (0.1,-2.65) {Antenne};
    \node[anchor=west,font=\footnotesize\bfseries] at (0.1,-3.05) {Assemblage, flash, test};
    % puce radio : importée, puis fabriquée à Grenoble (objectif)
    \cellimp{5.2}{-1.05}{importée}\cellimp{8.2}{-1.05}{importée}\cellok{11.2}{-1.05}{Grenoble (visé)}
    % processeur : importé dans les trois phases
    \cellimp{5.2}{-1.45}{importé}\cellimp{8.2}{-1.45}{importé}\cellimp{11.2}{-1.45}{importé}
    % carte
    \cellimp{5.2}{-1.85}{module Heltec}\cellok{8.2}{-1.85}{conçue ici}\cellok{11.2}{-1.85}{conçue ici}
    % boîtier
    \cellok{5.2}{-2.25}{impression 3D}\cellok{8.2}{-2.25}{impression 3D}\cellok{11.2}{-2.25}{injection}
    % antenne
    \cellimp{5.2}{-2.65}{d'origine Heltec}\cellok{8.2}{-2.65}{acier, à valider}\cellok{11.2}{-2.65}{acier français}
    % assemblage
    \cellok{5.2}{-3.05}{France}\cellok{8.2}{-3.05}{France}\cellok{11.2}{-3.05}{France}
  \end{tikzpicture}

  \vspace{-0.45cm}
  \begin{columns}[T]
    \begin{column}{0.325\textwidth}
      \begin{carte}[height=1.05cm,valign=top]
        \raggedright\gros{+ 5 à 6\,€ par kit}\\[1pt]
        {\footnotesize 35,30\,€ au lieu de 30,30\,€.}
      \end{carte}
    \end{column}
    \begin{column}{0.325\textwidth}
      \begin{carte}[height=1.05cm,valign=top]
        \raggedright\gros{≈ 2 emplois}\\[1pt]
        {\footnotesize d'assemblage à 22\,000 kits.}
      \end{carte}
    \end{column}
    \begin{column}{0.325\textwidth}
      \begin{carte}[height=1.05cm,valign=top]
        \raggedright\gros{Né à Grenoble}\\[1pt]
        {\footnotesize LoRa~: inventé par Cycléo.}
      \end{carte}
    \end{column}
  \end{columns}
  \vspace{0cm}
  \begin{pointcle}\centering\footnotesize
    Carte de 5\,cm sur 2,5\,cm, composants du commerce~: rien d'irréalisable chez nous.
  \end{pointcle}
  \srcnote{Estimations à confirmer par devis~; prix de 75 et 79\,€ calculés pour garder une marge de 20 à 22\,€ par kit, hors surcoût éventuel de puces et d'antennes françaises. Sources~: \cite{actutem_lora,heltec_v3}}
\end{frame}

%---------------------------------------------------------------------------
% Diapositive 18 : avec qui produire ? injection plastique, puces de Grenoble, acier français
%---------------------------------------------------------------------------
\begin{frame}{Boîtier, puces, antennes~: avec qui produire~?}
  \begin{columns}[T]
    \begin{column}{0.6\textwidth}
      % schéma d'une presse à injecter : trémie, vis chauffante, moule en acier, pièce
      \begin{tikzpicture}[font=\footnotesize,scale=0.85]
        % trémie remplie de granulés
        \draw[insagris,thick,fill=insagris!10] (0.1,1.1) -- (1.3,1.1) -- (0.95,0.5) -- (0.45,0.5) -- cycle;
        \foreach \x/\y in {0.4/0.95,0.65/1.0,0.9/0.93,0.55/0.74,0.8/0.7} {\fill[insaorange] (\x,\y) circle (0.045);}
        \node[anchor=west,inner sep=1pt] at (1.45,0.92) {granulés};
        % fourreau, vis et colliers chauffants (rouges)
        \draw[insagris,thick,fill=insagris!15,rounded corners=2pt] (0,0) rectangle (3.3,0.5);
        \foreach \i in {0,...,9} {\draw[insagris,thin] ({0.2+0.3*\i},0.07) -- ({0.35+0.3*\i},0.43);}
        \foreach \x in {1.45,2.1,2.75} {\fill[insarouge] (\x,-0.05) rectangle (\x+0.12,0.55);}
        \fill[insagris!60] (-0.4,0.12) rectangle (0,0.38);
        \node[inner sep=1pt] at (1.65,-0.3) {vis chauffante};
        % buse et moule en acier : deux plaques et une cavité (matière en rouge)
        \draw[insagris,thick,fill=insagris!15] (3.3,0.13) -- (3.62,0.2) -- (3.62,0.3) -- (3.3,0.37) -- cycle;
        \draw[insagris,thick,fill=insagris!40] (3.62,-0.3) rectangle (5.4,0.25);
        \draw[insagris,thick,fill=insagris!40] (3.62,0.25) rectangle (5.4,0.8);
        \fill[insarouge!65] (3.62,0.22) rectangle (3.95,0.28);
        \draw[insarouge,fill=insarouge!65,rounded corners=1.5pt] (3.95,0.08) rectangle (5.05,0.42);
        \node[inner sep=1pt] at (4.5,1.0) {moule en acier};
        % pièce démoulée : le boîtier
        \draw[fluxr] (5.55,0.25) -- (6.25,0.25);
        \draw[insarouge,thick,fill=insarose,rounded corners=2pt] (6.4,-0.02) rectangle (7.6,0.52);
        \draw[insarouge,thick] (6.4,0.36) -- (7.6,0.36);
        \fill[insarouge] (7.38,0.18) circle (0.04);
        \node[inner sep=1pt] at (7.0,0.8) {boîtier};
      \end{tikzpicture}

      \vspace{0cm}
      {\footnotesize\raggedright
       \textbf{Injection~:} granulés fondus par la vis, injectés sous pression dans un moule en acier,
       refroidis, puis éjectés. Partenaires~: les plasturgistes d'Oyonnax (Plastics Vallée).\par}
    \end{column}
    \begin{column}{0.38\textwidth}
      \vspace{-0.3cm}
      \begin{carte}[height=1.05cm,valign=top]
        \raggedright\gros{5 à 50\,k€}\\[1pt]
        {\footnotesize pour un moule en acier.}
      \end{carte}
      \vspace{-0.15cm}
      \begin{carte}[height=1.05cm,valign=top]
        \raggedright\gros{0,50 à 5\,€}\\[1pt]
        {\footnotesize la pièce, hors moule.}
      \end{carte}
    \end{column}
  \end{columns}
  \vspace{0.02cm}
  \begin{columns}[T]
    \begin{column}{0.48\textwidth}
      \begin{encadre}{Puces~: l'écosystème de Grenoble}
        \begin{puces}
          \item Crolles grave en 22 puis 18\,nm (IoT, 5G)~: une puce radio, pas un processeur de pointe.
          \item ST, GlobalFoundries, Soitec. Piste~: STM32WL de ST, radio LoRa + processeur, soutenue par la communauté Meshtastic.
          \item Limite~: 22\,000 puces par an, trop peu pour une usine~; passer par un fabricant.
        \end{puces}
      \end{encadre}
    \end{column}
    \begin{column}{0.5\textwidth}
      \begin{encadre}{Antennes~: l'acier français}
        \begin{puces}
          \item ArcelorMittal~: 608 postes supprimés (Dunkerque, Florange, Basse-Indre…), mais 1,3\,Md€ en four électrique (2029).
          \item Basse-Indre~: acier plat étamé, soudable~: de quoi faire réflecteurs et plans de masse.
          \item ≈ 110\,kg d'acier par an, via des transformateurs. Cuivre, alu~: à comparer.
        \end{puces}
      \end{encadre}
    \end{column}
  \end{columns}
  \srcnote{Hypothèses~: moule ≈ 33\,k€ (1,50\,€ par kit) et pièce à 1,50\,€, soit 3,00\,€ contre 2,50\,€ imprimé~; fouet en acier de 5\,g × 22\,000 kits ≈ 110\,kg. Sources~: \cite{injection_cout,injection_procede,plastics_vallee,st_crolles,lmi_crolles,soitec,stm32wl,meshtastic_community,am_plan,am_four,am_basseindre,ttn_antennes}}
\end{frame}
            % 7.  Produire en France (réindustrialisation)
%%%%%%%%%%%%%%%%%%%%%%%%%%%%%%%%%%%%%%%%%%%%%%%%%%%%%%%%%%%%%%%%%%%%%%%%%%%%%%
%  Partie « Étude documentaire » (6 diapositives) -- à importer avec %%%%%%%%%%%%%%%%%%%%%%%%%%%%%%%%%%%%%%%%%%%%%%%%%%%%%%%%%%%%%%%%%%%%%%%%%%%%%%
%  Partie « Étude documentaire » (6 diapositives) -- à importer avec %%%%%%%%%%%%%%%%%%%%%%%%%%%%%%%%%%%%%%%%%%%%%%%%%%%%%%%%%%%%%%%%%%%%%%%%%%%%%%
%  Partie « Étude documentaire » (6 diapositives) -- à importer avec \include{etude}
%  Fichier de contenu : ni préambule ni \begin{document}, commence par \section.
%  Les sources sont citées dans chaque diapositive par des liens cliquables (nom de domaine),
%  sans passer par sources.bib. Pages consultées le 5 octobre 2026.
%  Niveaux : les chiffres du modèle financier (69 € TTC, 22 000 kits, etc.) sont des hypothèses.
%%%%%%%%%%%%%%%%%%%%%%%%%%%%%%%%%%%%%%%%%%%%%%%%%%%%%%%%%%%%%%%%%%%%%%%%%%%%%%
\section{Étude documentaire}
\providecommand{\etl}[2]{\href{#1}{#2}}

%---------------------------------------------------------------------------
\begin{frame}{Que mettre en avant, et auprès de qui~?}
  \begin{columns}[T]
    \begin{column}{0.37\textwidth}
      \begin{puces}
        \item \textbf{Payé une fois}~: 69\,€ contre 319\,€ puis 15 à 20\,€ par mois (Garmin inReach Mini 2).
        \item \textbf{Marche en cas de panne}~: la 4G couvre 99,2 à 99,6\,\% de la population, « zone blanche » seul est faible.
        \item \textbf{Prêt à l'emploi, légal, en français} (869,4--869,65\,MHz, 27\,dBm, 10\,\%).
        \item \textbf{Jamais « balise de détresse »}~: texte seul, 1 à 3\,km en ville.
      \end{puces}
    \end{column}
    \begin{column}{0.61\textwidth}
      \scriptsize\renewcommand{\arraystretch}{1.05}
      \begin{tabularx}{\linewidth}{@{}p{2.2cm}Xc@{}}
        \rowcolor{insarose}\textcolor{insarouge}{\bfseries Segment} & \textcolor{insarouge}{\bfseries Taille, raison d'achat} & \textcolor{insarouge}{\bfseries Rang}\\
        Radioamateurs, ADRASEC & 15\,836 indicatifs~; FNRASEC ≈ 1\,600 bénévoles, agrément sécurité civile. Prescripteurs. & 1\\
        Communes & ≈ 10\,800 avec plan communal de sauvegarde~; achat sous 40\,000\,€ HT sans concurrence ≈ 695 kits. & 1\\
        Randonneurs & 27 millions~; 22\,000 kits = 0,08\,\%. & 2\\
        Chasseurs & 963\,571 (2023), âge médian 55 ans. Liaison de battue (hypothèse). & 2\\
        Préparation & Guide « Tous responsables » (21/11/2025), kit de 72\,h. & 3\\
        Makers du mesh & Déjà équipés~: 312 nœuds visibles sur le moniteur Gaulix. & 3\\
      \end{tabularx}
    \end{column}
  \end{columns}
  \vspace{0.1cm}
  \begin{pointcle}\centering\footnotesize
    Vendre de la densité, pas des objets~: lots de 4 à 10 kits, d'abord là où existe une ADRASEC ou un club.
  \end{pointcle}
  \srcnote{Sources~: \etl{https://www.tibison.com/produit/accessoires/garmin-inreach-mini-2/}{tibison.com}, \etl{https://www.arcep.fr/uploads/tx_gspublication/Arcep-RA2026-tome2-Panorama-connectivite-mobile-France_normeA.pdf}{arcep.fr}, \etl{https://meshtastic.org/docs/configuration/radio/lora/}{meshtastic.org}, \etl{https://www.ref66.fr/2026/01/25/les-radioamateurs-en-france-en-2025-etat-des-lieux-selon-lobservatoire-de-lanfr/}{ref66.fr}, \etl{https://fr.wikipedia.org/wiki/Chasse_en_France}{fr.wikipedia.org}, \etl{https://www.maire-info.com/sa\%C2\%A9curita\%C2\%A9-civile/le-gouvernement-veut-rendre-obligatoire-le-plan-communal-de-sauvegarde-dans-toutes-les-communes-article-31060}{maire-info.com}, \etl{https://donnees-publiques.fr/guides/seuils-marches-publics-2026/}{donnees-publiques.fr}, \etl{https://www.ffrandonnee.fr/s-informer/actualites/avec-27-millions-de-pratiquants-la-randonnee-est-le-sport-le-plus-pratique-en-france}{ffrandonnee.fr}, \etl{https://www.generation-nt.com/actualites/guide-tous-responsables-gouvernement-crise-kit-survie-resilience-2066493}{generation-nt.com}, \etl{https://map.gaulix.fr/stats}{map.gaulix.fr}}
\end{frame}

%---------------------------------------------------------------------------
\begin{frame}{Où se situe \nomproduit{} face à l'offre existante~?}
  \centering\scriptsize\renewcommand{\arraystretch}{1.1}
  \begin{tabularx}{\linewidth}{@{}p{3.0cm}p{2.0cm}p{2.7cm}p{1.6cm}X@{}}
    \rowcolor{insarose}\textcolor{insarouge}{\bfseries Offre} & \textcolor{insarouge}{\bfseries Prix TTC} & \textcolor{insarouge}{\bfseries Abonnement} & \textcolor{insarouge}{\bfseries Secours} & \textcolor{insarouge}{\bfseries Réseau utilisé}\\
    \rowcolor{insabeige}\textbf{\nomproduit} & 69\,€ & non & non & voisins~; prêt à l'emploi\\
    Carte nue Heltec V3 & 18 à 30\,€ & non & non & voisins~; à régler soi-même\\
    SenseCAP T1000-E & ≈ 58,80\,€ & non & non & voisins\\
    WisMesh Pocket V2 & 99\,€ & non & non & voisins~; pré-flashé\\
    Garmin inReach Mini 2 & 319 à 350\,€ & 15 à 20\,€ par mois & oui, 24\,h/24 & satellites\\
    SOS satellite du smartphone & iPhone 14 ou plus & Apple~: gratuit 2 ans~; Orange~: 5\,€/mois & Apple~: oui & satellites\\
  \end{tabularx}
  \vspace{0.15cm}
  \begin{columns}[T]
    \begin{column}{0.325\textwidth}
      \begin{carte}[height=1.05cm,valign=top]
        \raggedright\gros{≈ 1\,219\,€}\\[1pt]{\footnotesize sur 5 ans pour un inReach.}
      \end{carte}
    \end{column}
    \begin{column}{0.325\textwidth}
      \begin{carte}[height=1.05cm,valign=top]
        \raggedright\gros{4 kits = 276\,€}\\[1pt]{\footnotesize moins qu'un inReach seul.}
      \end{carte}
    \end{column}
    \begin{column}{0.325\textwidth}
      \begin{carte}[height=1.05cm,valign=top]
        \raggedright\gros{Pixel 9 et 10}\\[1pt]{\footnotesize seuls éligibles au SMS Orange.}
      \end{carte}
    \end{column}
  \end{columns}
  \srcnote{Sources~: \etl{https://www.tibison.com/produit/accessoires/garmin-inreach-mini-2/}{tibison.com}, \etl{https://www.tech-connect.info/meshtastic-meshcore-communiquer-sans-internet-le-guide-complet-2026/}{tech-connect.info}, \etl{https://www.passion-radio.fr/mesh-lora/114993369-3198.html}{passion-radio.fr}, \etl{https://hexaspot.com/products/wismesh-pocket-v2-ready-to-use-meshtastic-device}{hexaspot.com}, \etl{https://support.apple.com/fr-afri/101573}{support.apple.com}, \etl{https://newsroom.orange.com/orange-devient-le-1er-operateur-europeen-a-proposer-une-offre-de-sms-par-satellite-en-sappuyant-sur-la-technologie-direct-to-device/}{newsroom.orange.com}}
\end{frame}

%---------------------------------------------------------------------------
\begin{frame}{Acheter, produire~: à quel coût~?}
  \begin{columns}[T]
    \begin{column}{0.53\textwidth}
      \footnotesize\renewcommand{\arraystretch}{1.05}
      \begin{tabularx}{\linewidth}{@{}Xr@{}}
        \rowcolor{insarose}\textcolor{insarouge}{\bfseries Par kit (année 3)} & \textcolor{insarouge}{\bfseries €}\\
        Carte Heltec V3 + antenne (acheter)  & 17,00\\
        Batterie Li-Po (acheter)             & 4,50\\
        Câble, emballage, carte (acheter)    & 2,80\\
        Boîtier (produire~: 3D, puis injection) & 2,50\\
        Assemblage, flash, test (produire)   & 3,50\\
        \rowcolor{insabeige}\textbf{Matériel} & \textbf{30,30}\\
        Vente 4,00 + kit solidaire 3,03      & 37,33\\
        \rowcolor{insabeige}\textbf{Marge par kit (65\,\% du HT en coût)} & \textbf{20,17}\\
      \end{tabularx}
    \end{column}
    \begin{column}{0.45\textwidth}
      \begin{encadre}{Une fois, et obligations}
        \begin{puces}
          \item Tests CE/RED~: 4 à 11\,k€, 3 à 6 mois.
          \item Moule d'injection~: 5 à 50\,k€~; stock ≈ 100\,k€.
          \item Éco-contributions DEEE~: ≈ 0,5 à 2,5\,€ par kit (hypothèse).
          \item Besoin total ≈ 450\,k€~; Bourse French Tech~: 30\,k€ au plus.
        \end{puces}
      \end{encadre}
    \end{column}
  \end{columns}
  \vspace{0.05cm}
  \begin{pointcle}\centering\footnotesize
    Logiciel à développer~: aucun (Meshtastic, GPL-3.0). On produit assemblage, configuration, support, conformité~: là sont les 20\,€.
  \end{pointcle}
  \srcnote{Sources~: \etl{https://heltec.org/project/wifi-lora-32-v3/}{heltec.org}, \etl{https://rf-market.fr/meshtastic/1441-batterie-lipo-37v-1200mah-jst-125mm-2-broches.html}{rf-market.fr}, \etl{https://www.prototypage3d.com/impression-3d-vs-injection-moulee/}{prototypage3d.com}, \etl{https://codium.fr/en/news/ce-certification-electronics/}{codium.fr}, \etl{https://packdeclare.com/fr/blog/filiere-deee-ecologic-2026}{packdeclare.com}, \etl{https://pro.ecosystem.eco/profil/producteur/obligations-producteur}{pro.ecosystem.eco}, \etl{https://compass-financement.com/obtenir-bourse-french-tech/}{compass-financement.com}, \etl{https://github.com/meshtastic/firmware}{github.com/meshtastic}. Chiffres du kit~: hypothèses du modèle financier.}
\end{frame}

%---------------------------------------------------------------------------
\begin{frame}{Freins propres au projet}
  \footnotesize\renewcommand{\arraystretch}{1.0}
  \begin{tabularx}{\linewidth}{@{}p{2.9cm}X@{}}
    \rowcolor{insarose}\textcolor{insarouge}{\bfseries Frein} & \textcolor{insarouge}{\bfseries Éléments sourcés}\\
    F1 Effet de réseau~: « personne n'en a » & Une balise marche seule, pas un nœud. Moniteur Gaulix~: 312 nœuds visibles, 76 actifs. Portée 1 à 3\,km en ville, 3 sauts par défaut (7 au plus).\\
    F2 Le smartphone fait déjà « SOS » & iPhone 14 ou plus~: gratuit 2 ans (France, 12/2022)~; SMS satellite Orange (Pixel 9 et 10, 5\,€ par mois, 12/2025)~; Starlink Mobile pas encore en France.\\
    F3 Aucun secours, mais responsabilité & Directive 2024/2853 (9/12/2026)~: logiciel = produit, clauses d'exclusion inopposables~; exception du libre limitée au non commercial.\\
    F4 Conformité d'un vendeur visible & RED~: déclaration en français, amende jusqu'à 7\,500\,€~; cybersécurité RED depuis le 1/8/2025~; CRA~: notification depuis le 11/9/2026~; ANSSI pour les envois dans l'UE~; V4 à 28\,dBm pour 500\,mW permis.\\
    F5 Marque et logiciel & « Meshtastic » est une marque, interdite dans le nom du produit~; \emph{grant} révocable~; GPL-3.0~; MeshCore (MIT) en concurrent.\\
    F6 Matériel, batterie & Fournisseur unique (V3 vers V4)~; rapport UN38.3~; Colissimo~: route en Europe, 4 batteries par colis.\\
    F7 Pas de revenu récurrent & Point mort ≈ 21\,300 kits par an~; −127\,k€ à 15\,000 kits~; les makers achètent la carte à 18\,€.\\
  \end{tabularx}
  \srcnote{Sources~: \etl{https://map.gaulix.fr/stats}{map.gaulix.fr}, \etl{https://meshtastic.org/docs/configuration/radio/lora/}{meshtastic.org}, \etl{https://www.apple.com/fr/newsroom/2022/12/emergency-sos-via-satellite-available-in-france-germany-ireland-and-the-uk/}{apple.com}, \etl{https://newsroom.orange.com/orange-devient-le-1er-operateur-europeen-a-proposer-une-offre-de-sms-par-satellite-en-sappuyant-sur-la-technologie-direct-to-device/}{newsroom.orange.com}, \etl{https://www.spacexfrance.com/starlink-mobile-en-france-votre-smartphone-bientot-connecte-aux-satellites/}{spacexfrance.com}, \etl{https://forge-agency.fr/blog/responsabilite-produits-defectueux-logiciel-2026-pme-guide-pratique}{forge-agency.fr}, \etl{https://www.anfr.fr/fileadmin/medias/brochures-plaquettes/ANFR-brochure-directive-RED.pdf}{anfr.fr}, \etl{https://www.nemko.com/blog/cybersecurity-in-europe-en-18031-is-now-a-harmonized-standard}{nemko.com}, \etl{https://digital-strategy.ec.europa.eu/en/policies/cra-reporting}{digital-strategy.ec.europa.eu}, \etl{https://cyber.gouv.fr/controle-relatif-un-moyen-de-cryptologie}{cyber.gouv.fr}, \etl{https://heltec.org/project/wifi-lora-32-v4/}{heltec.org}, \etl{https://meshtastic.org/docs/legal/licensing-and-trademark/}{meshtastic.org}, \etl{https://en.wikipedia.org/wiki/MeshCore}{en.wikipedia.org}, \etl{https://6-napse.com/actualites/un38-3-certification-transport-batteries-lithium/}{6-napse.com}, \etl{https://www.colissimo.entreprise.laposte.fr/en/practical-advice/right-packaging/almost-perfect-parcel-web-series/lithium-batteries}{colissimo.entreprise.laposte.fr}}
\end{frame}

%---------------------------------------------------------------------------
\begin{frame}{Opportunités que le reste du marché n'a pas}
  \begin{columns}[T]
    \begin{column}{0.325\textwidth}
      \begin{carte}[height=2.3cm,valign=top]
        \raggedright\gros{Zéro licence}\\[1pt]{\footnotesize Firmware libre, plus de 100 cartes~: R\&D mutualisée (≈ 11\,€ par kit pour 4 développeurs).}
      \end{carte}
    \end{column}
    \begin{column}{0.325\textwidth}
      \begin{carte}[height=2.3cm,valign=top]
        \raggedright\gros{Acheteur public}\\[1pt]{\footnotesize Plan communal de sauvegarde pour toutes les communes en 2029 (projet)~; marchés sous 40\,000\,€ HT sans concurrence.}
      \end{carte}
    \end{column}
    \begin{column}{0.325\textwidth}
      \begin{carte}[height=2.3cm,valign=top]
        \raggedright\gros{Bénévoles}\\[1pt]{\footnotesize FNRASEC ≈ 1\,600 membres, agrément de sécurité civile~; canal d'urgence Gaulix (Fr\_EMCOM).}
      \end{carte}
    \end{column}
  \end{columns}
  \vspace{0.05cm}
  \begin{columns}[T]
    \begin{column}{0.325\textwidth}
      \begin{carte}[height=2.3cm,valign=top]
        \raggedright\gros{Conformité}\\[1pt]{\footnotesize Petits colis hors UE~: +3\,€ depuis le 1/7/2026 (≈ +17\,\% sur une carte), frais de gestion dès le 1/11.}
      \end{carte}
    \end{column}
    \begin{column}{0.325\textwidth}
      \begin{carte}[height=2.3cm,valign=top]
        \raggedright\gros{Fenêtre 2028}\\[1pt]{\footnotesize Starlink Mobile absent de France, SMS Orange réservé à deux Pixel~; Telekom~: début 2028, hors France.}
      \end{carte}
    \end{column}
    \begin{column}{0.325\textwidth}
      \begin{carte}[height=2.3cm,valign=top]
        \raggedright\gros{Relocaliser}\\[1pt]{\footnotesize Origine France Garantie~: 50 à 100\,\% du coût en France. Nous estimons 40 à 45\,\% en phase 3 (hypothèse).}
      \end{carte}
    \end{column}
  \end{columns}
  \srcnote{Sources~: \etl{https://github.com/meshtastic/firmware}{github.com/meshtastic}, \etl{https://www.tech-connect.info/meshtastic-meshcore-communiquer-sans-internet-le-guide-complet-2026/}{tech-connect.info}, \etl{https://donnees-publiques.fr/guides/seuils-marches-publics-2026/}{donnees-publiques.fr}, \etl{https://www.la-resilience.fr/2025/10/meshtastic-vs-gaulix-4-4/}{la-resilience.fr}, \etl{https://www.economie.gouv.fr/particuliers/mes-droits-conso/bien-consommer/tout-savoir-sur-la-taxe-sur-les-petits-colis}{economie.gouv.fr}, \etl{https://www.spacexfrance.com/starlink-mobile-en-france-votre-smartphone-bientot-connecte-aux-satellites/}{spacexfrance.com}, \etl{https://www.entreprises.cci-paris-idf.fr/fiches-pratiques/vous-fabriquez-en-france-demandez-le-label-origine-france-garantie}{entreprises.cci-paris-idf.fr}}
\end{frame}

%---------------------------------------------------------------------------
\begin{frame}{Solutions aux freins}
  \footnotesize\renewcommand{\arraystretch}{1.0}
  \begin{tabularx}{\linewidth}{@{}p{1.1cm}X@{}}
    \rowcolor{insarose}\textcolor{insarouge}{\bfseries Frein} & \textcolor{insarouge}{\bfseries Solution}\\
    F1, F7 & \textbf{Vendre de la densité}~: lots de 4 (276\,€) et de 10~; démarrer où existe une ADRASEC ou un club~; revenus hors matériel (installation, formation, accessoires).\\
    F2 & \textbf{Complément, pas remplaçant}~: tout smartphone (application Android officielle), sans opérateur, en groupe privé~; aucune promesse de secours.\\
    F3 & \textbf{Prévention}~: test à 100\,\%, protection de la batterie, notice « ne remplace ni le 112 ni une balise », assurance responsabilité produit avant la première vente.\\
    F4 & \textbf{Conformité dès l'année 1} (25\,k€)~: laboratoire en pré-conformité, rapports de Heltec, EN 18031, procédure de vulnérabilités~; vendre d'abord en France~; brider à 27\,dBm.\\
    F5 & \textbf{« \nomproduit, compatible Meshtastic »} avec les logos M-Powered~; demander un \emph{grant}~; publier nos modifications~; matériel ouvert à MeshCore.\\
    F6 & \textbf{Double source V3/V4}, stock tampon ≈ 3\,300 kits, fret en lots, UN38.3 exigé, route en Europe, 4 kits par colis.\\
  \end{tabularx}
  \vspace{0.1cm}
  \begin{pointcle}\centering\footnotesize
    À valider~: devis de laboratoire, entretiens (ADRASEC, 5 mairies, 5 clubs), déclaration CE de Heltec, éco-contribution, prix de la batterie.
  \end{pointcle}
  \srcnote{Sources~: voir les diapositives précédentes~; \etl{https://github.com/meshtastic/Meshtastic-Android}{github.com/meshtastic}, \etl{https://meshtastic.org/docs/legal/trademark-grants/}{meshtastic.org}, \etl{https://innovel.fr/actualites/certification-red-da-en-18031/}{innovel.fr}. Les chiffres du modèle sont des hypothèses.}
\end{frame}

%  Fichier de contenu : ni préambule ni \begin{document}, commence par \section.
%  Les sources sont citées dans chaque diapositive par des liens cliquables (nom de domaine),
%  sans passer par sources.bib. Pages consultées le 5 octobre 2026.
%  Niveaux : les chiffres du modèle financier (69 € TTC, 22 000 kits, etc.) sont des hypothèses.
%%%%%%%%%%%%%%%%%%%%%%%%%%%%%%%%%%%%%%%%%%%%%%%%%%%%%%%%%%%%%%%%%%%%%%%%%%%%%%
\section{Étude documentaire}
\providecommand{\etl}[2]{\href{#1}{#2}}

%---------------------------------------------------------------------------
\begin{frame}{Que mettre en avant, et auprès de qui~?}
  \begin{columns}[T]
    \begin{column}{0.37\textwidth}
      \begin{puces}
        \item \textbf{Payé une fois}~: 69\,€ contre 319\,€ puis 15 à 20\,€ par mois (Garmin inReach Mini 2).
        \item \textbf{Marche en cas de panne}~: la 4G couvre 99,2 à 99,6\,\% de la population, « zone blanche » seul est faible.
        \item \textbf{Prêt à l'emploi, légal, en français} (869,4--869,65\,MHz, 27\,dBm, 10\,\%).
        \item \textbf{Jamais « balise de détresse »}~: texte seul, 1 à 3\,km en ville.
      \end{puces}
    \end{column}
    \begin{column}{0.61\textwidth}
      \scriptsize\renewcommand{\arraystretch}{1.05}
      \begin{tabularx}{\linewidth}{@{}p{2.2cm}Xc@{}}
        \rowcolor{insarose}\textcolor{insarouge}{\bfseries Segment} & \textcolor{insarouge}{\bfseries Taille, raison d'achat} & \textcolor{insarouge}{\bfseries Rang}\\
        Radioamateurs, ADRASEC & 15\,836 indicatifs~; FNRASEC ≈ 1\,600 bénévoles, agrément sécurité civile. Prescripteurs. & 1\\
        Communes & ≈ 10\,800 avec plan communal de sauvegarde~; achat sous 40\,000\,€ HT sans concurrence ≈ 695 kits. & 1\\
        Randonneurs & 27 millions~; 22\,000 kits = 0,08\,\%. & 2\\
        Chasseurs & 963\,571 (2023), âge médian 55 ans. Liaison de battue (hypothèse). & 2\\
        Préparation & Guide « Tous responsables » (21/11/2025), kit de 72\,h. & 3\\
        Makers du mesh & Déjà équipés~: 312 nœuds visibles sur le moniteur Gaulix. & 3\\
      \end{tabularx}
    \end{column}
  \end{columns}
  \vspace{0.1cm}
  \begin{pointcle}\centering\footnotesize
    Vendre de la densité, pas des objets~: lots de 4 à 10 kits, d'abord là où existe une ADRASEC ou un club.
  \end{pointcle}
  \srcnote{Sources~: \etl{https://www.tibison.com/produit/accessoires/garmin-inreach-mini-2/}{tibison.com}, \etl{https://www.arcep.fr/uploads/tx_gspublication/Arcep-RA2026-tome2-Panorama-connectivite-mobile-France_normeA.pdf}{arcep.fr}, \etl{https://meshtastic.org/docs/configuration/radio/lora/}{meshtastic.org}, \etl{https://www.ref66.fr/2026/01/25/les-radioamateurs-en-france-en-2025-etat-des-lieux-selon-lobservatoire-de-lanfr/}{ref66.fr}, \etl{https://fr.wikipedia.org/wiki/Chasse_en_France}{fr.wikipedia.org}, \etl{https://www.maire-info.com/sa\%C2\%A9curita\%C2\%A9-civile/le-gouvernement-veut-rendre-obligatoire-le-plan-communal-de-sauvegarde-dans-toutes-les-communes-article-31060}{maire-info.com}, \etl{https://donnees-publiques.fr/guides/seuils-marches-publics-2026/}{donnees-publiques.fr}, \etl{https://www.ffrandonnee.fr/s-informer/actualites/avec-27-millions-de-pratiquants-la-randonnee-est-le-sport-le-plus-pratique-en-france}{ffrandonnee.fr}, \etl{https://www.generation-nt.com/actualites/guide-tous-responsables-gouvernement-crise-kit-survie-resilience-2066493}{generation-nt.com}, \etl{https://map.gaulix.fr/stats}{map.gaulix.fr}}
\end{frame}

%---------------------------------------------------------------------------
\begin{frame}{Où se situe \nomproduit{} face à l'offre existante~?}
  \centering\scriptsize\renewcommand{\arraystretch}{1.1}
  \begin{tabularx}{\linewidth}{@{}p{3.0cm}p{2.0cm}p{2.7cm}p{1.6cm}X@{}}
    \rowcolor{insarose}\textcolor{insarouge}{\bfseries Offre} & \textcolor{insarouge}{\bfseries Prix TTC} & \textcolor{insarouge}{\bfseries Abonnement} & \textcolor{insarouge}{\bfseries Secours} & \textcolor{insarouge}{\bfseries Réseau utilisé}\\
    \rowcolor{insabeige}\textbf{\nomproduit} & 69\,€ & non & non & voisins~; prêt à l'emploi\\
    Carte nue Heltec V3 & 18 à 30\,€ & non & non & voisins~; à régler soi-même\\
    SenseCAP T1000-E & ≈ 58,80\,€ & non & non & voisins\\
    WisMesh Pocket V2 & 99\,€ & non & non & voisins~; pré-flashé\\
    Garmin inReach Mini 2 & 319 à 350\,€ & 15 à 20\,€ par mois & oui, 24\,h/24 & satellites\\
    SOS satellite du smartphone & iPhone 14 ou plus & Apple~: gratuit 2 ans~; Orange~: 5\,€/mois & Apple~: oui & satellites\\
  \end{tabularx}
  \vspace{0.15cm}
  \begin{columns}[T]
    \begin{column}{0.325\textwidth}
      \begin{carte}[height=1.05cm,valign=top]
        \raggedright\gros{≈ 1\,219\,€}\\[1pt]{\footnotesize sur 5 ans pour un inReach.}
      \end{carte}
    \end{column}
    \begin{column}{0.325\textwidth}
      \begin{carte}[height=1.05cm,valign=top]
        \raggedright\gros{4 kits = 276\,€}\\[1pt]{\footnotesize moins qu'un inReach seul.}
      \end{carte}
    \end{column}
    \begin{column}{0.325\textwidth}
      \begin{carte}[height=1.05cm,valign=top]
        \raggedright\gros{Pixel 9 et 10}\\[1pt]{\footnotesize seuls éligibles au SMS Orange.}
      \end{carte}
    \end{column}
  \end{columns}
  \srcnote{Sources~: \etl{https://www.tibison.com/produit/accessoires/garmin-inreach-mini-2/}{tibison.com}, \etl{https://www.tech-connect.info/meshtastic-meshcore-communiquer-sans-internet-le-guide-complet-2026/}{tech-connect.info}, \etl{https://www.passion-radio.fr/mesh-lora/114993369-3198.html}{passion-radio.fr}, \etl{https://hexaspot.com/products/wismesh-pocket-v2-ready-to-use-meshtastic-device}{hexaspot.com}, \etl{https://support.apple.com/fr-afri/101573}{support.apple.com}, \etl{https://newsroom.orange.com/orange-devient-le-1er-operateur-europeen-a-proposer-une-offre-de-sms-par-satellite-en-sappuyant-sur-la-technologie-direct-to-device/}{newsroom.orange.com}}
\end{frame}

%---------------------------------------------------------------------------
\begin{frame}{Acheter, produire~: à quel coût~?}
  \begin{columns}[T]
    \begin{column}{0.53\textwidth}
      \footnotesize\renewcommand{\arraystretch}{1.05}
      \begin{tabularx}{\linewidth}{@{}Xr@{}}
        \rowcolor{insarose}\textcolor{insarouge}{\bfseries Par kit (année 3)} & \textcolor{insarouge}{\bfseries €}\\
        Carte Heltec V3 + antenne (acheter)  & 17,00\\
        Batterie Li-Po (acheter)             & 4,50\\
        Câble, emballage, carte (acheter)    & 2,80\\
        Boîtier (produire~: 3D, puis injection) & 2,50\\
        Assemblage, flash, test (produire)   & 3,50\\
        \rowcolor{insabeige}\textbf{Matériel} & \textbf{30,30}\\
        Vente 4,00 + kit solidaire 3,03      & 37,33\\
        \rowcolor{insabeige}\textbf{Marge par kit (65\,\% du HT en coût)} & \textbf{20,17}\\
      \end{tabularx}
    \end{column}
    \begin{column}{0.45\textwidth}
      \begin{encadre}{Une fois, et obligations}
        \begin{puces}
          \item Tests CE/RED~: 4 à 11\,k€, 3 à 6 mois.
          \item Moule d'injection~: 5 à 50\,k€~; stock ≈ 100\,k€.
          \item Éco-contributions DEEE~: ≈ 0,5 à 2,5\,€ par kit (hypothèse).
          \item Besoin total ≈ 450\,k€~; Bourse French Tech~: 30\,k€ au plus.
        \end{puces}
      \end{encadre}
    \end{column}
  \end{columns}
  \vspace{0.05cm}
  \begin{pointcle}\centering\footnotesize
    Logiciel à développer~: aucun (Meshtastic, GPL-3.0). On produit assemblage, configuration, support, conformité~: là sont les 20\,€.
  \end{pointcle}
  \srcnote{Sources~: \etl{https://heltec.org/project/wifi-lora-32-v3/}{heltec.org}, \etl{https://rf-market.fr/meshtastic/1441-batterie-lipo-37v-1200mah-jst-125mm-2-broches.html}{rf-market.fr}, \etl{https://www.prototypage3d.com/impression-3d-vs-injection-moulee/}{prototypage3d.com}, \etl{https://codium.fr/en/news/ce-certification-electronics/}{codium.fr}, \etl{https://packdeclare.com/fr/blog/filiere-deee-ecologic-2026}{packdeclare.com}, \etl{https://pro.ecosystem.eco/profil/producteur/obligations-producteur}{pro.ecosystem.eco}, \etl{https://compass-financement.com/obtenir-bourse-french-tech/}{compass-financement.com}, \etl{https://github.com/meshtastic/firmware}{github.com/meshtastic}. Chiffres du kit~: hypothèses du modèle financier.}
\end{frame}

%---------------------------------------------------------------------------
\begin{frame}{Freins propres au projet}
  \footnotesize\renewcommand{\arraystretch}{1.0}
  \begin{tabularx}{\linewidth}{@{}p{2.9cm}X@{}}
    \rowcolor{insarose}\textcolor{insarouge}{\bfseries Frein} & \textcolor{insarouge}{\bfseries Éléments sourcés}\\
    F1 Effet de réseau~: « personne n'en a » & Une balise marche seule, pas un nœud. Moniteur Gaulix~: 312 nœuds visibles, 76 actifs. Portée 1 à 3\,km en ville, 3 sauts par défaut (7 au plus).\\
    F2 Le smartphone fait déjà « SOS » & iPhone 14 ou plus~: gratuit 2 ans (France, 12/2022)~; SMS satellite Orange (Pixel 9 et 10, 5\,€ par mois, 12/2025)~; Starlink Mobile pas encore en France.\\
    F3 Aucun secours, mais responsabilité & Directive 2024/2853 (9/12/2026)~: logiciel = produit, clauses d'exclusion inopposables~; exception du libre limitée au non commercial.\\
    F4 Conformité d'un vendeur visible & RED~: déclaration en français, amende jusqu'à 7\,500\,€~; cybersécurité RED depuis le 1/8/2025~; CRA~: notification depuis le 11/9/2026~; ANSSI pour les envois dans l'UE~; V4 à 28\,dBm pour 500\,mW permis.\\
    F5 Marque et logiciel & « Meshtastic » est une marque, interdite dans le nom du produit~; \emph{grant} révocable~; GPL-3.0~; MeshCore (MIT) en concurrent.\\
    F6 Matériel, batterie & Fournisseur unique (V3 vers V4)~; rapport UN38.3~; Colissimo~: route en Europe, 4 batteries par colis.\\
    F7 Pas de revenu récurrent & Point mort ≈ 21\,300 kits par an~; −127\,k€ à 15\,000 kits~; les makers achètent la carte à 18\,€.\\
  \end{tabularx}
  \srcnote{Sources~: \etl{https://map.gaulix.fr/stats}{map.gaulix.fr}, \etl{https://meshtastic.org/docs/configuration/radio/lora/}{meshtastic.org}, \etl{https://www.apple.com/fr/newsroom/2022/12/emergency-sos-via-satellite-available-in-france-germany-ireland-and-the-uk/}{apple.com}, \etl{https://newsroom.orange.com/orange-devient-le-1er-operateur-europeen-a-proposer-une-offre-de-sms-par-satellite-en-sappuyant-sur-la-technologie-direct-to-device/}{newsroom.orange.com}, \etl{https://www.spacexfrance.com/starlink-mobile-en-france-votre-smartphone-bientot-connecte-aux-satellites/}{spacexfrance.com}, \etl{https://forge-agency.fr/blog/responsabilite-produits-defectueux-logiciel-2026-pme-guide-pratique}{forge-agency.fr}, \etl{https://www.anfr.fr/fileadmin/medias/brochures-plaquettes/ANFR-brochure-directive-RED.pdf}{anfr.fr}, \etl{https://www.nemko.com/blog/cybersecurity-in-europe-en-18031-is-now-a-harmonized-standard}{nemko.com}, \etl{https://digital-strategy.ec.europa.eu/en/policies/cra-reporting}{digital-strategy.ec.europa.eu}, \etl{https://cyber.gouv.fr/controle-relatif-un-moyen-de-cryptologie}{cyber.gouv.fr}, \etl{https://heltec.org/project/wifi-lora-32-v4/}{heltec.org}, \etl{https://meshtastic.org/docs/legal/licensing-and-trademark/}{meshtastic.org}, \etl{https://en.wikipedia.org/wiki/MeshCore}{en.wikipedia.org}, \etl{https://6-napse.com/actualites/un38-3-certification-transport-batteries-lithium/}{6-napse.com}, \etl{https://www.colissimo.entreprise.laposte.fr/en/practical-advice/right-packaging/almost-perfect-parcel-web-series/lithium-batteries}{colissimo.entreprise.laposte.fr}}
\end{frame}

%---------------------------------------------------------------------------
\begin{frame}{Opportunités que le reste du marché n'a pas}
  \begin{columns}[T]
    \begin{column}{0.325\textwidth}
      \begin{carte}[height=2.3cm,valign=top]
        \raggedright\gros{Zéro licence}\\[1pt]{\footnotesize Firmware libre, plus de 100 cartes~: R\&D mutualisée (≈ 11\,€ par kit pour 4 développeurs).}
      \end{carte}
    \end{column}
    \begin{column}{0.325\textwidth}
      \begin{carte}[height=2.3cm,valign=top]
        \raggedright\gros{Acheteur public}\\[1pt]{\footnotesize Plan communal de sauvegarde pour toutes les communes en 2029 (projet)~; marchés sous 40\,000\,€ HT sans concurrence.}
      \end{carte}
    \end{column}
    \begin{column}{0.325\textwidth}
      \begin{carte}[height=2.3cm,valign=top]
        \raggedright\gros{Bénévoles}\\[1pt]{\footnotesize FNRASEC ≈ 1\,600 membres, agrément de sécurité civile~; canal d'urgence Gaulix (Fr\_EMCOM).}
      \end{carte}
    \end{column}
  \end{columns}
  \vspace{0.05cm}
  \begin{columns}[T]
    \begin{column}{0.325\textwidth}
      \begin{carte}[height=2.3cm,valign=top]
        \raggedright\gros{Conformité}\\[1pt]{\footnotesize Petits colis hors UE~: +3\,€ depuis le 1/7/2026 (≈ +17\,\% sur une carte), frais de gestion dès le 1/11.}
      \end{carte}
    \end{column}
    \begin{column}{0.325\textwidth}
      \begin{carte}[height=2.3cm,valign=top]
        \raggedright\gros{Fenêtre 2028}\\[1pt]{\footnotesize Starlink Mobile absent de France, SMS Orange réservé à deux Pixel~; Telekom~: début 2028, hors France.}
      \end{carte}
    \end{column}
    \begin{column}{0.325\textwidth}
      \begin{carte}[height=2.3cm,valign=top]
        \raggedright\gros{Relocaliser}\\[1pt]{\footnotesize Origine France Garantie~: 50 à 100\,\% du coût en France. Nous estimons 40 à 45\,\% en phase 3 (hypothèse).}
      \end{carte}
    \end{column}
  \end{columns}
  \srcnote{Sources~: \etl{https://github.com/meshtastic/firmware}{github.com/meshtastic}, \etl{https://www.tech-connect.info/meshtastic-meshcore-communiquer-sans-internet-le-guide-complet-2026/}{tech-connect.info}, \etl{https://donnees-publiques.fr/guides/seuils-marches-publics-2026/}{donnees-publiques.fr}, \etl{https://www.la-resilience.fr/2025/10/meshtastic-vs-gaulix-4-4/}{la-resilience.fr}, \etl{https://www.economie.gouv.fr/particuliers/mes-droits-conso/bien-consommer/tout-savoir-sur-la-taxe-sur-les-petits-colis}{economie.gouv.fr}, \etl{https://www.spacexfrance.com/starlink-mobile-en-france-votre-smartphone-bientot-connecte-aux-satellites/}{spacexfrance.com}, \etl{https://www.entreprises.cci-paris-idf.fr/fiches-pratiques/vous-fabriquez-en-france-demandez-le-label-origine-france-garantie}{entreprises.cci-paris-idf.fr}}
\end{frame}

%---------------------------------------------------------------------------
\begin{frame}{Solutions aux freins}
  \footnotesize\renewcommand{\arraystretch}{1.0}
  \begin{tabularx}{\linewidth}{@{}p{1.1cm}X@{}}
    \rowcolor{insarose}\textcolor{insarouge}{\bfseries Frein} & \textcolor{insarouge}{\bfseries Solution}\\
    F1, F7 & \textbf{Vendre de la densité}~: lots de 4 (276\,€) et de 10~; démarrer où existe une ADRASEC ou un club~; revenus hors matériel (installation, formation, accessoires).\\
    F2 & \textbf{Complément, pas remplaçant}~: tout smartphone (application Android officielle), sans opérateur, en groupe privé~; aucune promesse de secours.\\
    F3 & \textbf{Prévention}~: test à 100\,\%, protection de la batterie, notice « ne remplace ni le 112 ni une balise », assurance responsabilité produit avant la première vente.\\
    F4 & \textbf{Conformité dès l'année 1} (25\,k€)~: laboratoire en pré-conformité, rapports de Heltec, EN 18031, procédure de vulnérabilités~; vendre d'abord en France~; brider à 27\,dBm.\\
    F5 & \textbf{« \nomproduit, compatible Meshtastic »} avec les logos M-Powered~; demander un \emph{grant}~; publier nos modifications~; matériel ouvert à MeshCore.\\
    F6 & \textbf{Double source V3/V4}, stock tampon ≈ 3\,300 kits, fret en lots, UN38.3 exigé, route en Europe, 4 kits par colis.\\
  \end{tabularx}
  \vspace{0.1cm}
  \begin{pointcle}\centering\footnotesize
    À valider~: devis de laboratoire, entretiens (ADRASEC, 5 mairies, 5 clubs), déclaration CE de Heltec, éco-contribution, prix de la batterie.
  \end{pointcle}
  \srcnote{Sources~: voir les diapositives précédentes~; \etl{https://github.com/meshtastic/Meshtastic-Android}{github.com/meshtastic}, \etl{https://meshtastic.org/docs/legal/trademark-grants/}{meshtastic.org}, \etl{https://innovel.fr/actualites/certification-red-da-en-18031/}{innovel.fr}. Les chiffres du modèle sont des hypothèses.}
\end{frame}

%  Fichier de contenu : ni préambule ni \begin{document}, commence par \section.
%  Les sources sont citées dans chaque diapositive par des liens cliquables (nom de domaine),
%  sans passer par sources.bib. Pages consultées le 5 octobre 2026.
%  Niveaux : les chiffres du modèle financier (69 € TTC, 22 000 kits, etc.) sont des hypothèses.
%%%%%%%%%%%%%%%%%%%%%%%%%%%%%%%%%%%%%%%%%%%%%%%%%%%%%%%%%%%%%%%%%%%%%%%%%%%%%%
\section{Étude documentaire}
\providecommand{\etl}[2]{\href{#1}{#2}}

%---------------------------------------------------------------------------
\begin{frame}{Que mettre en avant, et auprès de qui~?}
  \begin{columns}[T]
    \begin{column}{0.37\textwidth}
      \begin{puces}
        \item \textbf{Payé une fois}~: 69\,€ contre 319\,€ puis 15 à 20\,€ par mois (Garmin inReach Mini 2).
        \item \textbf{Marche en cas de panne}~: la 4G couvre 99,2 à 99,6\,\% de la population, « zone blanche » seul est faible.
        \item \textbf{Prêt à l'emploi, légal, en français} (869,4--869,65\,MHz, 27\,dBm, 10\,\%).
        \item \textbf{Jamais « balise de détresse »}~: texte seul, 1 à 3\,km en ville.
      \end{puces}
    \end{column}
    \begin{column}{0.61\textwidth}
      \scriptsize\renewcommand{\arraystretch}{1.05}
      \begin{tabularx}{\linewidth}{@{}p{2.2cm}Xc@{}}
        \rowcolor{insarose}\textcolor{insarouge}{\bfseries Segment} & \textcolor{insarouge}{\bfseries Taille, raison d'achat} & \textcolor{insarouge}{\bfseries Rang}\\
        Radioamateurs, ADRASEC & 15\,836 indicatifs~; FNRASEC ≈ 1\,600 bénévoles, agrément sécurité civile. Prescripteurs. & 1\\
        Communes & ≈ 10\,800 avec plan communal de sauvegarde~; achat sous 40\,000\,€ HT sans concurrence ≈ 695 kits. & 1\\
        Randonneurs & 27 millions~; 22\,000 kits = 0,08\,\%. & 2\\
        Chasseurs & 963\,571 (2023), âge médian 55 ans. Liaison de battue (hypothèse). & 2\\
        Préparation & Guide « Tous responsables » (21/11/2025), kit de 72\,h. & 3\\
        Makers du mesh & Déjà équipés~: 312 nœuds visibles sur le moniteur Gaulix. & 3\\
      \end{tabularx}
    \end{column}
  \end{columns}
  \vspace{0.1cm}
  \begin{pointcle}\centering\footnotesize
    Vendre de la densité, pas des objets~: lots de 4 à 10 kits, d'abord là où existe une ADRASEC ou un club.
  \end{pointcle}
  \srcnote{Sources~: \etl{https://www.tibison.com/produit/accessoires/garmin-inreach-mini-2/}{tibison.com}, \etl{https://www.arcep.fr/uploads/tx_gspublication/Arcep-RA2026-tome2-Panorama-connectivite-mobile-France_normeA.pdf}{arcep.fr}, \etl{https://meshtastic.org/docs/configuration/radio/lora/}{meshtastic.org}, \etl{https://www.ref66.fr/2026/01/25/les-radioamateurs-en-france-en-2025-etat-des-lieux-selon-lobservatoire-de-lanfr/}{ref66.fr}, \etl{https://fr.wikipedia.org/wiki/Chasse_en_France}{fr.wikipedia.org}, \etl{https://www.maire-info.com/sa\%C2\%A9curita\%C2\%A9-civile/le-gouvernement-veut-rendre-obligatoire-le-plan-communal-de-sauvegarde-dans-toutes-les-communes-article-31060}{maire-info.com}, \etl{https://donnees-publiques.fr/guides/seuils-marches-publics-2026/}{donnees-publiques.fr}, \etl{https://www.ffrandonnee.fr/s-informer/actualites/avec-27-millions-de-pratiquants-la-randonnee-est-le-sport-le-plus-pratique-en-france}{ffrandonnee.fr}, \etl{https://www.generation-nt.com/actualites/guide-tous-responsables-gouvernement-crise-kit-survie-resilience-2066493}{generation-nt.com}, \etl{https://map.gaulix.fr/stats}{map.gaulix.fr}}
\end{frame}

%---------------------------------------------------------------------------
\begin{frame}{Où se situe \nomproduit{} face à l'offre existante~?}
  \centering\scriptsize\renewcommand{\arraystretch}{1.1}
  \begin{tabularx}{\linewidth}{@{}p{3.0cm}p{2.0cm}p{2.7cm}p{1.6cm}X@{}}
    \rowcolor{insarose}\textcolor{insarouge}{\bfseries Offre} & \textcolor{insarouge}{\bfseries Prix TTC} & \textcolor{insarouge}{\bfseries Abonnement} & \textcolor{insarouge}{\bfseries Secours} & \textcolor{insarouge}{\bfseries Réseau utilisé}\\
    \rowcolor{insabeige}\textbf{\nomproduit} & 69\,€ & non & non & voisins~; prêt à l'emploi\\
    Carte nue Heltec V3 & 18 à 30\,€ & non & non & voisins~; à régler soi-même\\
    SenseCAP T1000-E & ≈ 58,80\,€ & non & non & voisins\\
    WisMesh Pocket V2 & 99\,€ & non & non & voisins~; pré-flashé\\
    Garmin inReach Mini 2 & 319 à 350\,€ & 15 à 20\,€ par mois & oui, 24\,h/24 & satellites\\
    SOS satellite du smartphone & iPhone 14 ou plus & Apple~: gratuit 2 ans~; Orange~: 5\,€/mois & Apple~: oui & satellites\\
  \end{tabularx}
  \vspace{0.15cm}
  \begin{columns}[T]
    \begin{column}{0.325\textwidth}
      \begin{carte}[height=1.05cm,valign=top]
        \raggedright\gros{≈ 1\,219\,€}\\[1pt]{\footnotesize sur 5 ans pour un inReach.}
      \end{carte}
    \end{column}
    \begin{column}{0.325\textwidth}
      \begin{carte}[height=1.05cm,valign=top]
        \raggedright\gros{4 kits = 276\,€}\\[1pt]{\footnotesize moins qu'un inReach seul.}
      \end{carte}
    \end{column}
    \begin{column}{0.325\textwidth}
      \begin{carte}[height=1.05cm,valign=top]
        \raggedright\gros{Pixel 9 et 10}\\[1pt]{\footnotesize seuls éligibles au SMS Orange.}
      \end{carte}
    \end{column}
  \end{columns}
  \srcnote{Sources~: \etl{https://www.tibison.com/produit/accessoires/garmin-inreach-mini-2/}{tibison.com}, \etl{https://www.tech-connect.info/meshtastic-meshcore-communiquer-sans-internet-le-guide-complet-2026/}{tech-connect.info}, \etl{https://www.passion-radio.fr/mesh-lora/114993369-3198.html}{passion-radio.fr}, \etl{https://hexaspot.com/products/wismesh-pocket-v2-ready-to-use-meshtastic-device}{hexaspot.com}, \etl{https://support.apple.com/fr-afri/101573}{support.apple.com}, \etl{https://newsroom.orange.com/orange-devient-le-1er-operateur-europeen-a-proposer-une-offre-de-sms-par-satellite-en-sappuyant-sur-la-technologie-direct-to-device/}{newsroom.orange.com}}
\end{frame}

%---------------------------------------------------------------------------
\begin{frame}{Acheter, produire~: à quel coût~?}
  \begin{columns}[T]
    \begin{column}{0.53\textwidth}
      \footnotesize\renewcommand{\arraystretch}{1.05}
      \begin{tabularx}{\linewidth}{@{}Xr@{}}
        \rowcolor{insarose}\textcolor{insarouge}{\bfseries Par kit (année 3)} & \textcolor{insarouge}{\bfseries €}\\
        Carte Heltec V3 + antenne (acheter)  & 17,00\\
        Batterie Li-Po (acheter)             & 4,50\\
        Câble, emballage, carte (acheter)    & 2,80\\
        Boîtier (produire~: 3D, puis injection) & 2,50\\
        Assemblage, flash, test (produire)   & 3,50\\
        \rowcolor{insabeige}\textbf{Matériel} & \textbf{30,30}\\
        Vente 4,00 + kit solidaire 3,03      & 37,33\\
        \rowcolor{insabeige}\textbf{Marge par kit (65\,\% du HT en coût)} & \textbf{20,17}\\
      \end{tabularx}
    \end{column}
    \begin{column}{0.45\textwidth}
      \begin{encadre}{Une fois, et obligations}
        \begin{puces}
          \item Tests CE/RED~: 4 à 11\,k€, 3 à 6 mois.
          \item Moule d'injection~: 5 à 50\,k€~; stock ≈ 100\,k€.
          \item Éco-contributions DEEE~: ≈ 0,5 à 2,5\,€ par kit (hypothèse).
          \item Besoin total ≈ 450\,k€~; Bourse French Tech~: 30\,k€ au plus.
        \end{puces}
      \end{encadre}
    \end{column}
  \end{columns}
  \vspace{0.05cm}
  \begin{pointcle}\centering\footnotesize
    Logiciel à développer~: aucun (Meshtastic, GPL-3.0). On produit assemblage, configuration, support, conformité~: là sont les 20\,€.
  \end{pointcle}
  \srcnote{Sources~: \etl{https://heltec.org/project/wifi-lora-32-v3/}{heltec.org}, \etl{https://rf-market.fr/meshtastic/1441-batterie-lipo-37v-1200mah-jst-125mm-2-broches.html}{rf-market.fr}, \etl{https://www.prototypage3d.com/impression-3d-vs-injection-moulee/}{prototypage3d.com}, \etl{https://codium.fr/en/news/ce-certification-electronics/}{codium.fr}, \etl{https://packdeclare.com/fr/blog/filiere-deee-ecologic-2026}{packdeclare.com}, \etl{https://pro.ecosystem.eco/profil/producteur/obligations-producteur}{pro.ecosystem.eco}, \etl{https://compass-financement.com/obtenir-bourse-french-tech/}{compass-financement.com}, \etl{https://github.com/meshtastic/firmware}{github.com/meshtastic}. Chiffres du kit~: hypothèses du modèle financier.}
\end{frame}

%---------------------------------------------------------------------------
\begin{frame}{Freins propres au projet}
  \footnotesize\renewcommand{\arraystretch}{1.0}
  \begin{tabularx}{\linewidth}{@{}p{2.9cm}X@{}}
    \rowcolor{insarose}\textcolor{insarouge}{\bfseries Frein} & \textcolor{insarouge}{\bfseries Éléments sourcés}\\
    F1 Effet de réseau~: « personne n'en a » & Une balise marche seule, pas un nœud. Moniteur Gaulix~: 312 nœuds visibles, 76 actifs. Portée 1 à 3\,km en ville, 3 sauts par défaut (7 au plus).\\
    F2 Le smartphone fait déjà « SOS » & iPhone 14 ou plus~: gratuit 2 ans (France, 12/2022)~; SMS satellite Orange (Pixel 9 et 10, 5\,€ par mois, 12/2025)~; Starlink Mobile pas encore en France.\\
    F3 Aucun secours, mais responsabilité & Directive 2024/2853 (9/12/2026)~: logiciel = produit, clauses d'exclusion inopposables~; exception du libre limitée au non commercial.\\
    F4 Conformité d'un vendeur visible & RED~: déclaration en français, amende jusqu'à 7\,500\,€~; cybersécurité RED depuis le 1/8/2025~; CRA~: notification depuis le 11/9/2026~; ANSSI pour les envois dans l'UE~; V4 à 28\,dBm pour 500\,mW permis.\\
    F5 Marque et logiciel & « Meshtastic » est une marque, interdite dans le nom du produit~; \emph{grant} révocable~; GPL-3.0~; MeshCore (MIT) en concurrent.\\
    F6 Matériel, batterie & Fournisseur unique (V3 vers V4)~; rapport UN38.3~; Colissimo~: route en Europe, 4 batteries par colis.\\
    F7 Pas de revenu récurrent & Point mort ≈ 21\,300 kits par an~; −127\,k€ à 15\,000 kits~; les makers achètent la carte à 18\,€.\\
  \end{tabularx}
  \srcnote{Sources~: \etl{https://map.gaulix.fr/stats}{map.gaulix.fr}, \etl{https://meshtastic.org/docs/configuration/radio/lora/}{meshtastic.org}, \etl{https://www.apple.com/fr/newsroom/2022/12/emergency-sos-via-satellite-available-in-france-germany-ireland-and-the-uk/}{apple.com}, \etl{https://newsroom.orange.com/orange-devient-le-1er-operateur-europeen-a-proposer-une-offre-de-sms-par-satellite-en-sappuyant-sur-la-technologie-direct-to-device/}{newsroom.orange.com}, \etl{https://www.spacexfrance.com/starlink-mobile-en-france-votre-smartphone-bientot-connecte-aux-satellites/}{spacexfrance.com}, \etl{https://forge-agency.fr/blog/responsabilite-produits-defectueux-logiciel-2026-pme-guide-pratique}{forge-agency.fr}, \etl{https://www.anfr.fr/fileadmin/medias/brochures-plaquettes/ANFR-brochure-directive-RED.pdf}{anfr.fr}, \etl{https://www.nemko.com/blog/cybersecurity-in-europe-en-18031-is-now-a-harmonized-standard}{nemko.com}, \etl{https://digital-strategy.ec.europa.eu/en/policies/cra-reporting}{digital-strategy.ec.europa.eu}, \etl{https://cyber.gouv.fr/controle-relatif-un-moyen-de-cryptologie}{cyber.gouv.fr}, \etl{https://heltec.org/project/wifi-lora-32-v4/}{heltec.org}, \etl{https://meshtastic.org/docs/legal/licensing-and-trademark/}{meshtastic.org}, \etl{https://en.wikipedia.org/wiki/MeshCore}{en.wikipedia.org}, \etl{https://6-napse.com/actualites/un38-3-certification-transport-batteries-lithium/}{6-napse.com}, \etl{https://www.colissimo.entreprise.laposte.fr/en/practical-advice/right-packaging/almost-perfect-parcel-web-series/lithium-batteries}{colissimo.entreprise.laposte.fr}}
\end{frame}

%---------------------------------------------------------------------------
\begin{frame}{Opportunités que le reste du marché n'a pas}
  \begin{columns}[T]
    \begin{column}{0.325\textwidth}
      \begin{carte}[height=2.3cm,valign=top]
        \raggedright\gros{Zéro licence}\\[1pt]{\footnotesize Firmware libre, plus de 100 cartes~: R\&D mutualisée (≈ 11\,€ par kit pour 4 développeurs).}
      \end{carte}
    \end{column}
    \begin{column}{0.325\textwidth}
      \begin{carte}[height=2.3cm,valign=top]
        \raggedright\gros{Acheteur public}\\[1pt]{\footnotesize Plan communal de sauvegarde pour toutes les communes en 2029 (projet)~; marchés sous 40\,000\,€ HT sans concurrence.}
      \end{carte}
    \end{column}
    \begin{column}{0.325\textwidth}
      \begin{carte}[height=2.3cm,valign=top]
        \raggedright\gros{Bénévoles}\\[1pt]{\footnotesize FNRASEC ≈ 1\,600 membres, agrément de sécurité civile~; canal d'urgence Gaulix (Fr\_EMCOM).}
      \end{carte}
    \end{column}
  \end{columns}
  \vspace{0.05cm}
  \begin{columns}[T]
    \begin{column}{0.325\textwidth}
      \begin{carte}[height=2.3cm,valign=top]
        \raggedright\gros{Conformité}\\[1pt]{\footnotesize Petits colis hors UE~: +3\,€ depuis le 1/7/2026 (≈ +17\,\% sur une carte), frais de gestion dès le 1/11.}
      \end{carte}
    \end{column}
    \begin{column}{0.325\textwidth}
      \begin{carte}[height=2.3cm,valign=top]
        \raggedright\gros{Fenêtre 2028}\\[1pt]{\footnotesize Starlink Mobile absent de France, SMS Orange réservé à deux Pixel~; Telekom~: début 2028, hors France.}
      \end{carte}
    \end{column}
    \begin{column}{0.325\textwidth}
      \begin{carte}[height=2.3cm,valign=top]
        \raggedright\gros{Relocaliser}\\[1pt]{\footnotesize Origine France Garantie~: 50 à 100\,\% du coût en France. Nous estimons 40 à 45\,\% en phase 3 (hypothèse).}
      \end{carte}
    \end{column}
  \end{columns}
  \srcnote{Sources~: \etl{https://github.com/meshtastic/firmware}{github.com/meshtastic}, \etl{https://www.tech-connect.info/meshtastic-meshcore-communiquer-sans-internet-le-guide-complet-2026/}{tech-connect.info}, \etl{https://donnees-publiques.fr/guides/seuils-marches-publics-2026/}{donnees-publiques.fr}, \etl{https://www.la-resilience.fr/2025/10/meshtastic-vs-gaulix-4-4/}{la-resilience.fr}, \etl{https://www.economie.gouv.fr/particuliers/mes-droits-conso/bien-consommer/tout-savoir-sur-la-taxe-sur-les-petits-colis}{economie.gouv.fr}, \etl{https://www.spacexfrance.com/starlink-mobile-en-france-votre-smartphone-bientot-connecte-aux-satellites/}{spacexfrance.com}, \etl{https://www.entreprises.cci-paris-idf.fr/fiches-pratiques/vous-fabriquez-en-france-demandez-le-label-origine-france-garantie}{entreprises.cci-paris-idf.fr}}
\end{frame}

%---------------------------------------------------------------------------
\begin{frame}{Solutions aux freins}
  \footnotesize\renewcommand{\arraystretch}{1.0}
  \begin{tabularx}{\linewidth}{@{}p{1.1cm}X@{}}
    \rowcolor{insarose}\textcolor{insarouge}{\bfseries Frein} & \textcolor{insarouge}{\bfseries Solution}\\
    F1, F7 & \textbf{Vendre de la densité}~: lots de 4 (276\,€) et de 10~; démarrer où existe une ADRASEC ou un club~; revenus hors matériel (installation, formation, accessoires).\\
    F2 & \textbf{Complément, pas remplaçant}~: tout smartphone (application Android officielle), sans opérateur, en groupe privé~; aucune promesse de secours.\\
    F3 & \textbf{Prévention}~: test à 100\,\%, protection de la batterie, notice « ne remplace ni le 112 ni une balise », assurance responsabilité produit avant la première vente.\\
    F4 & \textbf{Conformité dès l'année 1} (25\,k€)~: laboratoire en pré-conformité, rapports de Heltec, EN 18031, procédure de vulnérabilités~; vendre d'abord en France~; brider à 27\,dBm.\\
    F5 & \textbf{« \nomproduit, compatible Meshtastic »} avec les logos M-Powered~; demander un \emph{grant}~; publier nos modifications~; matériel ouvert à MeshCore.\\
    F6 & \textbf{Double source V3/V4}, stock tampon ≈ 3\,300 kits, fret en lots, UN38.3 exigé, route en Europe, 4 kits par colis.\\
  \end{tabularx}
  \vspace{0.1cm}
  \begin{pointcle}\centering\footnotesize
    À valider~: devis de laboratoire, entretiens (ADRASEC, 5 mairies, 5 clubs), déclaration CE de Heltec, éco-contribution, prix de la batterie.
  \end{pointcle}
  \srcnote{Sources~: voir les diapositives précédentes~; \etl{https://github.com/meshtastic/Meshtastic-Android}{github.com/meshtastic}, \etl{https://meshtastic.org/docs/legal/trademark-grants/}{meshtastic.org}, \etl{https://innovel.fr/actualites/certification-red-da-en-18031/}{innovel.fr}. Les chiffres du modèle sont des hypothèses.}
\end{frame}
                  % 8.  Étude documentaire : coûts, freins F1 à F7, opportunités
%%%%%%%%%%%%%%%%%%%%%%%%%%%%%%%%%%%%%%%%%%%%%%%%%%%%%%%%%%%%%%%%%%%%%%%%%%%%%%
%  Partie « SWOT et enjeux » (6 diapositives) -- à importer avec %%%%%%%%%%%%%%%%%%%%%%%%%%%%%%%%%%%%%%%%%%%%%%%%%%%%%%%%%%%%%%%%%%%%%%%%%%%%%%
%  Partie « SWOT et enjeux » (6 diapositives) -- à importer avec %%%%%%%%%%%%%%%%%%%%%%%%%%%%%%%%%%%%%%%%%%%%%%%%%%%%%%%%%%%%%%%%%%%%%%%%%%%%%%
%  Partie « SWOT et enjeux » (6 diapositives) -- à importer avec \include{chap_swot}
%  Fichier de contenu : ni préambule ni \begin{document}, commence par \section.
%
%  Remplace l'ancien SWOT en PDF (\includepdf supprimé) : le diagnostic est maintenant écrit
%  en LaTeX/TikZ, avec la charte du reste de la présentation, et complété par un SWOT croisé.
%  Ce qui a changé par rapport à la version PDF :
%   - la décomposition des 69 € n'est plus répétée ici (elle est dans la partie Finances) ;
%   - chaque point du diagnostic renvoie à un « frein » F1 à F7 de l'étude documentaire ;
%   - les marges des offres Solo / Duo / Groupe sont recalculées (marge par commande ET par kit,
%     avec le coût de vente de 4 € compté une seule fois par commande) ;
%   - le SWOT croisé (stratégies SO, WO, ST, WT) relie le diagnostic aux trois enjeux.
%  Hypothèses chiffrées : celles de la partie Finances (kit 69 € TTC, matériel 30,30 €, vente 4 €
%  par commande, 1 kit offert pour 10, charges fixes de l'année 3 = 430 k€).
%  Les offres Solo (79 €), Duo (139 €, 2 kits) et Groupe (249 €, 4 kits) sont celles du SWOT d'origine ;
%  le nombre de kits du Groupe (4) est une hypothèse.
%%%%%%%%%%%%%%%%%%%%%%%%%%%%%%%%%%%%%%%%%%%%%%%%%%%%%%%%%%%%%%%%%%%%%%%%%%%%%%
\section{SWOT et enjeux}

% Entrée de SWOT : \swit[couleur du code]{code}{titre}{texte}  (titre et texte sur le même paragraphe)
\newcommand{\swit}[4][insarouge]{%
  \begin{tcolorbox}[colback=insabeige,colframe=insabeige,boxrule=0pt,arc=2pt,left=4pt,right=4pt,
                    top=1.5pt,bottom=1.5pt,before skip=0pt,after skip=2pt]
    \scriptsize\raggedright{\sffamily\bfseries\color{#1}#2}\hspace{0.3em}{\sffamily\bfseries\footnotesize #3}\hspace{0.3em}#4\par
  \end{tcolorbox}}
% En-tête de colonne : \swhead{Forces}{couleur}
\newcommand{\swhead}[2]{%
  {\sffamily\bfseries\footnotesize\color{#2}\MakeUppercase{#1}}\\[-3pt]{\color{#2}\rule{\linewidth}{0.6pt}}\\[-1pt]}
% renvoi vers un enjeu ou un frein
\newcommand{\swref}[1]{\textcolor{insarouge}{\textbf{→\,#1}}}

%---------------------------------------------------------------------------
\begin{frame}{Diagnostic interne~: forces et faiblesses}
  \begin{columns}[T]
    \begin{column}{0.485\textwidth}
      \swhead{Forces}{insarouge}
      \swit{S1}{Proposition de valeur claire}{Une messagerie texte sans réseau mobile, sans compte ni abonnement.}
      \swit{S2}{Coût d'entrée faible}{69\,€ une fois, contre ≈ 830\,€ sur trois ans pour une balise satellite.}
      \swit{S3}{Écosystème libre existant}{Meshtastic, Heltec V3, applications~: l'effort porte sur l'usage.}
      \swit{S4}{Simplicité comme différence}{Réglages Europe, clé privée, QR code, guides en français.}
      \swit{S5}{Un socle, plusieurs usages}{Outdoor, campagne, ONG, crise~: un kit, quatre terrains.}
    \end{column}
    \begin{column}{0.485\textwidth}
      \swhead{Faiblesses}{insagris}
      \swit[insagris]{W1}{Faible barrière technologique}{Un bricoleur refait un kit proche pour 18 à 30\,€. \swref{F5}}
      \swit[insagris]{W2}{Effet de réseau indispensable}{Un kit isolé ne sert à rien~: la valeur vient des voisins. \swref{enjeu 1}}
      \swit[insagris]{W3}{Marge de contribution fine}{20,17\,€ par kit, mais 0,62\,€ de résultat par kit en année 3. \swref{enjeu 2}}
      \swit[insagris]{W4}{Financement avant rentabilité}{≈ 450\,k€ à financer, point mort ≈ 21\,300 kits par an. \swref{F7}}
      \swit[insagris]{W5}{Positionnement trop large}{ONG, randonneurs, makers, public~: attentes différentes. \swref{enjeu 3}}
    \end{column}
  \end{columns}
  \vspace{-0.2cm}
  \begin{pcompact}\centering\footnotesize
    Lecture clé~: la technologie est accessible~; la valeur viendra de l'usage, de la marque et de l'écosystème.
  \end{pcompact}
  \srcnote{Synthèse des parties Concurrence, Finances et Étude documentaire. F2 à F7 = freins de l'étude documentaire.}
\end{frame}

%---------------------------------------------------------------------------
\begin{frame}{Diagnostic externe~: opportunités et menaces}
  \begin{columns}[T]
    \begin{column}{0.485\textwidth}
      \swhead{Opportunités}{insarouge}
      \swit{O1}{Un marché outdoor important}{27 millions de pratiquants, dont 240\,000 licenciés FFRandonnée.}
      \swit{O2}{Besoin croissant de résilience}{Zones blanches, crises~: l'État invite à préparer un kit de 72\,h.}
      \swit{O3}{Réseaux locaux à construire}{Refuges, clubs, communes, ADRASEC~: points d'ancrage du maillage.}
      \swit{O4}{Précommandes adaptées}{Le financement participatif teste la demande et paie le premier stock.}
      \swit{O5}{Made in France, souveraineté}{Relocalisation par étapes~: un argument pour les collectivités.}
    \end{column}
    \begin{column}{0.485\textwidth}
      \swhead{Menaces}{insagris}
      \swit[insagris]{T1}{Fait-maison, kits concurrents}{Cartes dès 18\,€, kits prêts à l'emploi 59 à 153\,€. \swref{C}}
      \swit[insagris]{T2}{Substituts puissants}{SOS satellite des smartphones, balises, SMS satellite. \swref{F2}}
      \swit[insagris]{T3}{Contraintes réglementaires}{Marquage CE, cybersécurité, responsabilité produit. \swref{F3, F4}}
      \swit[insagris]{T4}{Dépendance amont}{Un seul fabricant de carte~; cap et marque de Meshtastic. \swref{F5, F6}}
      \swit[insagris]{T5}{Promesse mal comprise}{Ni téléphone, ni balise de détresse~: cadrer la promesse. \swref{F3}}
    \end{column}
  \end{columns}
  \vspace{-0.22cm}
  \begin{pcompact}\centering\footnotesize
    Opportunité réelle, mais fenêtre concurrentielle courte~: il faut bâtir vite une communauté et une marque.
  \end{pcompact}
  \srcnote{Sources~: \cite{ffrando,ffrando_licencies,guide_tous_resp}. F2 à F6 = freins de l'étude documentaire~; C = chapitre Concurrence.}
\end{frame}

%---------------------------------------------------------------------------
\newenvironment{pucesx}
  {\begingroup\setbeamerfont{itemize/enumerate body}{size=\scriptsize}\begin{itemize}\setlength{\itemsep}{0pt}\setlength{\topsep}{1pt}}
  {\end{itemize}\endgroup}
\begin{frame}{SWOT croisé~: quelles stratégies en déduire~?}
  \footnotesize\renewcommand{\arraystretch}{1.0}\setlength{\tabcolsep}{4pt}
  \hyphenpenalty=10000 \exhyphenpenalty=10000
  \noindent\begin{tabularx}{\textwidth}{@{}>{\raggedright\arraybackslash}p{1.65cm}>{\raggedright\arraybackslash}X>{\raggedright\arraybackslash}X@{}}
    & \cellcolor{insarose}\textcolor{insarouge}{\bfseries Opportunités (O)} & \cellcolor{insarose}\textcolor{insarouge}{\bfseries Menaces (T)} \\[2pt]
    \cellcolor{insarose}\textcolor{insarouge}{\bfseries Forces (S)} &
      \cellcolor{insabeige}{\sffamily\bfseries\color{insarouge}Attaquer}
      \begin{pucesx}
        \item Micro-réseaux prêts à l'emploi (Duo, Groupe) pour clubs, refuges, groupes de randonnée.
        \item Précommandes~: tester la demande et financer le premier stock.
        \item Pilotes avec des ADRASEC et des communes~: preuves d'usage, premiers nœuds.
      \end{pucesx} &
      \cellcolor{insabeige}{\sffamily\bfseries\color{insarouge}Défendre}
      \begin{pucesx}
        \item Vendre le prêt-à-l'emploi~: réglages, clé privée, guide, garantie, marquage CE.
        \item Rester un complément du smartphone et de la balise, sans promesse de secours.
        \item Contribuer au code libre, rester « compatible Meshtastic ».
      \end{pucesx} \\[3pt]
    \cellcolor{insarose}\textcolor{insarouge}{\bfseries Faiblesses (W)} &
      \cellcolor{insabeige}{\sffamily\bfseries\color{insarouge}Corriger}
      \begin{pucesx}
        \item Densifier autour de points d'ancrage (refuge, club, mairie) contre l'effet de réseau.
        \item Tester Solo, Duo et Groupe avant la série pour relever la marge.
        \item Bâtir le dossier de financement sur les précommandes.
      \end{pucesx} &
      \cellcolor{insabeige}{\sffamily\bfseries\color{insarouge}Éviter}
      \begin{pucesx}
        \item Notice claire, test à 100\,\%, assurance responsabilité produit.
        \item Double source de cartes (V3 et V4), stock tampon, vente d'abord en France.
        \item Fixer des seuils d'arrêt avant le moule et la série (étude quantitative).
      \end{pucesx} \\
  \end{tabularx}
  \srcnote{Croisement des deux diapositives précédentes (méthode TOWS)~: S, W = forces et faiblesses~; O, T = opportunités et menaces.}
\end{frame}

%---------------------------------------------------------------------------
% Liaisons possibles entre n kits : n(n-1)/2 (2 kits = 1, 4 = 6, 6 = 15, 10 = 45)
\begin{frame}{Enjeu 1~: résoudre l'effet de réseau}
  \begin{columns}[T]
    \begin{column}{0.46\textwidth}
      \centering
      \begin{tikzpicture}[font=\footnotesize]
        \node[pcli,text width=2.3cm,minimum height=0.7cm] (a) at (0,1.05) {Peu d'utilisateurs};
        \node[pgris,text width=1.6cm,minimum height=0.7cm] (b) at (3.55,1.05) {Peu de relais};
        \node[pgris,text width=1.6cm,minimum height=0.7cm] (c) at (1.78,-0.3) {Valeur limitée};
        \draw[fluxr] (a) -- (b);
        \draw[fluxr] (b) -- (c);
        \draw[fluxr] (c) -- (a);
        \node[font=\scriptsize\bfseries,text=insarouge,align=center] at (1.78,0.5) {cercle\\vicieux};
      \end{tikzpicture}

      \vspace{0.05cm}
      {\footnotesize\textbf{Un groupe vaut plus que la somme de ses kits}~: $n$ kits offrent $n(n{-}1)/2$ liaisons.\par}
      \begin{tikzpicture}[font=\scriptsize]
        \tikzset{noeud/.style={circle,fill=insarouge,inner sep=1.3pt}}
        % 2 kits
        \begin{scope}[shift={(0.0,0)}]
          \node[noeud] (u0) at (-0.28,0) {};\node[noeud] (u1) at (0.28,0) {};
          \draw[insagris,thick] (u0) -- (u1);
          \node[align=center,anchor=north] at (0,-0.55) {2 kits\\\textbf{1 liaison}};
        \end{scope}
        % 4 kits
        \begin{scope}[shift={(1.85,0)}]
          \foreach \i in {0,...,3} {\node[noeud] (v\i) at ({45+90*\i}:0.4) {};}
          \foreach \i/\j in {0/1,0/2,0/3,1/2,1/3,2/3} {\draw[insagris,thick] (v\i) -- (v\j);}
          \node[align=center,anchor=north] at (0,-0.55) {4 kits\\\textbf{6 liaisons}};
        \end{scope}
        % 6 kits
        \begin{scope}[shift={(3.7,0)}]
          \foreach \i in {0,...,5} {\node[noeud] (w\i) at ({90+60*\i}:0.42) {};}
          \foreach \i/\j in {0/1,0/2,0/3,0/4,0/5,1/2,1/3,1/4,1/5,2/3,2/4,2/5,3/4,3/5,4/5} {\draw[insagris,thin] (w\i) -- (w\j);}
          \node[align=center,anchor=north] at (0,-0.55) {6 kits\\\textbf{15 liaisons}};
        \end{scope}
        \node[align=center,font=\scriptsize] at (5.15,0) {10 kits\\\textbf{45}};
      \end{tikzpicture}
    \end{column}
    \begin{column}{0.51\textwidth}
      \begin{columns}[T]
        \begin{column}{0.47\linewidth}
          \begin{carte}[height=1.7cm,valign=top]
            \centering
            \begin{tikzpicture}\path (-0.5,-0.35) rectangle (0.5,0.35);\pic[insarouge,scale=0.9] at (0,0) {icone personne};\end{tikzpicture}\\[1pt]
            \gros{Duo, famille}\\[1pt]{\footnotesize 2 à 3 appareils}
          \end{carte}
        \end{column}
        \begin{column}{0.47\linewidth}
          \begin{carte}[height=1.7cm,valign=top]
            \centering
            \begin{tikzpicture}\path (-0.5,-0.35) rectangle (0.5,0.35);\pic[scale=0.9] at (0,0) {icone mont};\end{tikzpicture}\\[1pt]
            \gros{Randonnée}\\[1pt]{\footnotesize 4 appareils}
          \end{carte}
        \end{column}
      \end{columns}
      \vspace{0.1cm}
      \begin{columns}[T]
        \begin{column}{0.47\linewidth}
          \begin{carte}[height=1.7cm,valign=top]
            \centering
            \begin{tikzpicture}\path (-0.5,-0.35) rectangle (0.5,0.35);\pic[insarouge,scale=0.9] at (0,0) {icone tente};\end{tikzpicture}\\[1pt]
            \gros{Refuge, club}\\[1pt]{\footnotesize un relais et ses utilisateurs}
          \end{carte}
        \end{column}
        \begin{column}{0.47\linewidth}
          \begin{carte}[height=1.7cm,valign=top]
            \centering
            \begin{tikzpicture}\path (-0.5,-0.35) rectangle (0.5,0.35);\pic[scale=0.9] at (0,0) {icone croix};\end{tikzpicture}\\[1pt]
            \gros{ONG}\\[1pt]{\footnotesize secours~: 10 à 30 appareils}
          \end{carte}
        \end{column}
      \end{columns}
    \end{column}
  \end{columns}
  \vspace{0.05cm}
  \begin{pcompact}\centering\footnotesize
    Le bon produit est peut-être moins «\,une antenne\,» qu'un réseau prêt à l'emploi pour un groupe constitué.
  \end{pcompact}
  \srcnote{Micro-réseaux d'après le SWOT de l'équipe~; tailles de groupe à mesurer dans le questionnaire (questions 6 et 18).}
\end{frame}

%---------------------------------------------------------------------------
\begin{frame}{Enjeu 2~: un prix d'appel, une marge à protéger}
  \footnotesize\renewcommand{\arraystretch}{1.3}
  \hyphenpenalty=10000 \exhyphenpenalty=10000
  \newcolumntype{N}{>{\raggedleft\arraybackslash}p{1.5cm}}
  \setlength{\tabcolsep}{3pt}
  \noindent\begin{tabular}{@{}>{\raggedright\arraybackslash}p{2.9cm}NNNNNN@{}}
    \rowcolor{insarose}\textcolor{insarouge}{\bfseries Offre} & \textcolor{insarouge}{\bfseries Prix TTC} & \textcolor{insarouge}{\bfseries Par kit} & \textcolor{insarouge}{\bfseries Marge par commande} & \textcolor{insarouge}{\bfseries Marge par kit} & \textcolor{insarouge}{\bfseries Point mort (kits/an)} & \textcolor{insarouge}{\bfseries Résultat année 3}\\
    Kit seul (référence)    & 69\,€  & 69,00\,€ & 20,17\,€ & 20,17\,€ & 21\,300 & +13,7\,k€\\
    Solo                    & 79\,€  & 79,00\,€ & 28,50\,€ & 28,50\,€ & 15\,100 & +197\,k€\\
    \rowcolor{insabeige}Duo (2 kits)           & 139\,€ & 69,50\,€ & 45,17\,€ & 22,59\,€ & 19\,000 & +67\,k€\\
    Groupe (4 kits)         & 249\,€ & 62,25\,€ & 70,18\,€ & 17,54\,€ & 24\,500 & −44\,k€\\
  \end{tabular}
  \vspace{0.12cm}
  \begin{columns}[T]
    \begin{column}{0.325\textwidth}
      \begin{carte}[height=1.7cm,valign=top]
        \raggedright\gros{−73\,k€}\\[1pt]{\footnotesize sur le résultat de l'année 3 si le coût matériel monte de 10\,\% (+3\,€ par kit).}
      \end{carte}
    \end{column}
    \begin{column}{0.325\textwidth}
      \begin{carte}[height=1.7cm,valign=top]
        \raggedright\gros{−44\,k€}\\[1pt]{\footnotesize si la vente, le SAV et la logistique coûtent 2\,€ de plus par commande.}
      \end{carte}
    \end{column}
    \begin{column}{0.325\textwidth}
      \begin{carte}[height=1.7cm,valign=top]
        \raggedright\gros{±92\,k€}\\[1pt]{\footnotesize pour 5\,€ de prix en plus ou en moins sur chaque kit.}
      \end{carte}
    \end{column}
  \end{columns}
  \vspace{0.08cm}
  \begin{pcompact}\centering\footnotesize
    Le Duo relève la marge par kit~; le Groupe à 249\,€ la baisse (il faudrait ≥ 262\,€). Le prix se tranche par le test (questions 19 à 21 du questionnaire), pas par intuition.
  \end{pcompact}
  \srcnote{Calcul avec le modèle de la partie Finances~: matériel 30,30\,€, vente 4\,€ par commande, 1 kit offert pour 10, charges fixes 430\,k€, 22\,000 kits. Marge = prix HT moins coûts variables.}
\end{frame}

%---------------------------------------------------------------------------
\begin{frame}{Enjeu 3~: choisir un marché de départ}
  \begin{columns}[T]
    \begin{column}{0.35\textwidth}
      \begin{carte}[height=4.45cm,valign=top]
        \raggedright
        {\sffamily\bfseries\scriptsize\color{insarouge}CIBLE DE DÉPART (HYPOTHÈSE)}\\[2pt]
        {\sffamily\bfseries\large\color{insarouge}Randonneurs et groupes outdoor}\\[2pt]
        {\footnotesize\textbf{Pourquoi~?}}
        \begin{puces}
          \item besoin simple à comprendre
          \item usage réel en zone blanche
          \item marché large et accessible
          \item vente plus simple que le secours
        \end{puces}
        {\scriptsize\textcolor{insarouge}{À confirmer par l'étude quantitative.}}
      \end{carte}
    \end{column}
    \begin{column}{0.36\textwidth}
      {\sffamily\bfseries\footnotesize Progression naturelle}\\[3pt]
      \begin{tikzpicture}[font=\footnotesize]
        \draw[insarouge!40,thick] (0.2,0.05) -- (0.2,-3.55);
        \node[num] at (0.2,0)     {1};
        \node[num] at (0.2,-0.9) {2};
        \node[num] at (0.2,-1.8)  {3};
        \node[num] at (0.2,-2.7) {4};
        \node[num] at (0.2,-3.6)  {5};
        \node[anchor=west,text width=3.9cm,align=left] at (0.5,0)     {{\bfseries Outdoor}\\[-1pt]{\scriptsize randonneurs et groupes}};
        \node[anchor=west,text width=3.9cm,align=left] at (0.5,-0.9) {{\bfseries Refuges et clubs}\\[-1pt]{\scriptsize points de relais locaux}};
        \node[anchor=west,text width=3.9cm,align=left] at (0.5,-1.8)  {{\bfseries Associations}\\[-1pt]{\scriptsize ADRASEC, communautés structurées}};
        \node[anchor=west,text width=3.9cm,align=left] at (0.5,-2.7) {{\bfseries ONG, communes, pros}\\[-1pt]{\scriptsize déploiements plus larges}};
        \node[anchor=west,text width=3.9cm,align=left] at (0.5,-3.6)  {{\bfseries Grand public}\\[-1pt]{\scriptsize familles, voyageurs, une fois le réseau visible}};
      \end{tikzpicture}
    \end{column}
    \begin{column}{0.25\textwidth}
      \begin{idee}\footnotesize\centering
        {\sffamily\bfseries\color{insarouge}L'objectif}\\[3pt]
        Créer d'abord une zone où le produit est vraiment utile…\\[3pt]
        …puis utiliser cette densité, les retours terrain et la preuve d'usage pour élargir.
      \end{idee}
    \end{column}
  \end{columns}
  \vspace{0.05cm}
  \begin{pcompact}\centering\footnotesize
    Ne pas viser «\,tout le monde\,» au lancement~: gagner un premier territoire d'usage, puis élargir.
  \end{pcompact}
  \srcnote{Choix du marché de départ d'après le SWOT de l'équipe~; les seuils de décision sont dans la partie Étude quantitative.}
\end{frame}

%  Fichier de contenu : ni préambule ni \begin{document}, commence par \section.
%
%  Remplace l'ancien SWOT en PDF (\includepdf supprimé) : le diagnostic est maintenant écrit
%  en LaTeX/TikZ, avec la charte du reste de la présentation, et complété par un SWOT croisé.
%  Ce qui a changé par rapport à la version PDF :
%   - la décomposition des 69 € n'est plus répétée ici (elle est dans la partie Finances) ;
%   - chaque point du diagnostic renvoie à un « frein » F1 à F7 de l'étude documentaire ;
%   - les marges des offres Solo / Duo / Groupe sont recalculées (marge par commande ET par kit,
%     avec le coût de vente de 4 € compté une seule fois par commande) ;
%   - le SWOT croisé (stratégies SO, WO, ST, WT) relie le diagnostic aux trois enjeux.
%  Hypothèses chiffrées : celles de la partie Finances (kit 69 € TTC, matériel 30,30 €, vente 4 €
%  par commande, 1 kit offert pour 10, charges fixes de l'année 3 = 430 k€).
%  Les offres Solo (79 €), Duo (139 €, 2 kits) et Groupe (249 €, 4 kits) sont celles du SWOT d'origine ;
%  le nombre de kits du Groupe (4) est une hypothèse.
%%%%%%%%%%%%%%%%%%%%%%%%%%%%%%%%%%%%%%%%%%%%%%%%%%%%%%%%%%%%%%%%%%%%%%%%%%%%%%
\section{SWOT et enjeux}

% Entrée de SWOT : \swit[couleur du code]{code}{titre}{texte}  (titre et texte sur le même paragraphe)
\newcommand{\swit}[4][insarouge]{%
  \begin{tcolorbox}[colback=insabeige,colframe=insabeige,boxrule=0pt,arc=2pt,left=4pt,right=4pt,
                    top=1.5pt,bottom=1.5pt,before skip=0pt,after skip=2pt]
    \scriptsize\raggedright{\sffamily\bfseries\color{#1}#2}\hspace{0.3em}{\sffamily\bfseries\footnotesize #3}\hspace{0.3em}#4\par
  \end{tcolorbox}}
% En-tête de colonne : \swhead{Forces}{couleur}
\newcommand{\swhead}[2]{%
  {\sffamily\bfseries\footnotesize\color{#2}\MakeUppercase{#1}}\\[-3pt]{\color{#2}\rule{\linewidth}{0.6pt}}\\[-1pt]}
% renvoi vers un enjeu ou un frein
\newcommand{\swref}[1]{\textcolor{insarouge}{\textbf{→\,#1}}}

%---------------------------------------------------------------------------
\begin{frame}{Diagnostic interne~: forces et faiblesses}
  \begin{columns}[T]
    \begin{column}{0.485\textwidth}
      \swhead{Forces}{insarouge}
      \swit{S1}{Proposition de valeur claire}{Une messagerie texte sans réseau mobile, sans compte ni abonnement.}
      \swit{S2}{Coût d'entrée faible}{69\,€ une fois, contre ≈ 830\,€ sur trois ans pour une balise satellite.}
      \swit{S3}{Écosystème libre existant}{Meshtastic, Heltec V3, applications~: l'effort porte sur l'usage.}
      \swit{S4}{Simplicité comme différence}{Réglages Europe, clé privée, QR code, guides en français.}
      \swit{S5}{Un socle, plusieurs usages}{Outdoor, campagne, ONG, crise~: un kit, quatre terrains.}
    \end{column}
    \begin{column}{0.485\textwidth}
      \swhead{Faiblesses}{insagris}
      \swit[insagris]{W1}{Faible barrière technologique}{Un bricoleur refait un kit proche pour 18 à 30\,€. \swref{F5}}
      \swit[insagris]{W2}{Effet de réseau indispensable}{Un kit isolé ne sert à rien~: la valeur vient des voisins. \swref{enjeu 1}}
      \swit[insagris]{W3}{Marge de contribution fine}{20,17\,€ par kit, mais 0,62\,€ de résultat par kit en année 3. \swref{enjeu 2}}
      \swit[insagris]{W4}{Financement avant rentabilité}{≈ 450\,k€ à financer, point mort ≈ 21\,300 kits par an. \swref{F7}}
      \swit[insagris]{W5}{Positionnement trop large}{ONG, randonneurs, makers, public~: attentes différentes. \swref{enjeu 3}}
    \end{column}
  \end{columns}
  \vspace{-0.2cm}
  \begin{pcompact}\centering\footnotesize
    Lecture clé~: la technologie est accessible~; la valeur viendra de l'usage, de la marque et de l'écosystème.
  \end{pcompact}
  \srcnote{Synthèse des parties Concurrence, Finances et Étude documentaire. F2 à F7 = freins de l'étude documentaire.}
\end{frame}

%---------------------------------------------------------------------------
\begin{frame}{Diagnostic externe~: opportunités et menaces}
  \begin{columns}[T]
    \begin{column}{0.485\textwidth}
      \swhead{Opportunités}{insarouge}
      \swit{O1}{Un marché outdoor important}{27 millions de pratiquants, dont 240\,000 licenciés FFRandonnée.}
      \swit{O2}{Besoin croissant de résilience}{Zones blanches, crises~: l'État invite à préparer un kit de 72\,h.}
      \swit{O3}{Réseaux locaux à construire}{Refuges, clubs, communes, ADRASEC~: points d'ancrage du maillage.}
      \swit{O4}{Précommandes adaptées}{Le financement participatif teste la demande et paie le premier stock.}
      \swit{O5}{Made in France, souveraineté}{Relocalisation par étapes~: un argument pour les collectivités.}
    \end{column}
    \begin{column}{0.485\textwidth}
      \swhead{Menaces}{insagris}
      \swit[insagris]{T1}{Fait-maison, kits concurrents}{Cartes dès 18\,€, kits prêts à l'emploi 59 à 153\,€. \swref{C}}
      \swit[insagris]{T2}{Substituts puissants}{SOS satellite des smartphones, balises, SMS satellite. \swref{F2}}
      \swit[insagris]{T3}{Contraintes réglementaires}{Marquage CE, cybersécurité, responsabilité produit. \swref{F3, F4}}
      \swit[insagris]{T4}{Dépendance amont}{Un seul fabricant de carte~; cap et marque de Meshtastic. \swref{F5, F6}}
      \swit[insagris]{T5}{Promesse mal comprise}{Ni téléphone, ni balise de détresse~: cadrer la promesse. \swref{F3}}
    \end{column}
  \end{columns}
  \vspace{-0.22cm}
  \begin{pcompact}\centering\footnotesize
    Opportunité réelle, mais fenêtre concurrentielle courte~: il faut bâtir vite une communauté et une marque.
  \end{pcompact}
  \srcnote{Sources~: \cite{ffrando,ffrando_licencies,guide_tous_resp}. F2 à F6 = freins de l'étude documentaire~; C = chapitre Concurrence.}
\end{frame}

%---------------------------------------------------------------------------
\newenvironment{pucesx}
  {\begingroup\setbeamerfont{itemize/enumerate body}{size=\scriptsize}\begin{itemize}\setlength{\itemsep}{0pt}\setlength{\topsep}{1pt}}
  {\end{itemize}\endgroup}
\begin{frame}{SWOT croisé~: quelles stratégies en déduire~?}
  \footnotesize\renewcommand{\arraystretch}{1.0}\setlength{\tabcolsep}{4pt}
  \hyphenpenalty=10000 \exhyphenpenalty=10000
  \noindent\begin{tabularx}{\textwidth}{@{}>{\raggedright\arraybackslash}p{1.65cm}>{\raggedright\arraybackslash}X>{\raggedright\arraybackslash}X@{}}
    & \cellcolor{insarose}\textcolor{insarouge}{\bfseries Opportunités (O)} & \cellcolor{insarose}\textcolor{insarouge}{\bfseries Menaces (T)} \\[2pt]
    \cellcolor{insarose}\textcolor{insarouge}{\bfseries Forces (S)} &
      \cellcolor{insabeige}{\sffamily\bfseries\color{insarouge}Attaquer}
      \begin{pucesx}
        \item Micro-réseaux prêts à l'emploi (Duo, Groupe) pour clubs, refuges, groupes de randonnée.
        \item Précommandes~: tester la demande et financer le premier stock.
        \item Pilotes avec des ADRASEC et des communes~: preuves d'usage, premiers nœuds.
      \end{pucesx} &
      \cellcolor{insabeige}{\sffamily\bfseries\color{insarouge}Défendre}
      \begin{pucesx}
        \item Vendre le prêt-à-l'emploi~: réglages, clé privée, guide, garantie, marquage CE.
        \item Rester un complément du smartphone et de la balise, sans promesse de secours.
        \item Contribuer au code libre, rester « compatible Meshtastic ».
      \end{pucesx} \\[3pt]
    \cellcolor{insarose}\textcolor{insarouge}{\bfseries Faiblesses (W)} &
      \cellcolor{insabeige}{\sffamily\bfseries\color{insarouge}Corriger}
      \begin{pucesx}
        \item Densifier autour de points d'ancrage (refuge, club, mairie) contre l'effet de réseau.
        \item Tester Solo, Duo et Groupe avant la série pour relever la marge.
        \item Bâtir le dossier de financement sur les précommandes.
      \end{pucesx} &
      \cellcolor{insabeige}{\sffamily\bfseries\color{insarouge}Éviter}
      \begin{pucesx}
        \item Notice claire, test à 100\,\%, assurance responsabilité produit.
        \item Double source de cartes (V3 et V4), stock tampon, vente d'abord en France.
        \item Fixer des seuils d'arrêt avant le moule et la série (étude quantitative).
      \end{pucesx} \\
  \end{tabularx}
  \srcnote{Croisement des deux diapositives précédentes (méthode TOWS)~: S, W = forces et faiblesses~; O, T = opportunités et menaces.}
\end{frame}

%---------------------------------------------------------------------------
% Liaisons possibles entre n kits : n(n-1)/2 (2 kits = 1, 4 = 6, 6 = 15, 10 = 45)
\begin{frame}{Enjeu 1~: résoudre l'effet de réseau}
  \begin{columns}[T]
    \begin{column}{0.46\textwidth}
      \centering
      \begin{tikzpicture}[font=\footnotesize]
        \node[pcli,text width=2.3cm,minimum height=0.7cm] (a) at (0,1.05) {Peu d'utilisateurs};
        \node[pgris,text width=1.6cm,minimum height=0.7cm] (b) at (3.55,1.05) {Peu de relais};
        \node[pgris,text width=1.6cm,minimum height=0.7cm] (c) at (1.78,-0.3) {Valeur limitée};
        \draw[fluxr] (a) -- (b);
        \draw[fluxr] (b) -- (c);
        \draw[fluxr] (c) -- (a);
        \node[font=\scriptsize\bfseries,text=insarouge,align=center] at (1.78,0.5) {cercle\\vicieux};
      \end{tikzpicture}

      \vspace{0.05cm}
      {\footnotesize\textbf{Un groupe vaut plus que la somme de ses kits}~: $n$ kits offrent $n(n{-}1)/2$ liaisons.\par}
      \begin{tikzpicture}[font=\scriptsize]
        \tikzset{noeud/.style={circle,fill=insarouge,inner sep=1.3pt}}
        % 2 kits
        \begin{scope}[shift={(0.0,0)}]
          \node[noeud] (u0) at (-0.28,0) {};\node[noeud] (u1) at (0.28,0) {};
          \draw[insagris,thick] (u0) -- (u1);
          \node[align=center,anchor=north] at (0,-0.55) {2 kits\\\textbf{1 liaison}};
        \end{scope}
        % 4 kits
        \begin{scope}[shift={(1.85,0)}]
          \foreach \i in {0,...,3} {\node[noeud] (v\i) at ({45+90*\i}:0.4) {};}
          \foreach \i/\j in {0/1,0/2,0/3,1/2,1/3,2/3} {\draw[insagris,thick] (v\i) -- (v\j);}
          \node[align=center,anchor=north] at (0,-0.55) {4 kits\\\textbf{6 liaisons}};
        \end{scope}
        % 6 kits
        \begin{scope}[shift={(3.7,0)}]
          \foreach \i in {0,...,5} {\node[noeud] (w\i) at ({90+60*\i}:0.42) {};}
          \foreach \i/\j in {0/1,0/2,0/3,0/4,0/5,1/2,1/3,1/4,1/5,2/3,2/4,2/5,3/4,3/5,4/5} {\draw[insagris,thin] (w\i) -- (w\j);}
          \node[align=center,anchor=north] at (0,-0.55) {6 kits\\\textbf{15 liaisons}};
        \end{scope}
        \node[align=center,font=\scriptsize] at (5.15,0) {10 kits\\\textbf{45}};
      \end{tikzpicture}
    \end{column}
    \begin{column}{0.51\textwidth}
      \begin{columns}[T]
        \begin{column}{0.47\linewidth}
          \begin{carte}[height=1.7cm,valign=top]
            \centering
            \begin{tikzpicture}\path (-0.5,-0.35) rectangle (0.5,0.35);\pic[insarouge,scale=0.9] at (0,0) {icone personne};\end{tikzpicture}\\[1pt]
            \gros{Duo, famille}\\[1pt]{\footnotesize 2 à 3 appareils}
          \end{carte}
        \end{column}
        \begin{column}{0.47\linewidth}
          \begin{carte}[height=1.7cm,valign=top]
            \centering
            \begin{tikzpicture}\path (-0.5,-0.35) rectangle (0.5,0.35);\pic[scale=0.9] at (0,0) {icone mont};\end{tikzpicture}\\[1pt]
            \gros{Randonnée}\\[1pt]{\footnotesize 4 appareils}
          \end{carte}
        \end{column}
      \end{columns}
      \vspace{0.1cm}
      \begin{columns}[T]
        \begin{column}{0.47\linewidth}
          \begin{carte}[height=1.7cm,valign=top]
            \centering
            \begin{tikzpicture}\path (-0.5,-0.35) rectangle (0.5,0.35);\pic[insarouge,scale=0.9] at (0,0) {icone tente};\end{tikzpicture}\\[1pt]
            \gros{Refuge, club}\\[1pt]{\footnotesize un relais et ses utilisateurs}
          \end{carte}
        \end{column}
        \begin{column}{0.47\linewidth}
          \begin{carte}[height=1.7cm,valign=top]
            \centering
            \begin{tikzpicture}\path (-0.5,-0.35) rectangle (0.5,0.35);\pic[scale=0.9] at (0,0) {icone croix};\end{tikzpicture}\\[1pt]
            \gros{ONG}\\[1pt]{\footnotesize secours~: 10 à 30 appareils}
          \end{carte}
        \end{column}
      \end{columns}
    \end{column}
  \end{columns}
  \vspace{0.05cm}
  \begin{pcompact}\centering\footnotesize
    Le bon produit est peut-être moins «\,une antenne\,» qu'un réseau prêt à l'emploi pour un groupe constitué.
  \end{pcompact}
  \srcnote{Micro-réseaux d'après le SWOT de l'équipe~; tailles de groupe à mesurer dans le questionnaire (questions 6 et 18).}
\end{frame}

%---------------------------------------------------------------------------
\begin{frame}{Enjeu 2~: un prix d'appel, une marge à protéger}
  \footnotesize\renewcommand{\arraystretch}{1.3}
  \hyphenpenalty=10000 \exhyphenpenalty=10000
  \newcolumntype{N}{>{\raggedleft\arraybackslash}p{1.5cm}}
  \setlength{\tabcolsep}{3pt}
  \noindent\begin{tabular}{@{}>{\raggedright\arraybackslash}p{2.9cm}NNNNNN@{}}
    \rowcolor{insarose}\textcolor{insarouge}{\bfseries Offre} & \textcolor{insarouge}{\bfseries Prix TTC} & \textcolor{insarouge}{\bfseries Par kit} & \textcolor{insarouge}{\bfseries Marge par commande} & \textcolor{insarouge}{\bfseries Marge par kit} & \textcolor{insarouge}{\bfseries Point mort (kits/an)} & \textcolor{insarouge}{\bfseries Résultat année 3}\\
    Kit seul (référence)    & 69\,€  & 69,00\,€ & 20,17\,€ & 20,17\,€ & 21\,300 & +13,7\,k€\\
    Solo                    & 79\,€  & 79,00\,€ & 28,50\,€ & 28,50\,€ & 15\,100 & +197\,k€\\
    \rowcolor{insabeige}Duo (2 kits)           & 139\,€ & 69,50\,€ & 45,17\,€ & 22,59\,€ & 19\,000 & +67\,k€\\
    Groupe (4 kits)         & 249\,€ & 62,25\,€ & 70,18\,€ & 17,54\,€ & 24\,500 & −44\,k€\\
  \end{tabular}
  \vspace{0.12cm}
  \begin{columns}[T]
    \begin{column}{0.325\textwidth}
      \begin{carte}[height=1.7cm,valign=top]
        \raggedright\gros{−73\,k€}\\[1pt]{\footnotesize sur le résultat de l'année 3 si le coût matériel monte de 10\,\% (+3\,€ par kit).}
      \end{carte}
    \end{column}
    \begin{column}{0.325\textwidth}
      \begin{carte}[height=1.7cm,valign=top]
        \raggedright\gros{−44\,k€}\\[1pt]{\footnotesize si la vente, le SAV et la logistique coûtent 2\,€ de plus par commande.}
      \end{carte}
    \end{column}
    \begin{column}{0.325\textwidth}
      \begin{carte}[height=1.7cm,valign=top]
        \raggedright\gros{±92\,k€}\\[1pt]{\footnotesize pour 5\,€ de prix en plus ou en moins sur chaque kit.}
      \end{carte}
    \end{column}
  \end{columns}
  \vspace{0.08cm}
  \begin{pcompact}\centering\footnotesize
    Le Duo relève la marge par kit~; le Groupe à 249\,€ la baisse (il faudrait ≥ 262\,€). Le prix se tranche par le test (questions 19 à 21 du questionnaire), pas par intuition.
  \end{pcompact}
  \srcnote{Calcul avec le modèle de la partie Finances~: matériel 30,30\,€, vente 4\,€ par commande, 1 kit offert pour 10, charges fixes 430\,k€, 22\,000 kits. Marge = prix HT moins coûts variables.}
\end{frame}

%---------------------------------------------------------------------------
\begin{frame}{Enjeu 3~: choisir un marché de départ}
  \begin{columns}[T]
    \begin{column}{0.35\textwidth}
      \begin{carte}[height=4.45cm,valign=top]
        \raggedright
        {\sffamily\bfseries\scriptsize\color{insarouge}CIBLE DE DÉPART (HYPOTHÈSE)}\\[2pt]
        {\sffamily\bfseries\large\color{insarouge}Randonneurs et groupes outdoor}\\[2pt]
        {\footnotesize\textbf{Pourquoi~?}}
        \begin{puces}
          \item besoin simple à comprendre
          \item usage réel en zone blanche
          \item marché large et accessible
          \item vente plus simple que le secours
        \end{puces}
        {\scriptsize\textcolor{insarouge}{À confirmer par l'étude quantitative.}}
      \end{carte}
    \end{column}
    \begin{column}{0.36\textwidth}
      {\sffamily\bfseries\footnotesize Progression naturelle}\\[3pt]
      \begin{tikzpicture}[font=\footnotesize]
        \draw[insarouge!40,thick] (0.2,0.05) -- (0.2,-3.55);
        \node[num] at (0.2,0)     {1};
        \node[num] at (0.2,-0.9) {2};
        \node[num] at (0.2,-1.8)  {3};
        \node[num] at (0.2,-2.7) {4};
        \node[num] at (0.2,-3.6)  {5};
        \node[anchor=west,text width=3.9cm,align=left] at (0.5,0)     {{\bfseries Outdoor}\\[-1pt]{\scriptsize randonneurs et groupes}};
        \node[anchor=west,text width=3.9cm,align=left] at (0.5,-0.9) {{\bfseries Refuges et clubs}\\[-1pt]{\scriptsize points de relais locaux}};
        \node[anchor=west,text width=3.9cm,align=left] at (0.5,-1.8)  {{\bfseries Associations}\\[-1pt]{\scriptsize ADRASEC, communautés structurées}};
        \node[anchor=west,text width=3.9cm,align=left] at (0.5,-2.7) {{\bfseries ONG, communes, pros}\\[-1pt]{\scriptsize déploiements plus larges}};
        \node[anchor=west,text width=3.9cm,align=left] at (0.5,-3.6)  {{\bfseries Grand public}\\[-1pt]{\scriptsize familles, voyageurs, une fois le réseau visible}};
      \end{tikzpicture}
    \end{column}
    \begin{column}{0.25\textwidth}
      \begin{idee}\footnotesize\centering
        {\sffamily\bfseries\color{insarouge}L'objectif}\\[3pt]
        Créer d'abord une zone où le produit est vraiment utile…\\[3pt]
        …puis utiliser cette densité, les retours terrain et la preuve d'usage pour élargir.
      \end{idee}
    \end{column}
  \end{columns}
  \vspace{0.05cm}
  \begin{pcompact}\centering\footnotesize
    Ne pas viser «\,tout le monde\,» au lancement~: gagner un premier territoire d'usage, puis élargir.
  \end{pcompact}
  \srcnote{Choix du marché de départ d'après le SWOT de l'équipe~; les seuils de décision sont dans la partie Étude quantitative.}
\end{frame}

%  Fichier de contenu : ni préambule ni \begin{document}, commence par \section.
%
%  Remplace l'ancien SWOT en PDF (\includepdf supprimé) : le diagnostic est maintenant écrit
%  en LaTeX/TikZ, avec la charte du reste de la présentation, et complété par un SWOT croisé.
%  Ce qui a changé par rapport à la version PDF :
%   - la décomposition des 69 € n'est plus répétée ici (elle est dans la partie Finances) ;
%   - chaque point du diagnostic renvoie à un « frein » F1 à F7 de l'étude documentaire ;
%   - les marges des offres Solo / Duo / Groupe sont recalculées (marge par commande ET par kit,
%     avec le coût de vente de 4 € compté une seule fois par commande) ;
%   - le SWOT croisé (stratégies SO, WO, ST, WT) relie le diagnostic aux trois enjeux.
%  Hypothèses chiffrées : celles de la partie Finances (kit 69 € TTC, matériel 30,30 €, vente 4 €
%  par commande, 1 kit offert pour 10, charges fixes de l'année 3 = 430 k€).
%  Les offres Solo (79 €), Duo (139 €, 2 kits) et Groupe (249 €, 4 kits) sont celles du SWOT d'origine ;
%  le nombre de kits du Groupe (4) est une hypothèse.
%%%%%%%%%%%%%%%%%%%%%%%%%%%%%%%%%%%%%%%%%%%%%%%%%%%%%%%%%%%%%%%%%%%%%%%%%%%%%%
\section{SWOT et enjeux}

% Entrée de SWOT : \swit[couleur du code]{code}{titre}{texte}  (titre et texte sur le même paragraphe)
\newcommand{\swit}[4][insarouge]{%
  \begin{tcolorbox}[colback=insabeige,colframe=insabeige,boxrule=0pt,arc=2pt,left=4pt,right=4pt,
                    top=1.5pt,bottom=1.5pt,before skip=0pt,after skip=2pt]
    \scriptsize\raggedright{\sffamily\bfseries\color{#1}#2}\hspace{0.3em}{\sffamily\bfseries\footnotesize #3}\hspace{0.3em}#4\par
  \end{tcolorbox}}
% En-tête de colonne : \swhead{Forces}{couleur}
\newcommand{\swhead}[2]{%
  {\sffamily\bfseries\footnotesize\color{#2}\MakeUppercase{#1}}\\[-3pt]{\color{#2}\rule{\linewidth}{0.6pt}}\\[-1pt]}
% renvoi vers un enjeu ou un frein
\newcommand{\swref}[1]{\textcolor{insarouge}{\textbf{→\,#1}}}

%---------------------------------------------------------------------------
\begin{frame}{Diagnostic interne~: forces et faiblesses}
  \begin{columns}[T]
    \begin{column}{0.485\textwidth}
      \swhead{Forces}{insarouge}
      \swit{S1}{Proposition de valeur claire}{Une messagerie texte sans réseau mobile, sans compte ni abonnement.}
      \swit{S2}{Coût d'entrée faible}{69\,€ une fois, contre ≈ 830\,€ sur trois ans pour une balise satellite.}
      \swit{S3}{Écosystème libre existant}{Meshtastic, Heltec V3, applications~: l'effort porte sur l'usage.}
      \swit{S4}{Simplicité comme différence}{Réglages Europe, clé privée, QR code, guides en français.}
      \swit{S5}{Un socle, plusieurs usages}{Outdoor, campagne, ONG, crise~: un kit, quatre terrains.}
    \end{column}
    \begin{column}{0.485\textwidth}
      \swhead{Faiblesses}{insagris}
      \swit[insagris]{W1}{Faible barrière technologique}{Un bricoleur refait un kit proche pour 18 à 30\,€. \swref{F5}}
      \swit[insagris]{W2}{Effet de réseau indispensable}{Un kit isolé ne sert à rien~: la valeur vient des voisins. \swref{enjeu 1}}
      \swit[insagris]{W3}{Marge de contribution fine}{20,17\,€ par kit, mais 0,62\,€ de résultat par kit en année 3. \swref{enjeu 2}}
      \swit[insagris]{W4}{Financement avant rentabilité}{≈ 450\,k€ à financer, point mort ≈ 21\,300 kits par an. \swref{F7}}
      \swit[insagris]{W5}{Positionnement trop large}{ONG, randonneurs, makers, public~: attentes différentes. \swref{enjeu 3}}
    \end{column}
  \end{columns}
  \vspace{-0.2cm}
  \begin{pcompact}\centering\footnotesize
    Lecture clé~: la technologie est accessible~; la valeur viendra de l'usage, de la marque et de l'écosystème.
  \end{pcompact}
  \srcnote{Synthèse des parties Concurrence, Finances et Étude documentaire. F2 à F7 = freins de l'étude documentaire.}
\end{frame}

%---------------------------------------------------------------------------
\begin{frame}{Diagnostic externe~: opportunités et menaces}
  \begin{columns}[T]
    \begin{column}{0.485\textwidth}
      \swhead{Opportunités}{insarouge}
      \swit{O1}{Un marché outdoor important}{27 millions de pratiquants, dont 240\,000 licenciés FFRandonnée.}
      \swit{O2}{Besoin croissant de résilience}{Zones blanches, crises~: l'État invite à préparer un kit de 72\,h.}
      \swit{O3}{Réseaux locaux à construire}{Refuges, clubs, communes, ADRASEC~: points d'ancrage du maillage.}
      \swit{O4}{Précommandes adaptées}{Le financement participatif teste la demande et paie le premier stock.}
      \swit{O5}{Made in France, souveraineté}{Relocalisation par étapes~: un argument pour les collectivités.}
    \end{column}
    \begin{column}{0.485\textwidth}
      \swhead{Menaces}{insagris}
      \swit[insagris]{T1}{Fait-maison, kits concurrents}{Cartes dès 18\,€, kits prêts à l'emploi 59 à 153\,€. \swref{C}}
      \swit[insagris]{T2}{Substituts puissants}{SOS satellite des smartphones, balises, SMS satellite. \swref{F2}}
      \swit[insagris]{T3}{Contraintes réglementaires}{Marquage CE, cybersécurité, responsabilité produit. \swref{F3, F4}}
      \swit[insagris]{T4}{Dépendance amont}{Un seul fabricant de carte~; cap et marque de Meshtastic. \swref{F5, F6}}
      \swit[insagris]{T5}{Promesse mal comprise}{Ni téléphone, ni balise de détresse~: cadrer la promesse. \swref{F3}}
    \end{column}
  \end{columns}
  \vspace{-0.22cm}
  \begin{pcompact}\centering\footnotesize
    Opportunité réelle, mais fenêtre concurrentielle courte~: il faut bâtir vite une communauté et une marque.
  \end{pcompact}
  \srcnote{Sources~: \cite{ffrando,ffrando_licencies,guide_tous_resp}. F2 à F6 = freins de l'étude documentaire~; C = chapitre Concurrence.}
\end{frame}

%---------------------------------------------------------------------------
\newenvironment{pucesx}
  {\begingroup\setbeamerfont{itemize/enumerate body}{size=\scriptsize}\begin{itemize}\setlength{\itemsep}{0pt}\setlength{\topsep}{1pt}}
  {\end{itemize}\endgroup}
\begin{frame}{SWOT croisé~: quelles stratégies en déduire~?}
  \footnotesize\renewcommand{\arraystretch}{1.0}\setlength{\tabcolsep}{4pt}
  \hyphenpenalty=10000 \exhyphenpenalty=10000
  \noindent\begin{tabularx}{\textwidth}{@{}>{\raggedright\arraybackslash}p{1.65cm}>{\raggedright\arraybackslash}X>{\raggedright\arraybackslash}X@{}}
    & \cellcolor{insarose}\textcolor{insarouge}{\bfseries Opportunités (O)} & \cellcolor{insarose}\textcolor{insarouge}{\bfseries Menaces (T)} \\[2pt]
    \cellcolor{insarose}\textcolor{insarouge}{\bfseries Forces (S)} &
      \cellcolor{insabeige}{\sffamily\bfseries\color{insarouge}Attaquer}
      \begin{pucesx}
        \item Micro-réseaux prêts à l'emploi (Duo, Groupe) pour clubs, refuges, groupes de randonnée.
        \item Précommandes~: tester la demande et financer le premier stock.
        \item Pilotes avec des ADRASEC et des communes~: preuves d'usage, premiers nœuds.
      \end{pucesx} &
      \cellcolor{insabeige}{\sffamily\bfseries\color{insarouge}Défendre}
      \begin{pucesx}
        \item Vendre le prêt-à-l'emploi~: réglages, clé privée, guide, garantie, marquage CE.
        \item Rester un complément du smartphone et de la balise, sans promesse de secours.
        \item Contribuer au code libre, rester « compatible Meshtastic ».
      \end{pucesx} \\[3pt]
    \cellcolor{insarose}\textcolor{insarouge}{\bfseries Faiblesses (W)} &
      \cellcolor{insabeige}{\sffamily\bfseries\color{insarouge}Corriger}
      \begin{pucesx}
        \item Densifier autour de points d'ancrage (refuge, club, mairie) contre l'effet de réseau.
        \item Tester Solo, Duo et Groupe avant la série pour relever la marge.
        \item Bâtir le dossier de financement sur les précommandes.
      \end{pucesx} &
      \cellcolor{insabeige}{\sffamily\bfseries\color{insarouge}Éviter}
      \begin{pucesx}
        \item Notice claire, test à 100\,\%, assurance responsabilité produit.
        \item Double source de cartes (V3 et V4), stock tampon, vente d'abord en France.
        \item Fixer des seuils d'arrêt avant le moule et la série (étude quantitative).
      \end{pucesx} \\
  \end{tabularx}
  \srcnote{Croisement des deux diapositives précédentes (méthode TOWS)~: S, W = forces et faiblesses~; O, T = opportunités et menaces.}
\end{frame}

%---------------------------------------------------------------------------
% Liaisons possibles entre n kits : n(n-1)/2 (2 kits = 1, 4 = 6, 6 = 15, 10 = 45)
\begin{frame}{Enjeu 1~: résoudre l'effet de réseau}
  \begin{columns}[T]
    \begin{column}{0.46\textwidth}
      \centering
      \begin{tikzpicture}[font=\footnotesize]
        \node[pcli,text width=2.3cm,minimum height=0.7cm] (a) at (0,1.05) {Peu d'utilisateurs};
        \node[pgris,text width=1.6cm,minimum height=0.7cm] (b) at (3.55,1.05) {Peu de relais};
        \node[pgris,text width=1.6cm,minimum height=0.7cm] (c) at (1.78,-0.3) {Valeur limitée};
        \draw[fluxr] (a) -- (b);
        \draw[fluxr] (b) -- (c);
        \draw[fluxr] (c) -- (a);
        \node[font=\scriptsize\bfseries,text=insarouge,align=center] at (1.78,0.5) {cercle\\vicieux};
      \end{tikzpicture}

      \vspace{0.05cm}
      {\footnotesize\textbf{Un groupe vaut plus que la somme de ses kits}~: $n$ kits offrent $n(n{-}1)/2$ liaisons.\par}
      \begin{tikzpicture}[font=\scriptsize]
        \tikzset{noeud/.style={circle,fill=insarouge,inner sep=1.3pt}}
        % 2 kits
        \begin{scope}[shift={(0.0,0)}]
          \node[noeud] (u0) at (-0.28,0) {};\node[noeud] (u1) at (0.28,0) {};
          \draw[insagris,thick] (u0) -- (u1);
          \node[align=center,anchor=north] at (0,-0.55) {2 kits\\\textbf{1 liaison}};
        \end{scope}
        % 4 kits
        \begin{scope}[shift={(1.85,0)}]
          \foreach \i in {0,...,3} {\node[noeud] (v\i) at ({45+90*\i}:0.4) {};}
          \foreach \i/\j in {0/1,0/2,0/3,1/2,1/3,2/3} {\draw[insagris,thick] (v\i) -- (v\j);}
          \node[align=center,anchor=north] at (0,-0.55) {4 kits\\\textbf{6 liaisons}};
        \end{scope}
        % 6 kits
        \begin{scope}[shift={(3.7,0)}]
          \foreach \i in {0,...,5} {\node[noeud] (w\i) at ({90+60*\i}:0.42) {};}
          \foreach \i/\j in {0/1,0/2,0/3,0/4,0/5,1/2,1/3,1/4,1/5,2/3,2/4,2/5,3/4,3/5,4/5} {\draw[insagris,thin] (w\i) -- (w\j);}
          \node[align=center,anchor=north] at (0,-0.55) {6 kits\\\textbf{15 liaisons}};
        \end{scope}
        \node[align=center,font=\scriptsize] at (5.15,0) {10 kits\\\textbf{45}};
      \end{tikzpicture}
    \end{column}
    \begin{column}{0.51\textwidth}
      \begin{columns}[T]
        \begin{column}{0.47\linewidth}
          \begin{carte}[height=1.7cm,valign=top]
            \centering
            \begin{tikzpicture}\path (-0.5,-0.35) rectangle (0.5,0.35);\pic[insarouge,scale=0.9] at (0,0) {icone personne};\end{tikzpicture}\\[1pt]
            \gros{Duo, famille}\\[1pt]{\footnotesize 2 à 3 appareils}
          \end{carte}
        \end{column}
        \begin{column}{0.47\linewidth}
          \begin{carte}[height=1.7cm,valign=top]
            \centering
            \begin{tikzpicture}\path (-0.5,-0.35) rectangle (0.5,0.35);\pic[scale=0.9] at (0,0) {icone mont};\end{tikzpicture}\\[1pt]
            \gros{Randonnée}\\[1pt]{\footnotesize 4 appareils}
          \end{carte}
        \end{column}
      \end{columns}
      \vspace{0.1cm}
      \begin{columns}[T]
        \begin{column}{0.47\linewidth}
          \begin{carte}[height=1.7cm,valign=top]
            \centering
            \begin{tikzpicture}\path (-0.5,-0.35) rectangle (0.5,0.35);\pic[insarouge,scale=0.9] at (0,0) {icone tente};\end{tikzpicture}\\[1pt]
            \gros{Refuge, club}\\[1pt]{\footnotesize un relais et ses utilisateurs}
          \end{carte}
        \end{column}
        \begin{column}{0.47\linewidth}
          \begin{carte}[height=1.7cm,valign=top]
            \centering
            \begin{tikzpicture}\path (-0.5,-0.35) rectangle (0.5,0.35);\pic[scale=0.9] at (0,0) {icone croix};\end{tikzpicture}\\[1pt]
            \gros{ONG}\\[1pt]{\footnotesize secours~: 10 à 30 appareils}
          \end{carte}
        \end{column}
      \end{columns}
    \end{column}
  \end{columns}
  \vspace{0.05cm}
  \begin{pcompact}\centering\footnotesize
    Le bon produit est peut-être moins «\,une antenne\,» qu'un réseau prêt à l'emploi pour un groupe constitué.
  \end{pcompact}
  \srcnote{Micro-réseaux d'après le SWOT de l'équipe~; tailles de groupe à mesurer dans le questionnaire (questions 6 et 18).}
\end{frame}

%---------------------------------------------------------------------------
\begin{frame}{Enjeu 2~: un prix d'appel, une marge à protéger}
  \footnotesize\renewcommand{\arraystretch}{1.3}
  \hyphenpenalty=10000 \exhyphenpenalty=10000
  \newcolumntype{N}{>{\raggedleft\arraybackslash}p{1.5cm}}
  \setlength{\tabcolsep}{3pt}
  \noindent\begin{tabular}{@{}>{\raggedright\arraybackslash}p{2.9cm}NNNNNN@{}}
    \rowcolor{insarose}\textcolor{insarouge}{\bfseries Offre} & \textcolor{insarouge}{\bfseries Prix TTC} & \textcolor{insarouge}{\bfseries Par kit} & \textcolor{insarouge}{\bfseries Marge par commande} & \textcolor{insarouge}{\bfseries Marge par kit} & \textcolor{insarouge}{\bfseries Point mort (kits/an)} & \textcolor{insarouge}{\bfseries Résultat année 3}\\
    Kit seul (référence)    & 69\,€  & 69,00\,€ & 20,17\,€ & 20,17\,€ & 21\,300 & +13,7\,k€\\
    Solo                    & 79\,€  & 79,00\,€ & 28,50\,€ & 28,50\,€ & 15\,100 & +197\,k€\\
    \rowcolor{insabeige}Duo (2 kits)           & 139\,€ & 69,50\,€ & 45,17\,€ & 22,59\,€ & 19\,000 & +67\,k€\\
    Groupe (4 kits)         & 249\,€ & 62,25\,€ & 70,18\,€ & 17,54\,€ & 24\,500 & −44\,k€\\
  \end{tabular}
  \vspace{0.12cm}
  \begin{columns}[T]
    \begin{column}{0.325\textwidth}
      \begin{carte}[height=1.7cm,valign=top]
        \raggedright\gros{−73\,k€}\\[1pt]{\footnotesize sur le résultat de l'année 3 si le coût matériel monte de 10\,\% (+3\,€ par kit).}
      \end{carte}
    \end{column}
    \begin{column}{0.325\textwidth}
      \begin{carte}[height=1.7cm,valign=top]
        \raggedright\gros{−44\,k€}\\[1pt]{\footnotesize si la vente, le SAV et la logistique coûtent 2\,€ de plus par commande.}
      \end{carte}
    \end{column}
    \begin{column}{0.325\textwidth}
      \begin{carte}[height=1.7cm,valign=top]
        \raggedright\gros{±92\,k€}\\[1pt]{\footnotesize pour 5\,€ de prix en plus ou en moins sur chaque kit.}
      \end{carte}
    \end{column}
  \end{columns}
  \vspace{0.08cm}
  \begin{pcompact}\centering\footnotesize
    Le Duo relève la marge par kit~; le Groupe à 249\,€ la baisse (il faudrait ≥ 262\,€). Le prix se tranche par le test (questions 19 à 21 du questionnaire), pas par intuition.
  \end{pcompact}
  \srcnote{Calcul avec le modèle de la partie Finances~: matériel 30,30\,€, vente 4\,€ par commande, 1 kit offert pour 10, charges fixes 430\,k€, 22\,000 kits. Marge = prix HT moins coûts variables.}
\end{frame}

%---------------------------------------------------------------------------
\begin{frame}{Enjeu 3~: choisir un marché de départ}
  \begin{columns}[T]
    \begin{column}{0.35\textwidth}
      \begin{carte}[height=4.45cm,valign=top]
        \raggedright
        {\sffamily\bfseries\scriptsize\color{insarouge}CIBLE DE DÉPART (HYPOTHÈSE)}\\[2pt]
        {\sffamily\bfseries\large\color{insarouge}Randonneurs et groupes outdoor}\\[2pt]
        {\footnotesize\textbf{Pourquoi~?}}
        \begin{puces}
          \item besoin simple à comprendre
          \item usage réel en zone blanche
          \item marché large et accessible
          \item vente plus simple que le secours
        \end{puces}
        {\scriptsize\textcolor{insarouge}{À confirmer par l'étude quantitative.}}
      \end{carte}
    \end{column}
    \begin{column}{0.36\textwidth}
      {\sffamily\bfseries\footnotesize Progression naturelle}\\[3pt]
      \begin{tikzpicture}[font=\footnotesize]
        \draw[insarouge!40,thick] (0.2,0.05) -- (0.2,-3.55);
        \node[num] at (0.2,0)     {1};
        \node[num] at (0.2,-0.9) {2};
        \node[num] at (0.2,-1.8)  {3};
        \node[num] at (0.2,-2.7) {4};
        \node[num] at (0.2,-3.6)  {5};
        \node[anchor=west,text width=3.9cm,align=left] at (0.5,0)     {{\bfseries Outdoor}\\[-1pt]{\scriptsize randonneurs et groupes}};
        \node[anchor=west,text width=3.9cm,align=left] at (0.5,-0.9) {{\bfseries Refuges et clubs}\\[-1pt]{\scriptsize points de relais locaux}};
        \node[anchor=west,text width=3.9cm,align=left] at (0.5,-1.8)  {{\bfseries Associations}\\[-1pt]{\scriptsize ADRASEC, communautés structurées}};
        \node[anchor=west,text width=3.9cm,align=left] at (0.5,-2.7) {{\bfseries ONG, communes, pros}\\[-1pt]{\scriptsize déploiements plus larges}};
        \node[anchor=west,text width=3.9cm,align=left] at (0.5,-3.6)  {{\bfseries Grand public}\\[-1pt]{\scriptsize familles, voyageurs, une fois le réseau visible}};
      \end{tikzpicture}
    \end{column}
    \begin{column}{0.25\textwidth}
      \begin{idee}\footnotesize\centering
        {\sffamily\bfseries\color{insarouge}L'objectif}\\[3pt]
        Créer d'abord une zone où le produit est vraiment utile…\\[3pt]
        …puis utiliser cette densité, les retours terrain et la preuve d'usage pour élargir.
      \end{idee}
    \end{column}
  \end{columns}
  \vspace{0.05cm}
  \begin{pcompact}\centering\footnotesize
    Ne pas viser «\,tout le monde\,» au lancement~: gagner un premier territoire d'usage, puis élargir.
  \end{pcompact}
  \srcnote{Choix du marché de départ d'après le SWOT de l'équipe~; les seuils de décision sont dans la partie Étude quantitative.}
\end{frame}
              % 9.  SWOT et enjeux
%%%%%%%%%%%%%%%%%%%%%%%%%%%%%%%%%%%%%%%%%%%%%%%%%%%%%%%%%%%%%%%%%%%%%%%%%%%%%%
%  Partie « Étude quantitative » (7 diapositives) -- à importer avec %%%%%%%%%%%%%%%%%%%%%%%%%%%%%%%%%%%%%%%%%%%%%%%%%%%%%%%%%%%%%%%%%%%%%%%%%%%%%%
%  Partie « Étude quantitative » (7 diapositives) -- à importer avec %%%%%%%%%%%%%%%%%%%%%%%%%%%%%%%%%%%%%%%%%%%%%%%%%%%%%%%%%%%%%%%%%%%%%%%%%%%%%%
%  Partie « Étude quantitative » (7 diapositives) -- à importer avec \include{chap_quanti}
%  Fichier de contenu : ni préambule ni \begin{document}, commence par \section.
%
%  Contenu : (1) ce qu'il faut valider, (2) qui interroger et où, (3) plan du questionnaire,
%  (4) questions avant la présentation du concept, (5) questions après, (6) le prix,
%  (7) lecture des résultats et seuils de décision.
%  Cette partie reprend la méthodologie de l'équipe (étude quantitative, 3 segments, échantillon
%  non aléatoire par choix raisonné, 9 sections avant / après un concept neutre, prix maximal
%  puis « À 69 € … ? ») et la complète par : les questions rédigées, les lieux de diffusion,
%  des tailles d'échantillon avec marge d'erreur, et des seuils de décision.
%  Le questionnaire complet (formulations, échelles, logique de saut, variante secours) est
%  dans le fichier autonome questionnaire.tex ; les numéros de question sont les mêmes.
%
%  À DÉCIDER PAR L'ÉQUIPE (propositions de départ, signalées « proposé » dans les diapositives) :
%   - tailles d'échantillon : 150 (plein air), 100 (camping, voyage, famille), 60 (secours) ;
%   - seuils de décision de la dernière diapositive ;
%   - variante à trois prix (59, 69, 79 €) pour la question 20.
%  Marge d'erreur indicative : 1,96 × racine(0,25 / n), soit ± 8 points pour n = 150,
%  ± 9,8 pour n = 100, ± 12,7 pour n = 60 (échantillon aléatoire supposé : valeur indicative seulement).
%%%%%%%%%%%%%%%%%%%%%%%%%%%%%%%%%%%%%%%%%%%%%%%%%%%%%%%%%%%%%%%%%%%%%%%%%%%%%%
\section{Étude quantitative}

% numéro de question en rouge ; \qq{n°}{question}{réponses} : question suivie de ses réponses en gris
\newcommand{\qnum}[1]{{\sffamily\bfseries\color{insarouge}Q#1}}
\newcommand{\qq}[3]{\qnum{#1}\hspace{0.3em}#2\ {\color{insagris}\itshape#3}\par}
% titre de section de questionnaire, dans une carte
\newcommand{\qsec}[2]{\begingroup\footnotesize\sffamily\bfseries{\color{insarouge}#1}\hspace{0.4em}#2\par\endgroup\vspace{1.5pt}}
% cellule à remplir du tableau de résultats (diapositive 7)
\newcommand{\qvide}{\cellcolor{insabeige}}

%---------------------------------------------------------------------------
% Diapositive 1 : ce qu'il faut valider (la chaîne problème -> prix)
%---------------------------------------------------------------------------
\begin{frame}{Étude quantitative~: ce qu'il faut valider}
  \centering
  \begin{tikzpicture}[font=\scriptsize]
    \tikzset{etp/.style={minimum width=1.95cm,minimum height=0.62cm,inner xsep=3pt,font=\scriptsize\bfseries},
             detp/.style={anchor=north,text width=2.05cm,align=center,font=\scriptsize}}
    \node[pastille,etp,draw=insagris,fill=insagris!12,text=black]  (p1) at (0.98,0)  {PROBLÈME};
    \node[pastille,etp,draw=insagris,fill=insagris!12,text=black]  (p2) at (3.26,0)  {BESOIN};
    \node[pastille,etp,draw=insaorange,fill=insapeche,text=black]  (p3) at (5.54,0)  {SOLUTION};
    \node[pastille,etp,draw=insaorange,fill=insapeche,text=black]  (p4) at (7.82,0)  {INTÉRÊT};
    \node[pastille,etp,draw=insarouge,fill=insarose,text=black]    (p5) at (10.1,0) {ADOPTION};
    \node[pastille,etp,draw=insarouge,fill=insarose,text=black]    (p6) at (12.38,0) {PRIX};
    \foreach \a/\b in {p1/p2,p2/p3,p3/p4,p4/p5,p5/p6} {\draw[fluxr] (\a) -- (\b);}
    \node[detp] at ([yshift=-0.12cm]p1.south) {Se retrouve-\mbox{t-on} sans réseau~?\\{\color{insarouge}\qnum{7}}};
    \node[detp] at ([yshift=-0.12cm]p2.south) {Faut-il alors pouvoir communiquer~?\\{\color{insarouge}\qnum{10}, \qnum{11}}};
    \node[detp] at ([yshift=-0.12cm]p3.south) {La solution paraît-elle utile~?\\{\color{insarouge}\qnum{12}, \qnum{13}}};
    \node[detp] at ([yshift=-0.12cm]p4.south) {Dans quelles situations l'emploie-\mbox{rait-on}~?\\{\color{insarouge}\qnum{14}}};
    \node[detp] at ([yshift=-0.12cm]p5.south) {L'achèterait-\mbox{on}~? Qu'est-ce qui freine~?\\{\color{insarouge}\qnum{15} à \qnum{18}}};
    \node[detp] at ([yshift=-0.12cm]p6.south) {À quel prix~? 69\,€ passe-t-il~?\\{\color{insarouge}\qnum{19} à \qnum{22}}};
    % regroupements A, B, C
    \draw[decorate,decoration={brace,mirror,amplitude=4pt},insagris,thick] (-0.05,-1.95) -- (4.15,-1.95)
      node[midway,below=6pt,font=\scriptsize\bfseries,text=insagris] {A · Le problème est-il réel~?};
    \draw[decorate,decoration={brace,mirror,amplitude=4pt},insaorange,thick] (4.6,-1.95) -- (8.8,-1.95)
      node[midway,below=6pt,font=\scriptsize\bfseries,text=insaorange] {B · La solution plaît-elle~?};
    \draw[decorate,decoration={brace,mirror,amplitude=4pt},insarouge,thick] (9.1,-1.95) -- (13.3,-1.95)
      node[midway,below=6pt,font=\scriptsize\bfseries,text=insarouge] {C · Le potentiel commercial};
  \end{tikzpicture}

  \vspace{0.0cm}
  \begin{columns}[T]
    \begin{column}{0.48\textwidth}
      \begin{carte}[height=1.45cm,valign=top]
        \raggedright\gros{«\,Pas de réseau\,» ≠ besoin}\\[1pt]
        {\footnotesize Se retrouver sans réseau ne prouve pas qu'on veuille communiquer~: on mesure les deux.}
      \end{carte}
    \end{column}
    \begin{column}{0.48\textwidth}
      \begin{carte}[height=1.45cm,valign=top]
        \raggedright\gros{«\,Intéressant\,» ≠ achat}\\[1pt]
        {\footnotesize Un concept qui plaît ne se vend pas forcément~: on mesure l'intention, puis le prix.}
      \end{carte}
    \end{column}
  \end{columns}
  \vspace{0.05cm}
  \begin{pcompact}\centering\footnotesize
    Objectif final~: comparer les trois segments, retenir ceux où besoin et potentiel sont les plus forts.
  \end{pcompact}
  \srcnote{Méthodologie de l'équipe. Les entretiens qualitatifs (15 à 20, première vague) disent \emph{pourquoi} et \emph{comment}~; le questionnaire mesure \emph{combien}.}
\end{frame}

%---------------------------------------------------------------------------
% Diapositive 2 : qui interroger, où les trouver, combien
%---------------------------------------------------------------------------
\begin{frame}{Qui interroger, et où les trouver~?}
  \scriptsize\renewcommand{\arraystretch}{1.2}\setlength{\tabcolsep}{4pt}
  \hyphenpenalty=10000 \exhyphenpenalty=10000
  \noindent\begin{tabular}{@{}>{\raggedright\arraybackslash\bfseries}p{2.3cm}>{\raggedright\arraybackslash}p{3.6cm}>{\raggedright\arraybackslash}p{4.8cm}>{\centering\arraybackslash}p{0.8cm}>{\centering\arraybackslash}p{1.0cm}@{}}
    \rowcolor{insarose}\textcolor{insarouge}{Segment} & \textcolor{insarouge}{\bfseries Qui interroger} & \textcolor{insarouge}{\bfseries Où les trouver} & \textcolor{insarouge}{\bfseries Cible} & \textcolor{insarouge}{\bfseries Marge}\\
    1 Plein air, sport outdoor &
      randonneurs, trailers, vététistes, alpinistes, skieurs de randonnée &
      clubs de randonnée et de montagne, groupes de pratiquants en ligne, magasins outdoor, refuges (affiche avec QR code) & 150 & ±\,8\,pts\\
    \arrayrulecolor{insagris!30}\hline
    2 Campeurs, voyageurs, familles &
      campeurs, camping-caristes, voyageurs, familles qui sortent en zone isolée &
      campings (QR code à l'accueil), groupes de camping-caristes et de voyageurs, associations de parents, réseau personnel & 100 & ±\,10\,pts\\
    \hline
    3 Secours et gestion de crise &
      bénévoles et professionnels du secours, radioamateurs, élus et agents communaux, ONG &
      ADRASEC et associations de radioamateurs, Protection civile, Croix-Rouge, mairies, réseaux d'ONG & 60 & ±\,13\,pts\\
  \end{tabular}
  \vspace{0.2cm}
  \begin{columns}[T]
    \begin{column}{0.485\textwidth}
      \begin{carte}[height=1.55cm,valign=top]
        \raggedright\gros{Choix raisonné}\\[1pt]
        {\footnotesize Des personnes concernées, avec de la variété (âge, région, niveau de pratique). Question 1 = filtre~; «\,Autre\,» à part.}
      \end{carte}
    \end{column}
    \begin{column}{0.485\textwidth}
      \begin{carte}[height=1.55cm,valign=top]
        \raggedright\gros{Comparer, pas prévoir}\\[1pt]
        {\footnotesize Les segments entre eux~; pas de ventes ni de part de marché. Marge indicative ≈ ±\,98\,/\,$\sqrt{n}$ points.}
      \end{carte}
    \end{column}
  \end{columns}
  \srcnote{Tailles proposées par l'équipe (310 réponses au total), à ajuster~; marge à 95\,\% pour $p$ = 50\,\%, valable pour un tirage aléatoire~; la population du segment 3 étant petite, sa marge réelle est plus faible. Canaux~: ceux du recrutement des entretiens.}
\end{frame}

%---------------------------------------------------------------------------
% Diapositive 3 : plan du questionnaire (neuf sections, concept neutre au milieu)
%---------------------------------------------------------------------------
\begin{frame}{Plan du questionnaire~: neuf sections}
  \centering
  \begin{tikzpicture}[font=\scriptsize]
    \tikzset{sec/.style={draw=insagris,fill=insagris!10,rounded corners=3pt,minimum width=4.3cm,minimum height=0.55cm,
                         inner sep=3pt,align=left,anchor=west,text width=4.1cm},
             sec2/.style={draw=insaorange,fill=insapeche,rounded corners=3pt,minimum width=4.5cm,minimum height=0.55cm,
                         inner sep=3pt,align=left,anchor=west,text width=4.3cm}}
    % colonne de gauche : avant la présentation
    \node[font=\footnotesize\bfseries\sffamily,text=insagris,anchor=west] at (0,1.45) {AVANT~: problème et besoin};
    \node[sec] at (0,0.85)  {{\bfseries 1 Profil}~: quel segment~?};
    \node[sec] at (0,0.2)   {{\bfseries 2 Habitudes}~: pratiques actuelles};
    \node[sec] at (0,-0.45) {{\bfseries 3 Problème}~: sans réseau~?};
    \node[sec] at (0,-1.1)  {{\bfseries 4 Besoin}~: communiquer, pourquoi~?};
    % centre : concept neutre
    \node[fill=insarose,draw=insarouge,thick,rounded corners=4pt,text width=3.55cm,inner sep=5pt,align=left,anchor=north west]
      at (4.75,1.6)
      {{\sffamily\bfseries\color{insarouge}Concept présenté, sans nom ni prix}\\[2pt]
       \scriptsize Un petit appareil radio de poche, relié au téléphone par Bluetooth. Il envoie des messages texte à des proches
       ayant le même appareil, sans réseau mobile, sans internet, sans abonnement. Portée~: de 0,5 à 15\,km selon le terrain.
       Texte seulement~; ce n'est pas une balise de détresse.};
    % colonne de droite : après
    \node[font=\footnotesize\bfseries\sffamily,text=insaorange,anchor=west] at (8.95,1.45) {APRÈS~: solution et adoption};
    \node[sec2] at (8.95,0.85)  {{\bfseries 5 Intérêt}~: la solution est-elle utile~?};
    \node[sec2] at (8.95,0.2)   {{\bfseries 6 Usages}~: dans quelles situations~?};
    \node[sec2] at (8.95,-0.45) {{\bfseries 7 Freins}~: qu'est-ce qui bloque~?};
    \node[sec2] at (8.95,-1.1)  {{\bfseries 8 Intention}~: acheter ou adopter~?};
    \node[sec2] at (8.95,-1.75) {{\bfseries 9 Prix}~: maximal, puis 69\,€};
    % flèches
    \draw[fluxr] (4.35,-0.2) -- (4.7,-0.2);
    \draw[fluxr] (8.35,-0.2) -- (8.9,-0.2);
  \end{tikzpicture}

  \vspace{0.05cm}
  \begin{columns}[T]
    \begin{column}{0.325\textwidth}
      \begin{carte}[height=1.25cm,valign=top]
        \raggedright\gros{23 questions}\\[1pt]{\footnotesize 8 à 10 minutes, en ligne ou sur papier.}
      \end{carte}
    \end{column}
    \begin{column}{0.325\textwidth}
      \begin{carte}[height=1.25cm,valign=top]
        \raggedright\gros{Concept neutre}\\[1pt]{\footnotesize après la section 4, pour ne pas influencer.}
      \end{carte}
    \end{column}
    \begin{column}{0.325\textwidth}
      \begin{carte}[height=1.25cm,valign=top]
        \raggedright\gros{Pré-test}\\[1pt]{\footnotesize dix personnes~; réponses anonymes (RGPD).}
      \end{carte}
    \end{column}
  \end{columns}
  \srcnote{Structure de la méthodologie de l'équipe~; formulations complètes dans le fichier questionnaire.tex. Portée~: voir la partie Fonctionnement (\cite{techconnect}).}
\end{frame}

%---------------------------------------------------------------------------
% Diapositive 4 : questions 1 à 11 (avant la présentation)
%---------------------------------------------------------------------------
\begin{frame}{Avant la présentation~: questions 1 à 11}
  \begin{columns}[T]
    \begin{column}{0.485\textwidth}
      \begin{carte}[height=1.9cm,valign=top]
        \qsec{1}{Profil}\scriptsize\raggedright
        \qq{1}{À quel profil correspondez-vous le mieux~?}{plein air · camping, voyage, famille · secours et crise · autre}
        \qq{2}{Votre âge}{tranches}
        \qq{3}{Où vivez-vous~?}{ville · périurbain · campagne · montagne}
      \end{carte}
      \vspace{0.06cm}
      \begin{carte}[height=2.2cm,valign=top]
        \qsec{3}{Existence du problème}\scriptsize\raggedright
        \qq{7}{Vous retrouvez-vous dans une zone sans réseau mobile~?}{jamais → très souvent}
        \qq{8}{Où cela vous est-il arrivé~?}{montagne, campagne, étranger, ville, catastrophe…}
        \qq{9}{Comment communiquez-vous alors~?}{rien, talkies, balise, satellite…}
      \end{carte}
    \end{column}
    \begin{column}{0.485\textwidth}
      \begin{carte}[height=1.9cm,valign=top]
        \qsec{2}{Habitudes et pratiques}\scriptsize\raggedright
        \qq{4}{Quelles activités en pleine nature, ces 12 derniers mois~?}{choix multiple}
        \qq{5}{À quelle fréquence~?}{jamais → chaque semaine}
        \qq{6}{Avec qui sortez-vous~?}{seul · à deux · 3 à 5 · 6 ou plus}
      \end{carte}
      \vspace{0.06cm}
      \begin{carte}[height=2.2cm,valign=top]
        \qsec{4}{Importance du besoin}\scriptsize\raggedright
        \qq{10}{Sans réseau, importance de~: rester en contact avec son groupe, prévenir ses proches, partager sa position, coordonner une équipe, alerter les secours}{1 à 5}
        \qq{11}{Cela vous a-t-il déjà posé un vrai problème~?}{jamais · une fois · plusieurs fois}
      \end{carte}
    \end{column}
  \end{columns}
  \vspace{0.05cm}
  \begin{pcompact}\centering\footnotesize
    Secours~: mêmes questions, reformulées. Le concept n'est présenté qu'après la question 11.
  \end{pcompact}
  \srcnote{Sections 1 à 4 de la méthodologie de l'équipe~; réponses complètes et variante secours dans questionnaire.tex.}
\end{frame}

%---------------------------------------------------------------------------
% Diapositive 5 : questions 12 à 18 (après la présentation)
%---------------------------------------------------------------------------
\begin{frame}{Après la présentation~: questions 12 à 18}
  \begin{columns}[T]
    \begin{column}{0.485\textwidth}
      \begin{carte}[height=1.95cm,valign=top]
        \qsec{5}{Intérêt pour la solution}\scriptsize\raggedright
        \qq{12}{Dans quelle mesure cette solution vous paraît-elle utile~?}{pas du tout → très utile}
        \qq{13}{Qu'est-ce qui vous intéresse le plus~? (deux au plus)}{sans réseau · sans abonnement · messages privés · simplicité · prix}
      \end{carte}
      \vspace{0.1cm}
      \begin{carte}[height=1.95cm,valign=top]
        \qsec{7}{Freins à l'adoption}\scriptsize\raggedright
        \qq{15}{Qu'est-ce qui pourrait vous empêcher de l'acheter~?}{prix · équiper les autres · portée · autonomie · complexité · utilité · déjà équipé · marque inconnue}
        \qq{16}{Quel est le frein principal~?}{un seul choix}
      \end{carte}
    \end{column}
    \begin{column}{0.485\textwidth}
      \begin{carte}[height=1.95cm,valign=top]
        \qsec{6}{Usages envisagés}\scriptsize\raggedright
        \qq{14}{Dans quelles situations l'utiliseriez-vous~?}{randonnée · camping · voyage · sport · zones isolées · secours · catastrophe ou panne · événement · jamais}
      \end{carte}
      \vspace{0.1cm}
      \begin{carte}[height=1.95cm,valign=top]
        \qsec{8}{Intention d'achat ou d'adoption}\scriptsize\raggedright
        \qq{17}{Seriez-vous prêt à acheter ou adopter cette solution~?}{certainement pas → certainement}
        \qq{18}{Combien de personnes équiperiez-vous, vous compris~?}{1 · 2 · 3 à 4 · 5 à 9 · 10 ou plus}
      \end{carte}
    \end{column}
  \end{columns}
  \vspace{0.05cm}
  \begin{pcompact}\centering\footnotesize
    La question 18 donne la taille de groupe à équiper~: elle éclaire le choix kit seul, Duo ou Groupe.
  \end{pcompact}
  \srcnote{Sections 5 à 8 de la méthodologie de l'équipe~; échelles à cinq niveaux, ordre des réponses tiré au hasard quand c'est possible.}
\end{frame}

%---------------------------------------------------------------------------
% Diapositive 6 : le prix (questions 19 à 23)
%---------------------------------------------------------------------------
\begin{frame}{Le prix~: de la question ouverte à la question précise}
  \begin{columns}[T]
    \begin{column}{0.57\textwidth}
      \begin{tikzpicture}[font=\footnotesize]
        \draw[insarouge!40,thick] (0.2,0.2) -- (0.2,-3.35);
        \node[num] at (0.2,0)     {1};
        \node[num] at (0.2,-1.2)  {2};
        \node[num] at (0.2,-2.4)  {3};
        \node[anchor=west,text width=6.5cm,align=left] at (0.5,0.05)
          {\qnum{19}~: {\bfseries Quel prix maximal paieriez-vous, sans abonnement~?}\\[-1pt]
           {\scriptsize réponse libre, en euros, \emph{avant} de montrer 69\,€~: pas d'ancrage.}};
        \node[anchor=west,text width=6.5cm,align=left] at (0.5,-1.2)
          {\qnum{20}~: {\bfseries À 69\,€, sans abonnement, l'achèteriez-vous~?}\\[-1pt]
           {\scriptsize oui, peut-être, non~: teste l'hypothèse de prix.}};
        \node[anchor=west,text width=6.5cm,align=left] at (0.5,-2.4)
          {\qnum{21}~: {\bfseries Pour plusieurs personnes~: 2 kits à 139\,€, 4 à 249\,€, un seul à 69\,€, aucun~?}\\[-1pt]
           {\scriptsize teste les lots Duo et Groupe de la partie SWOT.}};
        \node[anchor=west,text width=6.5cm,align=left,font=\scriptsize,text=insagris] at (0.5,-3.3)
          {\qnum{22}~: où l'acheter (boutique, site, club, commune)~; \qnum{23}~: contact facultatif pour un test ou une précommande.};
      \end{tikzpicture}
    \end{column}
    \begin{column}{0.4\textwidth}
      \begin{encadre}{Comment lire}
        \begin{puces}
          \item \textbf{Part solvable à 69\,€}~: part des prix maximaux (Q19) au moins égaux à 69\,€.
          \item \textbf{Recette à un prix $p$}~: $p$ × part des prix maximaux ≥ $p$~; comparer 59, 69 et 79\,€.
          \item \textbf{Effet du prix}~: écart entre l'intention (Q17) et le «\,oui\,» à 69\,€ (Q20).
        \end{puces}
      \end{encadre}
    \end{column}
  \end{columns}
  \vspace{0.05cm}
  \begin{pcompact}\centering\footnotesize
    Variante proposée~: tirer au hasard 59, 69 ou 79\,€ à la question 20, pour lire trois points de la courbe de demande.
  \end{pcompact}
  \srcnote{Ordre « prix maximal puis 69\,€ » de la méthodologie de l'équipe. Sensibilité du modèle financier~: ±\,5\,€ de prix = ±\,92\,k€ de résultat en année 3 (partie SWOT).}
\end{frame}

%---------------------------------------------------------------------------
% Diapositive 7 : lire les résultats, seuils de décision
%---------------------------------------------------------------------------
\begin{frame}{Décider~: indicateurs et seuils, segment par segment}
  \scriptsize\renewcommand{\arraystretch}{1.25}\setlength{\tabcolsep}{4pt}
  \hyphenpenalty=10000 \exhyphenpenalty=10000
  \noindent\begin{tabular}{@{}>{\raggedright\arraybackslash}p{4.65cm}>{\centering\arraybackslash}p{0.7cm}>{\centering\arraybackslash}p{1.5cm}>{\centering\arraybackslash}p{1.75cm}!{\color{white}\vrule width 1.5pt}>{\centering\arraybackslash}p{1.75cm}!{\color{white}\vrule width 1.5pt}>{\centering\arraybackslash}p{1.75cm}@{}}
    \rowcolor{insarose}\textcolor{insarouge}{\bfseries Indicateur} & \textcolor{insarouge}{\bfseries Q} & \textcolor{insarouge}{\bfseries Seuil proposé} & \textcolor{insarouge}{\bfseries 1 Plein air} & \textcolor{insarouge}{\bfseries 2 Camping, famille} & \textcolor{insarouge}{\bfseries 3 Secours}\\
    Problème~: sans réseau «\,souvent\,» ou «\,très souvent\,» & 7 & ≥ 40\,\% & \qvide & \qvide & \qvide\\ \arrayrulecolor{white}\hline
    Besoin~: «\,important\,» ou «\,très important\,» (4 ou 5 sur 5) & 10 & ≥ 60\,\% & \qvide & \qvide & \qvide\\ \arrayrulecolor{white}\hline
    Problème déjà vécu au moins une fois & 11 & ≥ 30\,\% & \qvide & \qvide & \qvide\\ \arrayrulecolor{white}\hline
    Utilité~: «\,utile\,» ou «\,très utile\,» & 12 & ≥ 50\,\% & \qvide & \qvide & \qvide\\ \arrayrulecolor{white}\hline
    Intention~: «\,probablement\,» ou «\,certainement\,» & 17 & ≥ 30\,\% & \qvide & \qvide & \qvide\\ \arrayrulecolor{white}\hline
    Prix~: «\,oui\,» à 69\,€ & 20 & ≥ 20\,\% & \qvide & \qvide & \qvide\\ \arrayrulecolor{white}\hline
    Prix maximal, valeur médiane & 19 & ≥ 60\,€ & \qvide & \qvide & \qvide\\ \arrayrulecolor{white}\hline
    Appareils à équiper, valeur médiane & 18 & ≥ 2 & \qvide & \qvide & \qvide\\ \arrayrulecolor{white}\hline
  \end{tabular}
  \vspace{0.15cm}
  \begin{pcompact}\centering\footnotesize
    Règle proposée~: segment prioritaire s'il atteint 6 seuils sur 8~; si le «\,oui\,» à 69\,€ reste sous 10\,\% partout, revoir prix ou concept avant le moule.
  \end{pcompact}
  \srcnote{Seuils proposés par l'équipe, à valider \emph{avant} l'enquête (sinon on les ajuste après coup). L'échantillon n'étant pas aléatoire, ils servent à comparer les segments, pas à prévoir les ventes. Cases beiges~: à remplir avec les résultats.}
\end{frame}

%  Fichier de contenu : ni préambule ni \begin{document}, commence par \section.
%
%  Contenu : (1) ce qu'il faut valider, (2) qui interroger et où, (3) plan du questionnaire,
%  (4) questions avant la présentation du concept, (5) questions après, (6) le prix,
%  (7) lecture des résultats et seuils de décision.
%  Cette partie reprend la méthodologie de l'équipe (étude quantitative, 3 segments, échantillon
%  non aléatoire par choix raisonné, 9 sections avant / après un concept neutre, prix maximal
%  puis « À 69 € … ? ») et la complète par : les questions rédigées, les lieux de diffusion,
%  des tailles d'échantillon avec marge d'erreur, et des seuils de décision.
%  Le questionnaire complet (formulations, échelles, logique de saut, variante secours) est
%  dans le fichier autonome questionnaire.tex ; les numéros de question sont les mêmes.
%
%  À DÉCIDER PAR L'ÉQUIPE (propositions de départ, signalées « proposé » dans les diapositives) :
%   - tailles d'échantillon : 150 (plein air), 100 (camping, voyage, famille), 60 (secours) ;
%   - seuils de décision de la dernière diapositive ;
%   - variante à trois prix (59, 69, 79 €) pour la question 20.
%  Marge d'erreur indicative : 1,96 × racine(0,25 / n), soit ± 8 points pour n = 150,
%  ± 9,8 pour n = 100, ± 12,7 pour n = 60 (échantillon aléatoire supposé : valeur indicative seulement).
%%%%%%%%%%%%%%%%%%%%%%%%%%%%%%%%%%%%%%%%%%%%%%%%%%%%%%%%%%%%%%%%%%%%%%%%%%%%%%
\section{Étude quantitative}

% numéro de question en rouge ; \qq{n°}{question}{réponses} : question suivie de ses réponses en gris
\newcommand{\qnum}[1]{{\sffamily\bfseries\color{insarouge}Q#1}}
\newcommand{\qq}[3]{\qnum{#1}\hspace{0.3em}#2\ {\color{insagris}\itshape#3}\par}
% titre de section de questionnaire, dans une carte
\newcommand{\qsec}[2]{\begingroup\footnotesize\sffamily\bfseries{\color{insarouge}#1}\hspace{0.4em}#2\par\endgroup\vspace{1.5pt}}
% cellule à remplir du tableau de résultats (diapositive 7)
\newcommand{\qvide}{\cellcolor{insabeige}}

%---------------------------------------------------------------------------
% Diapositive 1 : ce qu'il faut valider (la chaîne problème -> prix)
%---------------------------------------------------------------------------
\begin{frame}{Étude quantitative~: ce qu'il faut valider}
  \centering
  \begin{tikzpicture}[font=\scriptsize]
    \tikzset{etp/.style={minimum width=1.95cm,minimum height=0.62cm,inner xsep=3pt,font=\scriptsize\bfseries},
             detp/.style={anchor=north,text width=2.05cm,align=center,font=\scriptsize}}
    \node[pastille,etp,draw=insagris,fill=insagris!12,text=black]  (p1) at (0.98,0)  {PROBLÈME};
    \node[pastille,etp,draw=insagris,fill=insagris!12,text=black]  (p2) at (3.26,0)  {BESOIN};
    \node[pastille,etp,draw=insaorange,fill=insapeche,text=black]  (p3) at (5.54,0)  {SOLUTION};
    \node[pastille,etp,draw=insaorange,fill=insapeche,text=black]  (p4) at (7.82,0)  {INTÉRÊT};
    \node[pastille,etp,draw=insarouge,fill=insarose,text=black]    (p5) at (10.1,0) {ADOPTION};
    \node[pastille,etp,draw=insarouge,fill=insarose,text=black]    (p6) at (12.38,0) {PRIX};
    \foreach \a/\b in {p1/p2,p2/p3,p3/p4,p4/p5,p5/p6} {\draw[fluxr] (\a) -- (\b);}
    \node[detp] at ([yshift=-0.12cm]p1.south) {Se retrouve-\mbox{t-on} sans réseau~?\\{\color{insarouge}\qnum{7}}};
    \node[detp] at ([yshift=-0.12cm]p2.south) {Faut-il alors pouvoir communiquer~?\\{\color{insarouge}\qnum{10}, \qnum{11}}};
    \node[detp] at ([yshift=-0.12cm]p3.south) {La solution paraît-elle utile~?\\{\color{insarouge}\qnum{12}, \qnum{13}}};
    \node[detp] at ([yshift=-0.12cm]p4.south) {Dans quelles situations l'emploie-\mbox{rait-on}~?\\{\color{insarouge}\qnum{14}}};
    \node[detp] at ([yshift=-0.12cm]p5.south) {L'achèterait-\mbox{on}~? Qu'est-ce qui freine~?\\{\color{insarouge}\qnum{15} à \qnum{18}}};
    \node[detp] at ([yshift=-0.12cm]p6.south) {À quel prix~? 69\,€ passe-t-il~?\\{\color{insarouge}\qnum{19} à \qnum{22}}};
    % regroupements A, B, C
    \draw[decorate,decoration={brace,mirror,amplitude=4pt},insagris,thick] (-0.05,-1.95) -- (4.15,-1.95)
      node[midway,below=6pt,font=\scriptsize\bfseries,text=insagris] {A · Le problème est-il réel~?};
    \draw[decorate,decoration={brace,mirror,amplitude=4pt},insaorange,thick] (4.6,-1.95) -- (8.8,-1.95)
      node[midway,below=6pt,font=\scriptsize\bfseries,text=insaorange] {B · La solution plaît-elle~?};
    \draw[decorate,decoration={brace,mirror,amplitude=4pt},insarouge,thick] (9.1,-1.95) -- (13.3,-1.95)
      node[midway,below=6pt,font=\scriptsize\bfseries,text=insarouge] {C · Le potentiel commercial};
  \end{tikzpicture}

  \vspace{0.0cm}
  \begin{columns}[T]
    \begin{column}{0.48\textwidth}
      \begin{carte}[height=1.45cm,valign=top]
        \raggedright\gros{«\,Pas de réseau\,» ≠ besoin}\\[1pt]
        {\footnotesize Se retrouver sans réseau ne prouve pas qu'on veuille communiquer~: on mesure les deux.}
      \end{carte}
    \end{column}
    \begin{column}{0.48\textwidth}
      \begin{carte}[height=1.45cm,valign=top]
        \raggedright\gros{«\,Intéressant\,» ≠ achat}\\[1pt]
        {\footnotesize Un concept qui plaît ne se vend pas forcément~: on mesure l'intention, puis le prix.}
      \end{carte}
    \end{column}
  \end{columns}
  \vspace{0.05cm}
  \begin{pcompact}\centering\footnotesize
    Objectif final~: comparer les trois segments, retenir ceux où besoin et potentiel sont les plus forts.
  \end{pcompact}
  \srcnote{Méthodologie de l'équipe. Les entretiens qualitatifs (15 à 20, première vague) disent \emph{pourquoi} et \emph{comment}~; le questionnaire mesure \emph{combien}.}
\end{frame}

%---------------------------------------------------------------------------
% Diapositive 2 : qui interroger, où les trouver, combien
%---------------------------------------------------------------------------
\begin{frame}{Qui interroger, et où les trouver~?}
  \scriptsize\renewcommand{\arraystretch}{1.2}\setlength{\tabcolsep}{4pt}
  \hyphenpenalty=10000 \exhyphenpenalty=10000
  \noindent\begin{tabular}{@{}>{\raggedright\arraybackslash\bfseries}p{2.3cm}>{\raggedright\arraybackslash}p{3.6cm}>{\raggedright\arraybackslash}p{4.8cm}>{\centering\arraybackslash}p{0.8cm}>{\centering\arraybackslash}p{1.0cm}@{}}
    \rowcolor{insarose}\textcolor{insarouge}{Segment} & \textcolor{insarouge}{\bfseries Qui interroger} & \textcolor{insarouge}{\bfseries Où les trouver} & \textcolor{insarouge}{\bfseries Cible} & \textcolor{insarouge}{\bfseries Marge}\\
    1 Plein air, sport outdoor &
      randonneurs, trailers, vététistes, alpinistes, skieurs de randonnée &
      clubs de randonnée et de montagne, groupes de pratiquants en ligne, magasins outdoor, refuges (affiche avec QR code) & 150 & ±\,8\,pts\\
    \arrayrulecolor{insagris!30}\hline
    2 Campeurs, voyageurs, familles &
      campeurs, camping-caristes, voyageurs, familles qui sortent en zone isolée &
      campings (QR code à l'accueil), groupes de camping-caristes et de voyageurs, associations de parents, réseau personnel & 100 & ±\,10\,pts\\
    \hline
    3 Secours et gestion de crise &
      bénévoles et professionnels du secours, radioamateurs, élus et agents communaux, ONG &
      ADRASEC et associations de radioamateurs, Protection civile, Croix-Rouge, mairies, réseaux d'ONG & 60 & ±\,13\,pts\\
  \end{tabular}
  \vspace{0.2cm}
  \begin{columns}[T]
    \begin{column}{0.485\textwidth}
      \begin{carte}[height=1.55cm,valign=top]
        \raggedright\gros{Choix raisonné}\\[1pt]
        {\footnotesize Des personnes concernées, avec de la variété (âge, région, niveau de pratique). Question 1 = filtre~; «\,Autre\,» à part.}
      \end{carte}
    \end{column}
    \begin{column}{0.485\textwidth}
      \begin{carte}[height=1.55cm,valign=top]
        \raggedright\gros{Comparer, pas prévoir}\\[1pt]
        {\footnotesize Les segments entre eux~; pas de ventes ni de part de marché. Marge indicative ≈ ±\,98\,/\,$\sqrt{n}$ points.}
      \end{carte}
    \end{column}
  \end{columns}
  \srcnote{Tailles proposées par l'équipe (310 réponses au total), à ajuster~; marge à 95\,\% pour $p$ = 50\,\%, valable pour un tirage aléatoire~; la population du segment 3 étant petite, sa marge réelle est plus faible. Canaux~: ceux du recrutement des entretiens.}
\end{frame}

%---------------------------------------------------------------------------
% Diapositive 3 : plan du questionnaire (neuf sections, concept neutre au milieu)
%---------------------------------------------------------------------------
\begin{frame}{Plan du questionnaire~: neuf sections}
  \centering
  \begin{tikzpicture}[font=\scriptsize]
    \tikzset{sec/.style={draw=insagris,fill=insagris!10,rounded corners=3pt,minimum width=4.3cm,minimum height=0.55cm,
                         inner sep=3pt,align=left,anchor=west,text width=4.1cm},
             sec2/.style={draw=insaorange,fill=insapeche,rounded corners=3pt,minimum width=4.5cm,minimum height=0.55cm,
                         inner sep=3pt,align=left,anchor=west,text width=4.3cm}}
    % colonne de gauche : avant la présentation
    \node[font=\footnotesize\bfseries\sffamily,text=insagris,anchor=west] at (0,1.45) {AVANT~: problème et besoin};
    \node[sec] at (0,0.85)  {{\bfseries 1 Profil}~: quel segment~?};
    \node[sec] at (0,0.2)   {{\bfseries 2 Habitudes}~: pratiques actuelles};
    \node[sec] at (0,-0.45) {{\bfseries 3 Problème}~: sans réseau~?};
    \node[sec] at (0,-1.1)  {{\bfseries 4 Besoin}~: communiquer, pourquoi~?};
    % centre : concept neutre
    \node[fill=insarose,draw=insarouge,thick,rounded corners=4pt,text width=3.55cm,inner sep=5pt,align=left,anchor=north west]
      at (4.75,1.6)
      {{\sffamily\bfseries\color{insarouge}Concept présenté, sans nom ni prix}\\[2pt]
       \scriptsize Un petit appareil radio de poche, relié au téléphone par Bluetooth. Il envoie des messages texte à des proches
       ayant le même appareil, sans réseau mobile, sans internet, sans abonnement. Portée~: de 0,5 à 15\,km selon le terrain.
       Texte seulement~; ce n'est pas une balise de détresse.};
    % colonne de droite : après
    \node[font=\footnotesize\bfseries\sffamily,text=insaorange,anchor=west] at (8.95,1.45) {APRÈS~: solution et adoption};
    \node[sec2] at (8.95,0.85)  {{\bfseries 5 Intérêt}~: la solution est-elle utile~?};
    \node[sec2] at (8.95,0.2)   {{\bfseries 6 Usages}~: dans quelles situations~?};
    \node[sec2] at (8.95,-0.45) {{\bfseries 7 Freins}~: qu'est-ce qui bloque~?};
    \node[sec2] at (8.95,-1.1)  {{\bfseries 8 Intention}~: acheter ou adopter~?};
    \node[sec2] at (8.95,-1.75) {{\bfseries 9 Prix}~: maximal, puis 69\,€};
    % flèches
    \draw[fluxr] (4.35,-0.2) -- (4.7,-0.2);
    \draw[fluxr] (8.35,-0.2) -- (8.9,-0.2);
  \end{tikzpicture}

  \vspace{0.05cm}
  \begin{columns}[T]
    \begin{column}{0.325\textwidth}
      \begin{carte}[height=1.25cm,valign=top]
        \raggedright\gros{23 questions}\\[1pt]{\footnotesize 8 à 10 minutes, en ligne ou sur papier.}
      \end{carte}
    \end{column}
    \begin{column}{0.325\textwidth}
      \begin{carte}[height=1.25cm,valign=top]
        \raggedright\gros{Concept neutre}\\[1pt]{\footnotesize après la section 4, pour ne pas influencer.}
      \end{carte}
    \end{column}
    \begin{column}{0.325\textwidth}
      \begin{carte}[height=1.25cm,valign=top]
        \raggedright\gros{Pré-test}\\[1pt]{\footnotesize dix personnes~; réponses anonymes (RGPD).}
      \end{carte}
    \end{column}
  \end{columns}
  \srcnote{Structure de la méthodologie de l'équipe~; formulations complètes dans le fichier questionnaire.tex. Portée~: voir la partie Fonctionnement (\cite{techconnect}).}
\end{frame}

%---------------------------------------------------------------------------
% Diapositive 4 : questions 1 à 11 (avant la présentation)
%---------------------------------------------------------------------------
\begin{frame}{Avant la présentation~: questions 1 à 11}
  \begin{columns}[T]
    \begin{column}{0.485\textwidth}
      \begin{carte}[height=1.9cm,valign=top]
        \qsec{1}{Profil}\scriptsize\raggedright
        \qq{1}{À quel profil correspondez-vous le mieux~?}{plein air · camping, voyage, famille · secours et crise · autre}
        \qq{2}{Votre âge}{tranches}
        \qq{3}{Où vivez-vous~?}{ville · périurbain · campagne · montagne}
      \end{carte}
      \vspace{0.06cm}
      \begin{carte}[height=2.2cm,valign=top]
        \qsec{3}{Existence du problème}\scriptsize\raggedright
        \qq{7}{Vous retrouvez-vous dans une zone sans réseau mobile~?}{jamais → très souvent}
        \qq{8}{Où cela vous est-il arrivé~?}{montagne, campagne, étranger, ville, catastrophe…}
        \qq{9}{Comment communiquez-vous alors~?}{rien, talkies, balise, satellite…}
      \end{carte}
    \end{column}
    \begin{column}{0.485\textwidth}
      \begin{carte}[height=1.9cm,valign=top]
        \qsec{2}{Habitudes et pratiques}\scriptsize\raggedright
        \qq{4}{Quelles activités en pleine nature, ces 12 derniers mois~?}{choix multiple}
        \qq{5}{À quelle fréquence~?}{jamais → chaque semaine}
        \qq{6}{Avec qui sortez-vous~?}{seul · à deux · 3 à 5 · 6 ou plus}
      \end{carte}
      \vspace{0.06cm}
      \begin{carte}[height=2.2cm,valign=top]
        \qsec{4}{Importance du besoin}\scriptsize\raggedright
        \qq{10}{Sans réseau, importance de~: rester en contact avec son groupe, prévenir ses proches, partager sa position, coordonner une équipe, alerter les secours}{1 à 5}
        \qq{11}{Cela vous a-t-il déjà posé un vrai problème~?}{jamais · une fois · plusieurs fois}
      \end{carte}
    \end{column}
  \end{columns}
  \vspace{0.05cm}
  \begin{pcompact}\centering\footnotesize
    Secours~: mêmes questions, reformulées. Le concept n'est présenté qu'après la question 11.
  \end{pcompact}
  \srcnote{Sections 1 à 4 de la méthodologie de l'équipe~; réponses complètes et variante secours dans questionnaire.tex.}
\end{frame}

%---------------------------------------------------------------------------
% Diapositive 5 : questions 12 à 18 (après la présentation)
%---------------------------------------------------------------------------
\begin{frame}{Après la présentation~: questions 12 à 18}
  \begin{columns}[T]
    \begin{column}{0.485\textwidth}
      \begin{carte}[height=1.95cm,valign=top]
        \qsec{5}{Intérêt pour la solution}\scriptsize\raggedright
        \qq{12}{Dans quelle mesure cette solution vous paraît-elle utile~?}{pas du tout → très utile}
        \qq{13}{Qu'est-ce qui vous intéresse le plus~? (deux au plus)}{sans réseau · sans abonnement · messages privés · simplicité · prix}
      \end{carte}
      \vspace{0.1cm}
      \begin{carte}[height=1.95cm,valign=top]
        \qsec{7}{Freins à l'adoption}\scriptsize\raggedright
        \qq{15}{Qu'est-ce qui pourrait vous empêcher de l'acheter~?}{prix · équiper les autres · portée · autonomie · complexité · utilité · déjà équipé · marque inconnue}
        \qq{16}{Quel est le frein principal~?}{un seul choix}
      \end{carte}
    \end{column}
    \begin{column}{0.485\textwidth}
      \begin{carte}[height=1.95cm,valign=top]
        \qsec{6}{Usages envisagés}\scriptsize\raggedright
        \qq{14}{Dans quelles situations l'utiliseriez-vous~?}{randonnée · camping · voyage · sport · zones isolées · secours · catastrophe ou panne · événement · jamais}
      \end{carte}
      \vspace{0.1cm}
      \begin{carte}[height=1.95cm,valign=top]
        \qsec{8}{Intention d'achat ou d'adoption}\scriptsize\raggedright
        \qq{17}{Seriez-vous prêt à acheter ou adopter cette solution~?}{certainement pas → certainement}
        \qq{18}{Combien de personnes équiperiez-vous, vous compris~?}{1 · 2 · 3 à 4 · 5 à 9 · 10 ou plus}
      \end{carte}
    \end{column}
  \end{columns}
  \vspace{0.05cm}
  \begin{pcompact}\centering\footnotesize
    La question 18 donne la taille de groupe à équiper~: elle éclaire le choix kit seul, Duo ou Groupe.
  \end{pcompact}
  \srcnote{Sections 5 à 8 de la méthodologie de l'équipe~; échelles à cinq niveaux, ordre des réponses tiré au hasard quand c'est possible.}
\end{frame}

%---------------------------------------------------------------------------
% Diapositive 6 : le prix (questions 19 à 23)
%---------------------------------------------------------------------------
\begin{frame}{Le prix~: de la question ouverte à la question précise}
  \begin{columns}[T]
    \begin{column}{0.57\textwidth}
      \begin{tikzpicture}[font=\footnotesize]
        \draw[insarouge!40,thick] (0.2,0.2) -- (0.2,-3.35);
        \node[num] at (0.2,0)     {1};
        \node[num] at (0.2,-1.2)  {2};
        \node[num] at (0.2,-2.4)  {3};
        \node[anchor=west,text width=6.5cm,align=left] at (0.5,0.05)
          {\qnum{19}~: {\bfseries Quel prix maximal paieriez-vous, sans abonnement~?}\\[-1pt]
           {\scriptsize réponse libre, en euros, \emph{avant} de montrer 69\,€~: pas d'ancrage.}};
        \node[anchor=west,text width=6.5cm,align=left] at (0.5,-1.2)
          {\qnum{20}~: {\bfseries À 69\,€, sans abonnement, l'achèteriez-vous~?}\\[-1pt]
           {\scriptsize oui, peut-être, non~: teste l'hypothèse de prix.}};
        \node[anchor=west,text width=6.5cm,align=left] at (0.5,-2.4)
          {\qnum{21}~: {\bfseries Pour plusieurs personnes~: 2 kits à 139\,€, 4 à 249\,€, un seul à 69\,€, aucun~?}\\[-1pt]
           {\scriptsize teste les lots Duo et Groupe de la partie SWOT.}};
        \node[anchor=west,text width=6.5cm,align=left,font=\scriptsize,text=insagris] at (0.5,-3.3)
          {\qnum{22}~: où l'acheter (boutique, site, club, commune)~; \qnum{23}~: contact facultatif pour un test ou une précommande.};
      \end{tikzpicture}
    \end{column}
    \begin{column}{0.4\textwidth}
      \begin{encadre}{Comment lire}
        \begin{puces}
          \item \textbf{Part solvable à 69\,€}~: part des prix maximaux (Q19) au moins égaux à 69\,€.
          \item \textbf{Recette à un prix $p$}~: $p$ × part des prix maximaux ≥ $p$~; comparer 59, 69 et 79\,€.
          \item \textbf{Effet du prix}~: écart entre l'intention (Q17) et le «\,oui\,» à 69\,€ (Q20).
        \end{puces}
      \end{encadre}
    \end{column}
  \end{columns}
  \vspace{0.05cm}
  \begin{pcompact}\centering\footnotesize
    Variante proposée~: tirer au hasard 59, 69 ou 79\,€ à la question 20, pour lire trois points de la courbe de demande.
  \end{pcompact}
  \srcnote{Ordre « prix maximal puis 69\,€ » de la méthodologie de l'équipe. Sensibilité du modèle financier~: ±\,5\,€ de prix = ±\,92\,k€ de résultat en année 3 (partie SWOT).}
\end{frame}

%---------------------------------------------------------------------------
% Diapositive 7 : lire les résultats, seuils de décision
%---------------------------------------------------------------------------
\begin{frame}{Décider~: indicateurs et seuils, segment par segment}
  \scriptsize\renewcommand{\arraystretch}{1.25}\setlength{\tabcolsep}{4pt}
  \hyphenpenalty=10000 \exhyphenpenalty=10000
  \noindent\begin{tabular}{@{}>{\raggedright\arraybackslash}p{4.65cm}>{\centering\arraybackslash}p{0.7cm}>{\centering\arraybackslash}p{1.5cm}>{\centering\arraybackslash}p{1.75cm}!{\color{white}\vrule width 1.5pt}>{\centering\arraybackslash}p{1.75cm}!{\color{white}\vrule width 1.5pt}>{\centering\arraybackslash}p{1.75cm}@{}}
    \rowcolor{insarose}\textcolor{insarouge}{\bfseries Indicateur} & \textcolor{insarouge}{\bfseries Q} & \textcolor{insarouge}{\bfseries Seuil proposé} & \textcolor{insarouge}{\bfseries 1 Plein air} & \textcolor{insarouge}{\bfseries 2 Camping, famille} & \textcolor{insarouge}{\bfseries 3 Secours}\\
    Problème~: sans réseau «\,souvent\,» ou «\,très souvent\,» & 7 & ≥ 40\,\% & \qvide & \qvide & \qvide\\ \arrayrulecolor{white}\hline
    Besoin~: «\,important\,» ou «\,très important\,» (4 ou 5 sur 5) & 10 & ≥ 60\,\% & \qvide & \qvide & \qvide\\ \arrayrulecolor{white}\hline
    Problème déjà vécu au moins une fois & 11 & ≥ 30\,\% & \qvide & \qvide & \qvide\\ \arrayrulecolor{white}\hline
    Utilité~: «\,utile\,» ou «\,très utile\,» & 12 & ≥ 50\,\% & \qvide & \qvide & \qvide\\ \arrayrulecolor{white}\hline
    Intention~: «\,probablement\,» ou «\,certainement\,» & 17 & ≥ 30\,\% & \qvide & \qvide & \qvide\\ \arrayrulecolor{white}\hline
    Prix~: «\,oui\,» à 69\,€ & 20 & ≥ 20\,\% & \qvide & \qvide & \qvide\\ \arrayrulecolor{white}\hline
    Prix maximal, valeur médiane & 19 & ≥ 60\,€ & \qvide & \qvide & \qvide\\ \arrayrulecolor{white}\hline
    Appareils à équiper, valeur médiane & 18 & ≥ 2 & \qvide & \qvide & \qvide\\ \arrayrulecolor{white}\hline
  \end{tabular}
  \vspace{0.15cm}
  \begin{pcompact}\centering\footnotesize
    Règle proposée~: segment prioritaire s'il atteint 6 seuils sur 8~; si le «\,oui\,» à 69\,€ reste sous 10\,\% partout, revoir prix ou concept avant le moule.
  \end{pcompact}
  \srcnote{Seuils proposés par l'équipe, à valider \emph{avant} l'enquête (sinon on les ajuste après coup). L'échantillon n'étant pas aléatoire, ils servent à comparer les segments, pas à prévoir les ventes. Cases beiges~: à remplir avec les résultats.}
\end{frame}

%  Fichier de contenu : ni préambule ni \begin{document}, commence par \section.
%
%  Contenu : (1) ce qu'il faut valider, (2) qui interroger et où, (3) plan du questionnaire,
%  (4) questions avant la présentation du concept, (5) questions après, (6) le prix,
%  (7) lecture des résultats et seuils de décision.
%  Cette partie reprend la méthodologie de l'équipe (étude quantitative, 3 segments, échantillon
%  non aléatoire par choix raisonné, 9 sections avant / après un concept neutre, prix maximal
%  puis « À 69 € … ? ») et la complète par : les questions rédigées, les lieux de diffusion,
%  des tailles d'échantillon avec marge d'erreur, et des seuils de décision.
%  Le questionnaire complet (formulations, échelles, logique de saut, variante secours) est
%  dans le fichier autonome questionnaire.tex ; les numéros de question sont les mêmes.
%
%  À DÉCIDER PAR L'ÉQUIPE (propositions de départ, signalées « proposé » dans les diapositives) :
%   - tailles d'échantillon : 150 (plein air), 100 (camping, voyage, famille), 60 (secours) ;
%   - seuils de décision de la dernière diapositive ;
%   - variante à trois prix (59, 69, 79 €) pour la question 20.
%  Marge d'erreur indicative : 1,96 × racine(0,25 / n), soit ± 8 points pour n = 150,
%  ± 9,8 pour n = 100, ± 12,7 pour n = 60 (échantillon aléatoire supposé : valeur indicative seulement).
%%%%%%%%%%%%%%%%%%%%%%%%%%%%%%%%%%%%%%%%%%%%%%%%%%%%%%%%%%%%%%%%%%%%%%%%%%%%%%
\section{Étude quantitative}

% numéro de question en rouge ; \qq{n°}{question}{réponses} : question suivie de ses réponses en gris
\newcommand{\qnum}[1]{{\sffamily\bfseries\color{insarouge}Q#1}}
\newcommand{\qq}[3]{\qnum{#1}\hspace{0.3em}#2\ {\color{insagris}\itshape#3}\par}
% titre de section de questionnaire, dans une carte
\newcommand{\qsec}[2]{\begingroup\footnotesize\sffamily\bfseries{\color{insarouge}#1}\hspace{0.4em}#2\par\endgroup\vspace{1.5pt}}
% cellule à remplir du tableau de résultats (diapositive 7)
\newcommand{\qvide}{\cellcolor{insabeige}}

%---------------------------------------------------------------------------
% Diapositive 1 : ce qu'il faut valider (la chaîne problème -> prix)
%---------------------------------------------------------------------------
\begin{frame}{Étude quantitative~: ce qu'il faut valider}
  \centering
  \begin{tikzpicture}[font=\scriptsize]
    \tikzset{etp/.style={minimum width=1.95cm,minimum height=0.62cm,inner xsep=3pt,font=\scriptsize\bfseries},
             detp/.style={anchor=north,text width=2.05cm,align=center,font=\scriptsize}}
    \node[pastille,etp,draw=insagris,fill=insagris!12,text=black]  (p1) at (0.98,0)  {PROBLÈME};
    \node[pastille,etp,draw=insagris,fill=insagris!12,text=black]  (p2) at (3.26,0)  {BESOIN};
    \node[pastille,etp,draw=insaorange,fill=insapeche,text=black]  (p3) at (5.54,0)  {SOLUTION};
    \node[pastille,etp,draw=insaorange,fill=insapeche,text=black]  (p4) at (7.82,0)  {INTÉRÊT};
    \node[pastille,etp,draw=insarouge,fill=insarose,text=black]    (p5) at (10.1,0) {ADOPTION};
    \node[pastille,etp,draw=insarouge,fill=insarose,text=black]    (p6) at (12.38,0) {PRIX};
    \foreach \a/\b in {p1/p2,p2/p3,p3/p4,p4/p5,p5/p6} {\draw[fluxr] (\a) -- (\b);}
    \node[detp] at ([yshift=-0.12cm]p1.south) {Se retrouve-\mbox{t-on} sans réseau~?\\{\color{insarouge}\qnum{7}}};
    \node[detp] at ([yshift=-0.12cm]p2.south) {Faut-il alors pouvoir communiquer~?\\{\color{insarouge}\qnum{10}, \qnum{11}}};
    \node[detp] at ([yshift=-0.12cm]p3.south) {La solution paraît-elle utile~?\\{\color{insarouge}\qnum{12}, \qnum{13}}};
    \node[detp] at ([yshift=-0.12cm]p4.south) {Dans quelles situations l'emploie-\mbox{rait-on}~?\\{\color{insarouge}\qnum{14}}};
    \node[detp] at ([yshift=-0.12cm]p5.south) {L'achèterait-\mbox{on}~? Qu'est-ce qui freine~?\\{\color{insarouge}\qnum{15} à \qnum{18}}};
    \node[detp] at ([yshift=-0.12cm]p6.south) {À quel prix~? 69\,€ passe-t-il~?\\{\color{insarouge}\qnum{19} à \qnum{22}}};
    % regroupements A, B, C
    \draw[decorate,decoration={brace,mirror,amplitude=4pt},insagris,thick] (-0.05,-1.95) -- (4.15,-1.95)
      node[midway,below=6pt,font=\scriptsize\bfseries,text=insagris] {A · Le problème est-il réel~?};
    \draw[decorate,decoration={brace,mirror,amplitude=4pt},insaorange,thick] (4.6,-1.95) -- (8.8,-1.95)
      node[midway,below=6pt,font=\scriptsize\bfseries,text=insaorange] {B · La solution plaît-elle~?};
    \draw[decorate,decoration={brace,mirror,amplitude=4pt},insarouge,thick] (9.1,-1.95) -- (13.3,-1.95)
      node[midway,below=6pt,font=\scriptsize\bfseries,text=insarouge] {C · Le potentiel commercial};
  \end{tikzpicture}

  \vspace{0.0cm}
  \begin{columns}[T]
    \begin{column}{0.48\textwidth}
      \begin{carte}[height=1.45cm,valign=top]
        \raggedright\gros{«\,Pas de réseau\,» ≠ besoin}\\[1pt]
        {\footnotesize Se retrouver sans réseau ne prouve pas qu'on veuille communiquer~: on mesure les deux.}
      \end{carte}
    \end{column}
    \begin{column}{0.48\textwidth}
      \begin{carte}[height=1.45cm,valign=top]
        \raggedright\gros{«\,Intéressant\,» ≠ achat}\\[1pt]
        {\footnotesize Un concept qui plaît ne se vend pas forcément~: on mesure l'intention, puis le prix.}
      \end{carte}
    \end{column}
  \end{columns}
  \vspace{0.05cm}
  \begin{pcompact}\centering\footnotesize
    Objectif final~: comparer les trois segments, retenir ceux où besoin et potentiel sont les plus forts.
  \end{pcompact}
  \srcnote{Méthodologie de l'équipe. Les entretiens qualitatifs (15 à 20, première vague) disent \emph{pourquoi} et \emph{comment}~; le questionnaire mesure \emph{combien}.}
\end{frame}

%---------------------------------------------------------------------------
% Diapositive 2 : qui interroger, où les trouver, combien
%---------------------------------------------------------------------------
\begin{frame}{Qui interroger, et où les trouver~?}
  \scriptsize\renewcommand{\arraystretch}{1.2}\setlength{\tabcolsep}{4pt}
  \hyphenpenalty=10000 \exhyphenpenalty=10000
  \noindent\begin{tabular}{@{}>{\raggedright\arraybackslash\bfseries}p{2.3cm}>{\raggedright\arraybackslash}p{3.6cm}>{\raggedright\arraybackslash}p{4.8cm}>{\centering\arraybackslash}p{0.8cm}>{\centering\arraybackslash}p{1.0cm}@{}}
    \rowcolor{insarose}\textcolor{insarouge}{Segment} & \textcolor{insarouge}{\bfseries Qui interroger} & \textcolor{insarouge}{\bfseries Où les trouver} & \textcolor{insarouge}{\bfseries Cible} & \textcolor{insarouge}{\bfseries Marge}\\
    1 Plein air, sport outdoor &
      randonneurs, trailers, vététistes, alpinistes, skieurs de randonnée &
      clubs de randonnée et de montagne, groupes de pratiquants en ligne, magasins outdoor, refuges (affiche avec QR code) & 150 & ±\,8\,pts\\
    \arrayrulecolor{insagris!30}\hline
    2 Campeurs, voyageurs, familles &
      campeurs, camping-caristes, voyageurs, familles qui sortent en zone isolée &
      campings (QR code à l'accueil), groupes de camping-caristes et de voyageurs, associations de parents, réseau personnel & 100 & ±\,10\,pts\\
    \hline
    3 Secours et gestion de crise &
      bénévoles et professionnels du secours, radioamateurs, élus et agents communaux, ONG &
      ADRASEC et associations de radioamateurs, Protection civile, Croix-Rouge, mairies, réseaux d'ONG & 60 & ±\,13\,pts\\
  \end{tabular}
  \vspace{0.2cm}
  \begin{columns}[T]
    \begin{column}{0.485\textwidth}
      \begin{carte}[height=1.55cm,valign=top]
        \raggedright\gros{Choix raisonné}\\[1pt]
        {\footnotesize Des personnes concernées, avec de la variété (âge, région, niveau de pratique). Question 1 = filtre~; «\,Autre\,» à part.}
      \end{carte}
    \end{column}
    \begin{column}{0.485\textwidth}
      \begin{carte}[height=1.55cm,valign=top]
        \raggedright\gros{Comparer, pas prévoir}\\[1pt]
        {\footnotesize Les segments entre eux~; pas de ventes ni de part de marché. Marge indicative ≈ ±\,98\,/\,$\sqrt{n}$ points.}
      \end{carte}
    \end{column}
  \end{columns}
  \srcnote{Tailles proposées par l'équipe (310 réponses au total), à ajuster~; marge à 95\,\% pour $p$ = 50\,\%, valable pour un tirage aléatoire~; la population du segment 3 étant petite, sa marge réelle est plus faible. Canaux~: ceux du recrutement des entretiens.}
\end{frame}

%---------------------------------------------------------------------------
% Diapositive 3 : plan du questionnaire (neuf sections, concept neutre au milieu)
%---------------------------------------------------------------------------
\begin{frame}{Plan du questionnaire~: neuf sections}
  \centering
  \begin{tikzpicture}[font=\scriptsize]
    \tikzset{sec/.style={draw=insagris,fill=insagris!10,rounded corners=3pt,minimum width=4.3cm,minimum height=0.55cm,
                         inner sep=3pt,align=left,anchor=west,text width=4.1cm},
             sec2/.style={draw=insaorange,fill=insapeche,rounded corners=3pt,minimum width=4.5cm,minimum height=0.55cm,
                         inner sep=3pt,align=left,anchor=west,text width=4.3cm}}
    % colonne de gauche : avant la présentation
    \node[font=\footnotesize\bfseries\sffamily,text=insagris,anchor=west] at (0,1.45) {AVANT~: problème et besoin};
    \node[sec] at (0,0.85)  {{\bfseries 1 Profil}~: quel segment~?};
    \node[sec] at (0,0.2)   {{\bfseries 2 Habitudes}~: pratiques actuelles};
    \node[sec] at (0,-0.45) {{\bfseries 3 Problème}~: sans réseau~?};
    \node[sec] at (0,-1.1)  {{\bfseries 4 Besoin}~: communiquer, pourquoi~?};
    % centre : concept neutre
    \node[fill=insarose,draw=insarouge,thick,rounded corners=4pt,text width=3.55cm,inner sep=5pt,align=left,anchor=north west]
      at (4.75,1.6)
      {{\sffamily\bfseries\color{insarouge}Concept présenté, sans nom ni prix}\\[2pt]
       \scriptsize Un petit appareil radio de poche, relié au téléphone par Bluetooth. Il envoie des messages texte à des proches
       ayant le même appareil, sans réseau mobile, sans internet, sans abonnement. Portée~: de 0,5 à 15\,km selon le terrain.
       Texte seulement~; ce n'est pas une balise de détresse.};
    % colonne de droite : après
    \node[font=\footnotesize\bfseries\sffamily,text=insaorange,anchor=west] at (8.95,1.45) {APRÈS~: solution et adoption};
    \node[sec2] at (8.95,0.85)  {{\bfseries 5 Intérêt}~: la solution est-elle utile~?};
    \node[sec2] at (8.95,0.2)   {{\bfseries 6 Usages}~: dans quelles situations~?};
    \node[sec2] at (8.95,-0.45) {{\bfseries 7 Freins}~: qu'est-ce qui bloque~?};
    \node[sec2] at (8.95,-1.1)  {{\bfseries 8 Intention}~: acheter ou adopter~?};
    \node[sec2] at (8.95,-1.75) {{\bfseries 9 Prix}~: maximal, puis 69\,€};
    % flèches
    \draw[fluxr] (4.35,-0.2) -- (4.7,-0.2);
    \draw[fluxr] (8.35,-0.2) -- (8.9,-0.2);
  \end{tikzpicture}

  \vspace{0.05cm}
  \begin{columns}[T]
    \begin{column}{0.325\textwidth}
      \begin{carte}[height=1.25cm,valign=top]
        \raggedright\gros{23 questions}\\[1pt]{\footnotesize 8 à 10 minutes, en ligne ou sur papier.}
      \end{carte}
    \end{column}
    \begin{column}{0.325\textwidth}
      \begin{carte}[height=1.25cm,valign=top]
        \raggedright\gros{Concept neutre}\\[1pt]{\footnotesize après la section 4, pour ne pas influencer.}
      \end{carte}
    \end{column}
    \begin{column}{0.325\textwidth}
      \begin{carte}[height=1.25cm,valign=top]
        \raggedright\gros{Pré-test}\\[1pt]{\footnotesize dix personnes~; réponses anonymes (RGPD).}
      \end{carte}
    \end{column}
  \end{columns}
  \srcnote{Structure de la méthodologie de l'équipe~; formulations complètes dans le fichier questionnaire.tex. Portée~: voir la partie Fonctionnement (\cite{techconnect}).}
\end{frame}

%---------------------------------------------------------------------------
% Diapositive 4 : questions 1 à 11 (avant la présentation)
%---------------------------------------------------------------------------
\begin{frame}{Avant la présentation~: questions 1 à 11}
  \begin{columns}[T]
    \begin{column}{0.485\textwidth}
      \begin{carte}[height=1.9cm,valign=top]
        \qsec{1}{Profil}\scriptsize\raggedright
        \qq{1}{À quel profil correspondez-vous le mieux~?}{plein air · camping, voyage, famille · secours et crise · autre}
        \qq{2}{Votre âge}{tranches}
        \qq{3}{Où vivez-vous~?}{ville · périurbain · campagne · montagne}
      \end{carte}
      \vspace{0.06cm}
      \begin{carte}[height=2.2cm,valign=top]
        \qsec{3}{Existence du problème}\scriptsize\raggedright
        \qq{7}{Vous retrouvez-vous dans une zone sans réseau mobile~?}{jamais → très souvent}
        \qq{8}{Où cela vous est-il arrivé~?}{montagne, campagne, étranger, ville, catastrophe…}
        \qq{9}{Comment communiquez-vous alors~?}{rien, talkies, balise, satellite…}
      \end{carte}
    \end{column}
    \begin{column}{0.485\textwidth}
      \begin{carte}[height=1.9cm,valign=top]
        \qsec{2}{Habitudes et pratiques}\scriptsize\raggedright
        \qq{4}{Quelles activités en pleine nature, ces 12 derniers mois~?}{choix multiple}
        \qq{5}{À quelle fréquence~?}{jamais → chaque semaine}
        \qq{6}{Avec qui sortez-vous~?}{seul · à deux · 3 à 5 · 6 ou plus}
      \end{carte}
      \vspace{0.06cm}
      \begin{carte}[height=2.2cm,valign=top]
        \qsec{4}{Importance du besoin}\scriptsize\raggedright
        \qq{10}{Sans réseau, importance de~: rester en contact avec son groupe, prévenir ses proches, partager sa position, coordonner une équipe, alerter les secours}{1 à 5}
        \qq{11}{Cela vous a-t-il déjà posé un vrai problème~?}{jamais · une fois · plusieurs fois}
      \end{carte}
    \end{column}
  \end{columns}
  \vspace{0.05cm}
  \begin{pcompact}\centering\footnotesize
    Secours~: mêmes questions, reformulées. Le concept n'est présenté qu'après la question 11.
  \end{pcompact}
  \srcnote{Sections 1 à 4 de la méthodologie de l'équipe~; réponses complètes et variante secours dans questionnaire.tex.}
\end{frame}

%---------------------------------------------------------------------------
% Diapositive 5 : questions 12 à 18 (après la présentation)
%---------------------------------------------------------------------------
\begin{frame}{Après la présentation~: questions 12 à 18}
  \begin{columns}[T]
    \begin{column}{0.485\textwidth}
      \begin{carte}[height=1.95cm,valign=top]
        \qsec{5}{Intérêt pour la solution}\scriptsize\raggedright
        \qq{12}{Dans quelle mesure cette solution vous paraît-elle utile~?}{pas du tout → très utile}
        \qq{13}{Qu'est-ce qui vous intéresse le plus~? (deux au plus)}{sans réseau · sans abonnement · messages privés · simplicité · prix}
      \end{carte}
      \vspace{0.1cm}
      \begin{carte}[height=1.95cm,valign=top]
        \qsec{7}{Freins à l'adoption}\scriptsize\raggedright
        \qq{15}{Qu'est-ce qui pourrait vous empêcher de l'acheter~?}{prix · équiper les autres · portée · autonomie · complexité · utilité · déjà équipé · marque inconnue}
        \qq{16}{Quel est le frein principal~?}{un seul choix}
      \end{carte}
    \end{column}
    \begin{column}{0.485\textwidth}
      \begin{carte}[height=1.95cm,valign=top]
        \qsec{6}{Usages envisagés}\scriptsize\raggedright
        \qq{14}{Dans quelles situations l'utiliseriez-vous~?}{randonnée · camping · voyage · sport · zones isolées · secours · catastrophe ou panne · événement · jamais}
      \end{carte}
      \vspace{0.1cm}
      \begin{carte}[height=1.95cm,valign=top]
        \qsec{8}{Intention d'achat ou d'adoption}\scriptsize\raggedright
        \qq{17}{Seriez-vous prêt à acheter ou adopter cette solution~?}{certainement pas → certainement}
        \qq{18}{Combien de personnes équiperiez-vous, vous compris~?}{1 · 2 · 3 à 4 · 5 à 9 · 10 ou plus}
      \end{carte}
    \end{column}
  \end{columns}
  \vspace{0.05cm}
  \begin{pcompact}\centering\footnotesize
    La question 18 donne la taille de groupe à équiper~: elle éclaire le choix kit seul, Duo ou Groupe.
  \end{pcompact}
  \srcnote{Sections 5 à 8 de la méthodologie de l'équipe~; échelles à cinq niveaux, ordre des réponses tiré au hasard quand c'est possible.}
\end{frame}

%---------------------------------------------------------------------------
% Diapositive 6 : le prix (questions 19 à 23)
%---------------------------------------------------------------------------
\begin{frame}{Le prix~: de la question ouverte à la question précise}
  \begin{columns}[T]
    \begin{column}{0.57\textwidth}
      \begin{tikzpicture}[font=\footnotesize]
        \draw[insarouge!40,thick] (0.2,0.2) -- (0.2,-3.35);
        \node[num] at (0.2,0)     {1};
        \node[num] at (0.2,-1.2)  {2};
        \node[num] at (0.2,-2.4)  {3};
        \node[anchor=west,text width=6.5cm,align=left] at (0.5,0.05)
          {\qnum{19}~: {\bfseries Quel prix maximal paieriez-vous, sans abonnement~?}\\[-1pt]
           {\scriptsize réponse libre, en euros, \emph{avant} de montrer 69\,€~: pas d'ancrage.}};
        \node[anchor=west,text width=6.5cm,align=left] at (0.5,-1.2)
          {\qnum{20}~: {\bfseries À 69\,€, sans abonnement, l'achèteriez-vous~?}\\[-1pt]
           {\scriptsize oui, peut-être, non~: teste l'hypothèse de prix.}};
        \node[anchor=west,text width=6.5cm,align=left] at (0.5,-2.4)
          {\qnum{21}~: {\bfseries Pour plusieurs personnes~: 2 kits à 139\,€, 4 à 249\,€, un seul à 69\,€, aucun~?}\\[-1pt]
           {\scriptsize teste les lots Duo et Groupe de la partie SWOT.}};
        \node[anchor=west,text width=6.5cm,align=left,font=\scriptsize,text=insagris] at (0.5,-3.3)
          {\qnum{22}~: où l'acheter (boutique, site, club, commune)~; \qnum{23}~: contact facultatif pour un test ou une précommande.};
      \end{tikzpicture}
    \end{column}
    \begin{column}{0.4\textwidth}
      \begin{encadre}{Comment lire}
        \begin{puces}
          \item \textbf{Part solvable à 69\,€}~: part des prix maximaux (Q19) au moins égaux à 69\,€.
          \item \textbf{Recette à un prix $p$}~: $p$ × part des prix maximaux ≥ $p$~; comparer 59, 69 et 79\,€.
          \item \textbf{Effet du prix}~: écart entre l'intention (Q17) et le «\,oui\,» à 69\,€ (Q20).
        \end{puces}
      \end{encadre}
    \end{column}
  \end{columns}
  \vspace{0.05cm}
  \begin{pcompact}\centering\footnotesize
    Variante proposée~: tirer au hasard 59, 69 ou 79\,€ à la question 20, pour lire trois points de la courbe de demande.
  \end{pcompact}
  \srcnote{Ordre « prix maximal puis 69\,€ » de la méthodologie de l'équipe. Sensibilité du modèle financier~: ±\,5\,€ de prix = ±\,92\,k€ de résultat en année 3 (partie SWOT).}
\end{frame}

%---------------------------------------------------------------------------
% Diapositive 7 : lire les résultats, seuils de décision
%---------------------------------------------------------------------------
\begin{frame}{Décider~: indicateurs et seuils, segment par segment}
  \scriptsize\renewcommand{\arraystretch}{1.25}\setlength{\tabcolsep}{4pt}
  \hyphenpenalty=10000 \exhyphenpenalty=10000
  \noindent\begin{tabular}{@{}>{\raggedright\arraybackslash}p{4.65cm}>{\centering\arraybackslash}p{0.7cm}>{\centering\arraybackslash}p{1.5cm}>{\centering\arraybackslash}p{1.75cm}!{\color{white}\vrule width 1.5pt}>{\centering\arraybackslash}p{1.75cm}!{\color{white}\vrule width 1.5pt}>{\centering\arraybackslash}p{1.75cm}@{}}
    \rowcolor{insarose}\textcolor{insarouge}{\bfseries Indicateur} & \textcolor{insarouge}{\bfseries Q} & \textcolor{insarouge}{\bfseries Seuil proposé} & \textcolor{insarouge}{\bfseries 1 Plein air} & \textcolor{insarouge}{\bfseries 2 Camping, famille} & \textcolor{insarouge}{\bfseries 3 Secours}\\
    Problème~: sans réseau «\,souvent\,» ou «\,très souvent\,» & 7 & ≥ 40\,\% & \qvide & \qvide & \qvide\\ \arrayrulecolor{white}\hline
    Besoin~: «\,important\,» ou «\,très important\,» (4 ou 5 sur 5) & 10 & ≥ 60\,\% & \qvide & \qvide & \qvide\\ \arrayrulecolor{white}\hline
    Problème déjà vécu au moins une fois & 11 & ≥ 30\,\% & \qvide & \qvide & \qvide\\ \arrayrulecolor{white}\hline
    Utilité~: «\,utile\,» ou «\,très utile\,» & 12 & ≥ 50\,\% & \qvide & \qvide & \qvide\\ \arrayrulecolor{white}\hline
    Intention~: «\,probablement\,» ou «\,certainement\,» & 17 & ≥ 30\,\% & \qvide & \qvide & \qvide\\ \arrayrulecolor{white}\hline
    Prix~: «\,oui\,» à 69\,€ & 20 & ≥ 20\,\% & \qvide & \qvide & \qvide\\ \arrayrulecolor{white}\hline
    Prix maximal, valeur médiane & 19 & ≥ 60\,€ & \qvide & \qvide & \qvide\\ \arrayrulecolor{white}\hline
    Appareils à équiper, valeur médiane & 18 & ≥ 2 & \qvide & \qvide & \qvide\\ \arrayrulecolor{white}\hline
  \end{tabular}
  \vspace{0.15cm}
  \begin{pcompact}\centering\footnotesize
    Règle proposée~: segment prioritaire s'il atteint 6 seuils sur 8~; si le «\,oui\,» à 69\,€ reste sous 10\,\% partout, revoir prix ou concept avant le moule.
  \end{pcompact}
  \srcnote{Seuils proposés par l'équipe, à valider \emph{avant} l'enquête (sinon on les ajuste après coup). L'échantillon n'étant pas aléatoire, ils servent à comparer les segments, pas à prévoir les ventes. Cases beiges~: à remplir avec les résultats.}
\end{frame}
            % 10. Étude quantitative : cibles, questionnaire, seuils de décision
%%%%%%%%%%%%%%%%%%%%%%%%%%%%%%%%%%%%%%%%%%%%%%%%%%%%%%%%%%%%%%%%%%%%%%%%%%%%%%
%  Partie 7 -- Risques et conformité, conclusion et prochaines étapes (2 diapositives)
%%%%%%%%%%%%%%%%%%%%%%%%%%%%%%%%%%%%%%%%%%%%%%%%%%%%%%%%%%%%%%%%%%%%%%%%%%%%%%
\section{Risques et suite}

%---------------------------------------------------------------------------
% Diapositive 19 : risques et conformité
%---------------------------------------------------------------------------
\begin{frame}{Risques et conformité~: ce qu'il faut maîtriser}
  \renewcommand{\arraystretch}{1.22}
  \footnotesize\hyphenpenalty=10000 \exhyphenpenalty=10000
  \noindent\begin{tabularx}{\textwidth}{@{}>{\bfseries\raggedright\arraybackslash}p{3.0cm}>{\raggedright\arraybackslash}X>{\raggedright\arraybackslash}X@{}}
    \rowcolor{insarose}
    \textcolor{insarouge}{Risque} & \textcolor{insarouge}{Pourquoi} & \textcolor{insarouge}{Notre réponse} \\
    Conformité et droit & Marquage CE (directive RED)~; chiffrement à vérifier hors UE. &
      Tests en laboratoire dès l'an 1 (25\,k€)~; avis juridique avant export. \\
    \arrayrulecolor{insagris!30}\hline
    Règles radio & 10\,\% de temps d'émission, 500\,mW ERP au plus (869,4 à 869,65\,MHz). &
      Région Europe préréglée, antennes d'origine, guide d'usage. \\
    \hline
    Densité du maillage & Sans voisins équipés, le réseau ne sert à rien. &
      Packs de 2 à 3 kits, démarrage par zones, cartes de couverture. \\
    \hline
    Projet libre amont & Meshtastic peut changer de cap~; des alternatives existent (MeshCore). &
      Code GPL~: fork possible, développeurs contributeurs, veille. \\
    \hline
    Concurrence du fait-maison & La carte seule coûte ≈ 17\,€~; les passionnés assemblent eux-mêmes. &
      Valeur = prêt à l'emploi, support, garantie~; cible non technique. \\
    \hline
    Promesses et limites & Débit faible, aucun secours relié~: risque d'attentes excessives. &
      Messagerie texte, jamais présentée comme une balise de détresse. \\
    \hline
    Trésorerie & Le stock importé se paie avant de vendre. &
      Précommandes, achats par lots, financement participatif. \\
  \end{tabularx}
  \srcnote{Sources~: \cite{red_2014_53,ttn_eu868,meshtastic_crypto,techconnect}}
\end{frame}

%---------------------------------------------------------------------------
% Diapositive 20 : conclusion et prochaines étapes
%---------------------------------------------------------------------------
\begin{frame}{Conclusion~: l'essentiel et la suite}
  \begin{columns}[T]
    \begin{column}{0.47\textwidth}
      \begin{carte}[height=1.25cm,valign=top]
        \raggedright\gros{Un besoin réel}\\[1pt]
        {\footnotesize Réseau coupé, cher, dépendant d'opérateurs.}
      \end{carte}
      \vspace{0.1cm}
      \begin{carte}[height=1.25cm,valign=top]
        \raggedright\gros{Une réponse simple}\\[1pt]
        {\footnotesize Kit radio + logiciel libre, 69\,€ sans abonnement.}
      \end{carte}
      \vspace{0.1cm}
      \begin{carte}[height=1.25cm,valign=top]
        \raggedright\gros{Un modèle viable}\\[1pt]
        {\footnotesize Point mort à 21\,300 kits/an, 450\,k€ à financer.}
      \end{carte}
    \end{column}
    \begin{column}{0.5\textwidth}
      \vspace{0.2cm}
      \begin{tikzpicture}[font=\footnotesize]
        \draw[insarouge!40,thick] (0.2,0.25) -- (0.2,-3.65);
        \node[num] at (0.2,0)     {1};
        \node[num] at (0.2,-1.2)  {2};
        \node[num] at (0.2,-2.4)  {3};
        \node[num] at (0.2,-3.6)  {4};
        \node[anchor=west,text width=5.6cm,align=left] at (0.5,0)
          {{\bfseries 0 à 3 mois}\\Prototype et dix kits testés sur le terrain (montagne, ville, ONG).};
        \node[anchor=west,text width=5.6cm,align=left] at (0.5,-1.2)
          {{\bfseries 3 à 6 mois}\\Nom et marque à vérifier, précommandes, dossier de financement.};
        \node[anchor=west,text width=5.6cm,align=left] at (0.5,-2.4)
          {{\bfseries 6 à 12 mois}\\Tests de conformité CE, première série de 2\,000 kits.};
        \node[anchor=west,text width=5.6cm,align=left] at (0.5,-3.6)
          {{\bfseries Année 2}\\ONG pilotes, relocalisation de la carte en France.};
      \end{tikzpicture}
    \end{column}
  \end{columns}
  \vspace{0cm}
  \begin{idee}\centering
    {\sffamily\bfseries\color{insarouge}Là où le réseau s'arrête, le tien commence.}
  \end{idee}
\end{frame}
             % 11. Risques, prochaines étapes, conclusion


%%%%% Sources (sources.bib : appel \cite{cle} dans les diapositives) %%%%%
% Liste sur 3 colonnes, un peu plus large que le texte courant, coupée en trois diapositives (filtres srcA, srcB, srcC).
\newcommand{\diaposources}[2]{%
  \begin{frame}{#2}
    \etiquette{Sources}
    \makebox[\textwidth][c]{\begin{minipage}{13.6cm}
      \setlength{\columnsep}{0.3cm}
      \begin{multicols}{3}
        \printbibliography[heading=none,check=#1]
      \end{multicols}
    \end{minipage}}
  \end{frame}}
\diaposources{srcA}{Sources 1/3~: pages consultées du 5 au 8 octobre 2026}
\diaposources{srcB}{Sources 2/3}
\diaposources{srcC}{Sources 3/3}

\end{document}
